\subsection{Protected structural calculations}
\begin{lemma}[Lower structural transitions and the size of the protected test]\label{lem:transition}
On a protected unexceptional structural step a component is measured inside one exact extended source during Grow, or one ordinary home afterwards. On a step with no birth or ballot and no address normalization, a strict interior of such a component survives. On a recruitment step an apparent vacancy can fail to take the proposed birth for one step if its old nominal state is live but unserved.
\end{lemma}
\begin{proof}


A strict interior is between two positive matching served neighbors with consecutive exact integer addresses. Their three nominal states are live and matching since any served live state equals its nominal state. The lower triple forces the middle proper. It is therefore computing, and its next nominal state is its native transition. Its own left and right bonds have grade two. Purge has no weak right bond for a kill, Shrink sees no zero face, and temporary and true withdrawals likewise have no zero face. Other nonballot maintenance operations hold a live incumbent. Proper successor regularity and proper-live service at the lower rank turn this positive nominal result into an actual served result. At grade one, even an occupied adjacent true home cannot provide strong Purge backing. For example if address \(P- 1- k\) is vacant for \(2\le k<q\), the k remaining sites immediately to its right receive zero in the q propagations and the rightmost has a weak right bond at the wrap. They can all die in the killing pass.

If x is live but y is D, the definition gives \(g=0\). On this step x' is D by the noncomputing clause even if the native map on y requests birth. The actual-support clause at the new cut makes this tile proper dead. Its served value is also D. This accounts for one failed transition at that slot. A subsequent recruitment still needs its full lower advance hypotheses. If x' is instead live from \(y=D\), then \(g=1\), so \(x=D\). Classification makes this old slot proper dead. It has a proper successor and the live output is served. An undisturbed clear preparation row followed by a complete Purge before the attempted birth is sufficient to remove the first obstruction as shown in the buffered clearance proof.

Every nondead nominal output has a genuine previous served incumbent or donor within radius one, by the noncomputing clause, exact computation, and the three-dead quiescent rule. Its spatial residue and clock identify the same parent phase on tracing this edge backward. A proper next state is served when live. A matching lower member uses its extended address modulo P, including in deciding the occupancy of a displaced geographical home. Undisturbed means incoming-clear in the \(64A\) pad.

When invoking classification on a shifted home, take its own preceding work, its three suppliers' preceding transactions and their one-neighbor input lanes. Adjoining inputs or gates are queried inside the outer protected box. Each work interval includes its own service lookahead. The transfer of disturbance requires one preceding predecessor comparison and the two issuing epochs for the overlap argument. In a single terminal comparison only the two adjacent authentic root cores are relevant to a remote marker: a farther source cannot reach it across an intervening core and never has an overhang exceeding \(1.1P\). The marker, two source origins, and \(300B\) and \(0.2v\) margins occupy fewer than eight parent home widths since \((300B+v)/P\) tends to zero. Paired seed and two-period tests cover six input homes and two further periods. The directed recipient checks nine home widths and at most five source homes for each queried displaced home. These are finite shifts in units of the current parent tile; \(608<m=\lfloor w/5\rfloor\) at the global choice \(w\ge10^{10}\). Their lower-rank calls take their own m lower tiles of margin, of width \(m/P\) and duration \(m/T\) in current units, and are not enlarged by m \emph{parent} tiles recursively. Large \(T/P\) in the unrestricted cap ancestry cone is accounted for separately in the centered tests of Lemma~\ref{lem:ascent}.

\textbf{Use.} The no-split argument, all stencil birth calls, the actual cap ascent, and the initialized two-period service assertion use these conventions. The exceptional numerical allowances are specified where the corresponding geometry is proved.

For the localization count, a direct same-rank query spans at most 16 additional owner widths: the largest directed stretch is \([- 2P- 8B,3P+8B]\) plus two widths on each side for the potential prior source in a normalized home, four lower guard positions and the neighboring home clipped by rounding. Six seed homes, two adjacent message routes, lookup at the three central positions, and all terminal, \(300B\) and \(0.2v\) margins fit this same bound. The rear-wall radius \(3v\) \(d_0+2w^6+12d_0+20\) and the polynomial early-seed cone are \(<P/10\) for sufficiently large w; these are the unrestricted front cones. A direct temporal lookup takes at most six periods: two past weak levels with the neighbors' computations, the prior work for disturbance, and two periods for the initialized descendants, including a rounding period where necessary. Count the following ordered within-rank calculations separately: observed rows, ordinary relaxation and elementary wound response, common full seed and transport bridge, clean weak issue, general weak issue, strong issue, classification, authentic source, thin trace, wide trace, density, disturbance inheritance, overlap/triple, prospective contender exclusions, common withdrawal calendar, current service/proper successor, advance, two-period descent, and the protected query. The elementary response is applied first with a supplied long interval and subsequently with the actual issued source. Allow two extensions per entry for the pairwise comparisons or the two charged searches. Direct displacements therefore fit within 19\(\cdot\)2\(\cdot\)\(16=608\) owner widths and 19\(\cdot\)2\(\cdot\)\(6=228\) owner periods; no within-rank request starts the same induction all over again. A complete actual lower test needs m lower units, and the finer residual margins sum to at most \(2m/P\) and \(2m/T\) current units. These are less than one and the entire region fits in the own collar. In child units a common enlarged exception has both spans at most \(3w+2m+4+2\)\(\cdot\)\(608<4w+8\) for the chosen w; this even pays extra fixed stencil offsets if grouped in advance. The cap uses separately centered tests with the larger \(2T\)-child ancestry zone, which is not charged to one ordinary owner's collar.
\end{proof}

\subsection{Ordinary precommit}
The ordinary precommit adapts G\'acs \cite[Sections~12--13]{gacs86}.
\begin{lemma}[Ordinary relaxation from arbitrary rows]\label{lem:ordinary-start}
Work from an arbitrary normalized input in one home, with the lower transition convention and the common exception. An interior vacant band of width \(4B+4\), disjoint from K and the two true endpoint sites, at a retired strobe persists at subsequent retired strobes during withdrawal and ordinary work until the ballot, with a vacant center throughout. A home capable of genuine computation has a single genuine interval with both end gaps smaller than \(16B+2d+24\) at a clean sample at most \(50A d_1\) after normalization. At that sample the \(3A\) and \(6A\) threshold positions and their neighbors at both ends are proper strict interior sites. The other branch is erased except for immediate raw-only material.
\end{lemma}
\begin{proof}
Let \(H=[0,P)\) in child coordinates. Sample after a complete provisional retirement, immediately before each actual \(\operatorname{Shrink}(0)\) instruction. Thus the step after the sample is a true Shrink. All following samples include the q intervening Main cycles with full \(\mathsf{Temp}\) resets, B attempts and B removals. At a Grow-to-ordinary transition normalize first. Allow one whole initial cycle to clear marks and start fresh \(\mathsf{Cn}\). All statements concern one clock and address residue group actually supplied by the lower level. Lower optional omissions and failed births are allowed. No second source can put a positively consistent ordinary member in H across a true wrap: the proposed address must be in the donor's own ordinary home.

\textbf{Empty band.} A strip of \(4B+4\) consecutive vacant positions at a retired strobe, disjoint from K and ordinary endpoint positions, retains an empty center on every subsequent date of the final Grow withdrawal and ordinary work until the ballot. It returns to vacancy at each complete retirement. No true Grow births occur on this interval. In each probe all entrants into a currently empty strip are newly marked, except raw assignments which were excluded. At most B sites can be entered from each edge because recruitment moves by one per attempt. The two entering portions therefore never meet. At each removal the furthest surviving entrant from each side, if any, has no positively graded neighbor towards the center. A positively graded neighbor would itself belong to this same clock and residue group. This entrant is marked and its facing site is not its true endpoint (in Grow all potential true endpoints have the same pair of residues modulo P). It dies. Inducting through B removal rounds eliminates the entrants, even if an outside original donor has died. Holes in an entrance only accelerate this bound: the furthest coordinate retreats at least one until its original strip edge. Each true withdrawal, normalization, Shrink and ordinary computation between retired probes can only erase members of the group. A ballot is excluded. From an arbitrary partial phase, an initially empty \(8B+8\) band contains an empty central \(4B+4\) band at the next strobe: from either edge there are at most B remaining attempts before a complete retirement. The central protected site remains empty up to this next strobe, even if old flags were arbitrary outside.

\textbf{Clean gaps.} In an ordinary home a clean pre-ballot step cannot split the descendants of a single component (the no-split convention of Lemma~\ref{lem:transition}). In an incumbent's strict interior the two exact consecutive bonds and properness force a hold on the structural instructions. On a birth, an old end with an inner neighbor and an actual outer child stays between matching nominal successors. The nominal triple forces the old end served on the new row. A singleton with two children likewise stays served, and with one child its set of descendants remains connected. On a clean nonbirth step only ends of the same ordinary home/component can disappear. Components can merge but cannot pass one another: a holder cannot accept a proposal at an occupied position, and there is no swap through a nearest-neighbor pair. These arguments concern actual descendants up to the next step and allow an entire component to vanish.

Take two adjacent components still separate at a later inspection. Propagate their component classes back one sampled interval; if either descends from several components, keep the rightmost contributing old component on the left and leftmost contributing one on the right. They have ordered, disjoint ancestors. If there are intermediate old components, any that contribute to neither final class have vanished or have merged into another class. In the latter case a surviving intermediate class would spatially separate the supposed adjacent two, which is impossible. Fix any vacancy between the two ancestral groups at the previous strobe, for instance the gap immediately after the last contributing left component. Both extreme endpoints surrounding an old internal vacancy see a grade-zero face and fail the outward-exemption test on that face. \(\operatorname{Shrink}(0)\) removes each, enlarging the empty interval. If an entire edge bank dies, the enlarged vacant positions still exist. Before the next sample no old material moves into an empty coordinate; only marked provisional attempts can enter. To persist at a retired strobe, such a frontier directed into the vacancy must either reach a true end, or meet a matching incumbent or a facing probe on the far side. A true end cannot be reached while the other surviving ancestral class still separates it from the end of the home. A meeting of the two final ancestral classes is impossible because they could never separate again. A meeting with a formerly intervening third component makes that component an ancestor of the joining class if the joined interval has any surviving descendants: once joined, label the entire interval by the union of its component labels. Consequently choosing the gap immediately after the last contributing left component excludes this possibility for persistent advance from the left. On the right, choose the gap immediately before the first contributing right component if there are intervening disappearing components. A discarded middle component may be crossed only after it has disappeared; it cannot leave a lasting supporting front for either side without becoming a contributor. It follows that the final gap contains the chosen old gap enlarged by at least one cell of the bordering contributing component on its populated side. Equivalently its left edge is no farther right than the last old contributing left edge minus one, and its right edge is no farther left than the first old contributing right edge plus one. This latter statement follows separately for each class by looking at its outermost probe at each of the q complete cycles. Without a join to a new contributing class its marked outside extension has a zero unexempt face and the B removals peel it completely. Thus an internal gap at a later sample has an internal ancestor gap at every preceding clean sample; tracing repeatedly gives width at least 1 plus the number of clean \(\operatorname{Shrink}(0)\) calls. A single genuine interval, once reached, is preserved in clean ordinary work up to extinction or a cut.

\textbf{Early certificate.} After the initial cycle take two consecutive trial stretches of \(20A+1\) Shrink spacings, leaving two complete spacings between them. S has length at most \(d<d_1\), so at least one stretch is wholly clean including boundary samples. Its terminal sample is within \(50A\) \(d_1\) of normalization since \(2(20A+1)+2+2<50A\). If it has more than one component it has an internal gap larger than \(20A\). Such a gap contains a sterile band of width \(4B+4\) disjoint from K and from both true ends: remove K and the two true endpoints, splitting the interval into at most four pieces; \((20A- d- 2)/4>4B+4\). The same holds for a true-end vacant interval of width \(16B+2d+24\): after removing K and the one true-edge pair at the boundary there remain at most two nonempty pieces with total length at least \(16B+d+22\). One contains a sterile band. On every clean stage-zero inspection the strobe point nearest to that band on either side, if any, disappears. Its face toward the spine is broken at a strictly internal depth; the exposed site is eligible for Shrink. No forward provisional gain persists beyond the newly recessed strobe frontier: there is no other same-tribe material between it and the band, the band's center never becomes occupied, and no true endpoint lies between the frontier and the spine on its home side. Each clean inspection reduces the available distance to the far end of H by at least one per such side. Hence \(P+1\) consecutive clean inspections suffice to erase both sides, including a last singleton abutting the opposite true end (its spine-facing face is internal). More than \(3P+10\) calls remain because \(\lfloor 3^w/q \rfloor - 50A>3P+10\). If the pulse is still ahead, removing its touched spacing and the two adjacent partial spacings leaves more than \(3P+6\) calls in at most two stretches. At least one stretch contains \(P+1\) consecutive clean calls. If this stretch follows S every remaining occupant dies. If it precedes S the home is empty at S, and every later occupant traces to K. Before the next complete initialized clean Purge there are at most four recruitment instructions in this phase, so its whole span is at most \(d+8<q\). With fresh \(\mathsf{Cn}\), after q propagations it has no positive right tip, and all its components disappear from the right within q killing ticks. Ordinary recruits from outside H are forbidden. Thus the only late alternative on this erased branch is immediate raw-only material.

A home capable of a genuine productive computation must therefore have had a single genuine interval at the above sample, with \emph{both} true-end gaps smaller than \(16B+2d+24\). Each \(3A\) and \(6A\) threshold and its two neighboring positions on both sides of H are strictly interior, since \(16B+2d+25<3A\); at this clean sample each is the middle of a matching nominal triple, hence proper by the rank-\((j-1)\) clause of Proposition~\ref{prop:organization}\ref{org:classification}. Their proper incumbents then follow every native live hold by lower proper-successor and service. Any later genuine reoccupation starts with proper birth. This is the input to the one-event ledger, independent of nominal separation at the new rank.
\end{proof}

\begin{lemma}[Clearance between held flanks]\label{lem:flanks}
Suppose at the row just before S there is an unbroken canonical stretch between two intervals of 24 old positions held continuously from that row to \(t_E\). From \(t_E\) to the contemplated fresh Heal attempts suppose that each flank remains one canonical interval, has at least 24 positions, and in one tick its left edge can move right by at most one and its right edge left by at most one. Enlargement when it absorbs another compatible interval is allowed. Stop when flanks meet or either disappears. For the questioned nine-slot cross-section assume that its furthest spatial coordinates extended by three on either side, across the current and four preceding dates, are within or between the flanks, allowing use of their innermost 16 positions as a buffer. Then every foreign-or-wrong-phase served obstruction in that middle region descends from K\(\times\)S; on and after the last affected date plus \(4q+10\) there is no such obstruction. The same conclusion holds simultaneously for all the foreign nominal incoming stencils at dates from four later, until the stopped comparison ends. For lower native clean births a previous clean row gives properness of the input slots; the recipient's possibly improper old nominal incumbent may block one attempt. If the walls come from an earlier common full seed of a pair of homes, with an unmarked strip of \(q+9\) proper old sites disjoint from K on each side of an interior pulse (in the same home or, when K is within \(q+10\) of the true wrap, flush with the seam in the neighboring full home) and the survivor cores held, the paired corridor between those cores is repaired in one pass within \(d+4q+14+2d_0<w^6\) ticks after the start of S, with proper interior triples throughout the restored corridor.
\end{lemma}
\begin{proof}
Use Cartesian child rectangles of unit width and time. A shifted child grid has offsets a,b in \([0,1)\); its tile with indices i,t has rectangle \([i+a,i+a+1)\)\(\times\)\([t+b,t+b+1)\). On one fixed shifted grid, any nominal live output has a served live parent in the preceding row at spatial index distance at most one, and a served live state is nominal live. On the same geographical grid a different native phase cannot occupy a slot already occupied by a canonical state. On distinct geographical child grids the protected lower overlap exclusion forbids intersecting tiles both nominal live. All coordinates below are at that child scale. An incoming pad for a queried lower child is \(64A/P<1\) unit. At rank one there is only the one physical grid and native mandatory births can be tested directly, without a lower properness call.



\textbf{Wall test.} For times up to \(t_E\) choose a fixed interior site in each 24-site old run. Afterwards track an interior site in each flank by clamping the previous site into \([\mathrm{left}+8,\mathrm{right}- 8]\), inclusive. This interval has length at least 8. Its left boundary can move right by \(\le\)1, its right boundary can move left by \(\le\)1. The selected site therefore moves by at most one per canonical tick, including when a flank enlarges abruptly. It has at least eight uninterrupted occupied positions on both sides. For the region outward of the selected site but inside its own held flank, every possible competing served rectangle meets a canonical occupied rectangle and is excluded directly. For the region between the two selected tracks, consider any traced foreign tile path through successive rows, following a served parent each time. At the observed tile choose a time inside its intersection with the canonical queried time row. For the served parent of a nominal tile starting at \(\sigma\) with \(t- 1<\sigma <t+1\), its own interval \([\sigma - 1,\sigma )\) meets \([t- 1,t+1)\), so choose a comparison time in this intersection (for \(\sigma =t\) use the preceding row). Follow this same fractional offset inside every earlier tile of its lineage, subtracting one per backward step. Compare its spatial midpoint to the selected index on the canonical row containing that time. Successive comparison times differ by one and the index of the canonical row increases by one. Only the current and previous canonical rows are queried. The spatial midpoint moves by \(\le\)1 and the selected index by \(\le\)1; their signed difference changes by \(\le\)2. A path switching sides therefore includes a midpoint within 2 of the selected position. Every such foreign tile intersects an occupied canonical position in that row (half-site boundary rounding costs at most one). It is forbidden by lower exclusion, or by unique occupancy on the same grid. The same calculation handles a tile crossing the old bottom row just before S: the whole segment between the old selected tracks was occupied. Thus each permissible traced foreign lineage must encounter the exception. This also proves that none can be supplied from beyond a shortened but still enduring flank. After a join the comparison ends.

\textbf{Uniform killing time.} Follow a trapped age/residue class backwards beyond its last raw ancestor. On uncompromised steps a changed spatial position entails a successful birth; lower omissions alone cannot produce a new position. A fixed native clock can birth only on a handle instruction, once per \(H_0=2q+2\) ticks or less often. From first S through last \(S+4q+10\) there are \(\le\)\(d+4q\)+11\textless{}\(3H_0+1\) ticks for \(w\ge 100000\); even upon rounding both shifted temporal boundaries, allow four opportunities. Each descendant is therefore within K expanded by six shifted positions on either side, counting two more for the spatial rounding and the old site versus updated site. Its total width is \(\le\)\(d+12\)\textless{}q. For each alien clock after the last affected update, at most \(H_0- 1\) further ticks precede the next clean Purge initialization; the entire initialized Purge uses \(2q+1\) ticks. Two rounding rows suffice to translate these dates back to the canonical time grid, giving a common upper bound last \(S+4q+6\). No outside immigrant can supply strong backing during this operation. Wrong phases, residues or geographical grids at the two canonical walls cannot supply a strong link to the trapped alien. In a fixed alien age class with distinct ordinary sources, only exact consecutive addresses within one root can give grade two; the whole trapped age class in any case fits \(d+12\). If a normalization occurred earlier, bound the union of the original footprints, all in this same column. Each surviving \(\mathsf{Cn}\)=1 after propagation would require q strong left predecessors, which do not exist in this short span. The right break eats every remaining zero component in at most q killing ticks. No birth occurs during the Purge. The common date bound holds separately for every shift, with no sum over the number of origins. After this time a new alien interior site would need another raw pulse or a path across a held wall, neither of which is available.

\textbf{From served to nominal.} A foreign nominal non-proper-dead blocker is nominal live by lower classification. Its preceding served parent is at most one spatial shift index away and one time index earlier. A lower stencil querying four neighboring slots and a \(<1\) pad can intersect a shifted tile that extends by one further unit; use its parent, then one additional prior row to certify lower proper successor. The total extra spatial reach from the outermost stencil site is \(<1\) pad +1 tile width +1 parent step, hence at most three slots; the corresponding temporal reach is at most four rows, including the previous-row proper-successor test and both rounding endpoints. All these areas are within the buffered wall hypothesis on rows up to and including the queried row; no prospective future wall beyond a first actual join or disappearance is required. Canonical undisturbed live inputs at that previous date acquire proper successors. Canonical nominal D inputs are proper dead by lower classification, even if a foreign native residue lives in a different slot. A geometric competing origin is gone by the previous calculation. Thus the required nine inputs are simultaneously proper and incoming-clear at the actual recruitment step. For an old vacancy that could conceal a canonical unserved nominal-live state, any preceding nonbirth step with no canonical incumbent there has nominal output D regardless of its service indicator. The initialized Purge just before every attempt provides such a preceding date or a true prior matching incumbent. If the latter is a proper incumbent and outputs nominal live, its positive successor remains proper and served. An initially improper incumbent can be omitted for one attempt. Consequently the exact native birth into a qualified clean vacancy is served.

\textbf{Previously full paired seam.} When separated walls originally form a full seeded pair, static inner backings suffice. On each side of an interior pulse choose an unmarked contiguous strip of \(q+9\) proper old sites outside K in the same home. Its \(q+1\) interior survives the at most two Shrink calls before the next requested complete pass, by the fixed-backing proof in Lemma~\ref{lem:ordinary-response}. Take a one-tick properness and edge reserve. If K is within \(q+10\) of the true wrap, an inadequate strip on the near side of the wrap is replaced by one in the neighboring full home, flush with the seam and disjoint from K. If K straddles the seam take one strip in each home outside it. If a fully backed hole lies deep inside one home take both walls there. At the resulting in-home faces one reaches each damaged position or its legitimate true end from a \(q+1\)-backed strip in at most \(q+d+32\) attempts: the short old near-wrap portion is \(<q+10\), K has \(\le\)d positions, the two-possible-Shrink and stencil reserves use \(\le\)20 more, and at most two attempts are blocked by initially improper nominal incumbents. For a deep cut its width including these reserves is \(\le\)\(d+20\). Outside the two stationary \(q+1\) backings the source core used for transport lies in the already held part of each computing home. This is applied only to seams in the corridor between the two inner cores of the full pair; its two exterior true ends are irrelevant.

Wait for the common clearance above, and use the first subsequent reset / growth pass (at most another \(d_0\) until its reset and \(d_0\) until completion). On the B growth attempts all exposed in-home front sites adjacent to the hole have legal proposals, unoccupied recipients see either matching age or D one step farther, and all nine inputs are clear as just proved. A legal competing bid from the opposite bank in the same home has exactly the same canonical address at the recipient; a bid across the true seam is illegal for that bidder's ordinary home. Existing correctly occupied positions along an advancing front need no replacement. During this pass no real Shrink occurs, every new birth is proper, and each connected extension containing the old backing survives each intervening Purge: its rightmost site has q strong predecessors, and all earlier sites have a strong right neighbor. Allow at most two attempts blocked by initially improper nominal incumbents. Every other attempt fills the next vacant site of each backed front or finds that its target has already been reached. Since \(q+d+32<B\), the two legitimate ends or the joined same-home fronts are reached in one pass. Between the restored full seam and the two cores all interior lower triples are proper; a wrap is positively bonded but each home supplies its own strong Purge backing. Subsequent clean operations preserve that corridor. All response dates fit \(d+4q+14+2d_0<w^6\). A later good input transfer and all its shifts require that their whole transaction interval is disjoint from this charged slab, not merely that its first read date is disjoint.

\textbf{Dependencies.} The unconditional application to a general old single interval requires the two enduring, sufficiently buffered flanks and a bound on their cavity width until the chosen pass. Lemma~\ref{lem:ordinary-response} and Lemma~\ref{lem:attachment} discuss the marginal cases. The paired-seam result additionally needs an earlier actual common full seed as supplied by Lemma~\ref{lem:seed}, and the survivor cores of the computing homes. Thus this is a proven barrier/clearance mechanism conditional on the explicit spatial and temporal wall premises.
\end{proof}

\begin{lemma}[Local response from a supplied macroscopic interval]\label{lem:local-wound}
Assume the lower-rank hypotheses, the structural transitions of Lemma~\ref{lem:transition} and the held-wall conclusions of Lemma~\ref{lem:flanks} (not its paired-seam clause), with one exception K\(\times\)S of spans at most d. Supply an interval I of length at least \(2q+d+44\), consisting of one exact extended source during Grow or of one ordinary home after normalization, at its last fully retired strobe before S. Assume that its strict interiors are proper and that its evolution is clean from its latest scratch reset (a \(\mathsf{Temp}\) reset, or the fresh reset of a positive vote whose live decisions keep its proper incumbents outside K) until S. Times are measured from the start of S. Put \(t_c=2d+9q+40\); the selected pass is the first full B-growth pass after \(t_c\), or the remainder of a pass then within its first ten growth instructions. Stop at normalization. A ballot during the response does not stop it, provided that (i) every ballot decision differing from the native one lies in K, (ii) after cutting the pre-ballot old occupancy at K, each remaining old panel of length at least \(q-1\) loses at most one site per side and no interior site across the ballot, (iii) scratch is freshly reset at it and the handle cadence continues through it, and (iv) at most four birth opportunities occur between the start of S and the first complete clean Purge that follows both S and every non-native ballot decision, and the last of these exceptional dates plus \(4q+14\) is at most \(t_c\). Otherwise stop at the ballot. A reopening of the same column K at a later ballot after a completed restoration is a second wound, to which the lemma applies with the pre-ballot interval as I and the ballot's non-native decisions in K as the exception. To the first complete clean Purge the bounds are \(25d+100\) new or reassigned positions and \(3d+12\) components, with a continuously held 24-site old run in every surviving component of length at least \(q-6\). A persistent internal cavity has width at most \(B+30q+40\) and stays inside K enlarged by \(2B+40q+40\) on each side. Within \(2d_0+10H_0<w^6\) ticks after the start of S it is filled or a distal flank disappears, provided that time remains before the stopping boundary. On a truncated response only localization is asserted. During Grow, every run of \(q+1\) old proper positions with ordinary addresses, disjoint from K, carried continuously through the latest reset and trimmed by one site at each end, loses no site for the rest of that Grow: no Purge with arbitrary exterior, temporary retirement, native attack or true Grow withdrawal removes any of its sites. A run of extended addresses loses at most one outer layer per true Grow withdrawal. The same conclusion holds in ordinary precommit with no ballot during the response, using in-home bids and at most three real Shrinks in place of the Grow losses.
\end{lemma}
\begin{proof}
The structural facts used below are those of Lemma~\ref{lem:transition}; the only clearance call is the held-wall part of Lemma~\ref{lem:flanks}.

\paragraph{Count to the first clean Purge}
Fix one exact source and clock, before its ballot or a change of address convention; in ordinary work this is one ordinary home. Immediately before S its actual matching occupancy is one interval, possibly empty. Let \(t_E\) be the end of the first complete initialized Purge starting after S. With \(H_0=2q+2\) an instruction occurs only once per \(H_0\) ticks. Starting after the last exceptional transition one waits at most \(H_0- 1\) transitions for the first clean initialization and then takes \(2q+1\) transitions through its last killing tick. Including S gives at most \(d+4q+2\) transitions. Since \(d<q\), a sequence of four distinct instruction dates (first and last at distance \(3H_0\)) cannot fit entirely inside these transitions. Keeping four opportunities is conservative even if both bounding rows are counted.

During this time let k be the number of canonical recruitment opportunities elapsed; for purposes of the proof count a full instruction even when it does not recruit. Apart from K enlarged by k positions on both sides, each of the two remaining half-lines meets at most one interval of the actual phase. At \(k=0\) initially this follows by intersecting the single initial interval with the two rays. Consider one future update. On a nonbirth tick an old interval within a ray can lose only endpoints. Its strict interior has two matching neighbors, hence is proper by the lower nominal triple property, computes its true update, and persists when the native rule holds it. Purge kills only at a right break, Shrink only at a zero face, \(\mathsf{Temp}\) retreat only at a zero face, and a true Grow withdrawal only at a zero face (Lemma~\ref{lem:transition}). On a birth tick clip each ray by one additional inner position. A child of a site outside the old ray cannot reach the new ray. Descendants of the one old ray interval still constitute an interval, possibly empty. An old interior has a proper positive hold. If a previously exposed tip with an old next-inner neighbor issues an actual outer child, then the old tip, its inner neighbor and the child all have matching nominal outputs (the first two have actual old served states and compute positive native holds). The middle of this next nominal triple is therefore proper and served. A singleton with children on both sides has the same argument; with only one child its descendant set is connected whether or not the old tip persists. A right- or left-end loss with no outer child adds no component. At the old ray boundary, a previously excluded parent's child can occur at most in the one newly clipped position. This proves the induction.

Inside the enlarged band there are at most \(d+2k\) occupied coordinates, giving at most this many components there and at most \(d+2k+2\) components overall. For \(k\le 4\) this is \(\le\)\(d+10\)\(\le\)\(3d+12\). On each of the at most four opportunities a clean birth can only adjoin one of the at most two exposed ends per canonical component at the previous row. At most \(8(d+10)\) assignments are therefore clean new births, counted with multiplicity. All remaining forged canonical sites, or lost-and-recovered canonical sites whose recovery is forged, occupy one of d coordinates. At most \(9d+80\le 25d+100\) distinct coordinates in total can be newly attached or forged. If a surviving component at \(t_E\) has length at least \(q- 6\), delete these coordinates. Deleting k'\(\le\)\(25d+100\) coordinates leaves at most k'+1 runs. If every run had length at most 23, the total initial size would be at most k'\(+23(k\)'\(+1)=24k\)'+23. The inequality \(q- 6>24(25d+100)+23\) holds for \(w\ge 100000\) by division by w (right-hand linear coefficient 2400). Hence some 24 consecutive original coordinates in the component have survived continuously since before S. Each lost original site can only have returned by one of the counted assignments.

A canonical fragment with no outside flank at the start of that clean Purge is erased if its whole strong-link component has \(\le\)q sites. Along any rightward string of strong links a surviving \(\mathsf{Cn}\)=1 after q propagation ticks requires q strong predecessors on that interval: unfold the copies of positive left \(\mathsf{Cn}\) one tick at a time. The finite string lacks such predecessors for every site if its length is at most q. No new site is born in the Purge. Starting from the resulting zero flags, the rightmost remaining site of each component has a weak or broken right link and is removed at each killing tick. Optional lower omissions only shorten the strings. This proves extinction within q killing ticks. A genuinely raw-only independent phase confined throughout to K expanded by four has width \(\le\)\(d+8\)\textless{}q. A newly joined flank instead belongs to the resulting whole component.

\textbf{Run protection.} A run of \(q+1\) old proper ordinary positions, carried continuously through the latest reset except for the particular phase boundary being analyzed, cannot lose a site to a Purge with arbitrary exterior. In q propagations the last site's q predecessors all lie in this run and were positive at reset. During the killing ticks induction from the old row preserves the last site by its positive \(\mathsf{Cn}\) and every earlier one by its strong right neighbor. This also applies starting within a partial sweep when the run was present for its clean prefix. A raw alteration only outside this run cannot place a zero signal inside except by the same leftward-to-rightward propagation of at most q links counted from the reset. A positive cut is followed by fresh initialization. Native attack never replaces incumbents; true Grow withdrawal cannot remove an ordinary address. Originally unmarked proper positions cannot die in a temporary retirement. These assertions are simultaneous invariants of the full proper run, via lower proper-successor service. They prove protection throughout Grow of each original \(q+1\) run of old proper ordinary positions away from K, after trimming one site for initial endpoint properness. If J hits an extended arm, one true Grow may remove one outer layer of a run of extended addresses; trimming three layers per exposed face in the short interval considered below pays up to three true Grow withdrawals. Each still existing strict interior has a nominal matching triple and is therefore proper. On the native step before a \(\mathsf{Temp}\) reset, old strobe positions can still be marked from the preceding completed retirement. No further temporary retirement happens before the next unconditional reset, so these marks give no additional loss at old positions outside K. If J hits immediately after a positive vote, the vote resets all scratch and keeps the proper ordinary run. If it straddles the vote, the run had a common continuous ordinary history through the vote and the fresh first Grow reset.

\textbf{Bookkeeping at a single wound.} Write \(H_0=2q+2\) for the handle length. From the last exceptional transition to the end of one full clean canonical Purge takes at most \(4q+2\) ticks. The same bound applies independently to the first full clean Purge on each displaced grid; allow d more ticks of exceptional rounding and four previous rows for removing nominal ghosts. For definiteness put the clearance time after the start of J at \(t_c=2d+9q+40\), an upper bound on these waits even if a terminal reset restarts a purge at a phase boundary, since Purges do not straddle the boundary on the fixed clock. If a reset restarts, the nominally preceding handle still ends at its proper instruction, and initialization starts at the very next handle. Take the first full B-growth pass after \(t_c\), or the remaining part of a pass already in its first 10 growth instructions at \(t_c\). There are at most ten skipped instructions in this latter case. To see that no full extra growth pass occurs between the original interrupted cycle and this selected pass, note that \(t_c<5H_0\) for \(w\ge 100000\) and begins at the start of J: even when J begins at the very last instruction before a cycle boundary, it cannot consume eleven growth instructions in the next cycle. If it begins later than ten instructions into the interrupted pass, the next pass is selected. If it begins in a cleanup or idle part, one simply proceeds to the next pass, paying at most the same ten initial opportunities. The selected pass ends, and the following complete Purge ends, in less than \(2d_0+10H_0\) after J. This is \(<w^6\) for \(w\ge 100000\). At most three Main instructions intervene even when counting both ends. Close to normalization use only the valid prefix of this count and stop at the label change.

Take \([u,v\)'\(]=I\) at the last fully retired strobe preceding J in this phase. If the wound coincides with an incoming vote, use the preceding ordinary retired interval, followed by the ballot on its actual old occupants. If the raw was materially earlier and K reopens at the ballot after restoration, use the then actual pre-ballot interval for the second wound, with no new column. Before the selected growth opportunities begin, genuine descendants outside I may have advanced at most \(B+14\) on either side: at most the one interrupted B-probe pass, up to ten initial opportunities of the selected pass, and at most three Main births. Only the one interrupted pass is charged as a full pass in this count. Any part first forged in K and not advancing from a genuine donor either dies at its first full Purge or attaches by then to a genuine component; its amplitude before that Purge is \(\le\)\(d+8\). Accordingly no enduring member extends beyond \([u-B-d-22,v\)'\(+B+d+22]\) at the selected start. A remote raw-only group cannot certify a more distant old tip here. If J is spent early in the previous ordinary work, ordering starts anew at the genuine restored interval.

Write \(K=[a,b]\). Where K cuts two original pieces each of length \(\ge\)\(q+22\), there is a \(q+1\) fixed old proper subrun on each side reaching within 10 of K. For example use \([a-q-10,a-10]\) and \([b+10,b+q+10]\); when K is inside I these are wholly inside the corresponding piece after discarding one old extreme. The ten-site margin covers one initial omission and three possible native true Grow losses at each facing edge. Apply the run argument at each purge between these operations. Then the between-flank gap is at most \(d+20\). If the left piece is shorter (including length zero when K strikes the left tail outside I), use the first right protected run starting at \(\max (u+10,b+10)\), whenever the right old piece is present. In this case \(b\le u+q+22+d\) unless K is disjoint far to the right, which is handled by the opposite orientation. The distance from any lasting distal left-hand member to the fixed run is at most \(B+q+2d+55\), by the preceding support bound. Exchanging u with v' and left with right gives the same estimate for a short right piece. The argument uses \(q+1\) strong links ordered \emph{left to right} for both protected runs; it does not reflect the orientation of \(\mathsf{Cn}\). Both old sides cannot be short when K intersects the macroscopic old interval. When K misses I and also misses its \(B+14+d+8\) neighborhood, it cannot hit a genuine descendant during this time and produces no genuine split. Thus any persistent between-flank cavity has length \(\le\)\(B+30q\)+40, since \(q+2d+55\le 30q+40\). Its entire interior stays inside K enlarged by \(2B+40q+40\) even on partial earlier trials, since each relevant original fringe is within \(B+q+2d+55\) of K, and one incomplete trial moves a front by \(\le\)B. For a deep K use the two protected neighboring runs directly. This argument allows a genuine, macroscopically long far part outside K; only the missing interval is localized.

\textbf{Durability and actual refill.} Stop at first reunion, extinction of a substantiated side, or the extra full Purge after the chosen pass. The cone count to the first clean canonical purge is \(\le\)\(25d+100\) new or reassigned positions and \(\le\)\(3d+12\) components, by the count just proved, which uses only intra-source bonds and genuine birth/nominal-next rules. After this purge a wholly new separated component of sustained width at most \(d+8\) is impossible, even if pinned at a grade-one or true-end contact. Otherwise a lasting flank of size \(\le\)q-2 on any intermediate row will have at most two new sites before its next complete Purge (there is at most one intervening instruction, giving at most one birth at each end). It will disappear unless a genuine join occurs. On the branch where neither stopping event happens, both flanks consequently have size \(\ge\)q-1 on the needed rows, including those whose next such Purge is the extra final check. At the first clean purge each contains 24 continuous old positions, by the pigeonhole proof above. The positions between these two initially stationary panels contained only the authentic source before J. Thereafter each bounding interval moves its endpoints by at most one per structural tick and has width at least q-1. Successive slices thus share \(\ge\)q-3 sites. A displaced lower lineage moving at most one cell per one of its own ticks cannot enter the cavity across the panel without occupying a simultaneously protected tile of the panel; shifted endpoints use two units of space and one extra row. Lemma~\ref{lem:flanks}'s buffered clearance argument applies to both enduring inner faces (eight spatial guards and the four prior rows are within the panels or cleared interior). It supplies mandatory births on the actual nine-cell stencils. A genuinely competing exact source planted between these panels has only K enlarged by the four birth opportunities of its own phase for backing and dies in its own full clean purge. If its normalization or ballot intervenes, its aggregate width is still \(\le\)\(d+8\), and the next initialized purity sweep treats that width. If one side disappears earlier the cavity no longer requires a directed fill.

The two inward-facing donors propose legal intermediate extended integers because both face the same exact source label and the legal integers comprise an interval. Where one vacant slot separates them the proposals agree on its address and age; priority is harmless. During the selected pass each lasting face makes at least \(B- 12\) genuine inward births (allowing two attempts blocked by an initially improper nominal incumbent), or joins. Its inward \emph{net} movement need not equal its birth count. The following larger allowance replaces B-3 in uses of this Grow estimate. At the start of the clean pass all but its two initially exposed endpoints are proper, and each genuine birth thereafter is proper. Thus at most two optional losses can ever occur in that flank during the pass. On a fresh Purge with no such loss, a component of size \(\ge\)\(q+1\) stays intact by the strong padding argument, while one of size \(\le\)q disappears. If any of the at-most-two optional losses occurs within a Purge starting with \(s\ge q+3\) sites, all but the two possible losses persist: the rightmost original survivor has original index at least \(s-3\ge q\), and a new left break caused by a single loss of the old left endpoint cannot propagate farther in q ticks than original index q (if it disappears on reset before propagation starts). In the latter extremal case \(s=q+3\) and loss of \emph{both} original endpoints still leaves rightmost original index \(q+1\), outside that signal. In detail if the left original endpoint is removed at initialization, propagation on the shortened string starts with reset ones; its rightmost after loss of the other endpoint has at least \(s-3\ge q\) old predecessors in the remaining run. If the left loss occurs later, the zero front can only move less far. If no left loss occurs the zero front reaches at most original index q-1. On a Purge with \(s\le q+2\) and an optional loss, the enduring branch leaves at least q-1 on its final row, hence the total sites removed in that Purge is \(\le\)3. At most two Purges can be exceptional in this accounting. Allow two additional optional losses on instruction dates for a conservative bound of eight sites lost. Hence each still separate enduring inner face gains at least B-20 inward sites during the chosen pass. The two gains exceed the initial hole bound, since \(2(B-20)-(B+30q+40)=98q-80>0\). There must be a join or earlier extinction. After joining, the true internal bonds have grade two, so the temporary cleanup cannot reopen this genuine hole. After distal extinction, a new detached genuine piece would require another pulse or an authentic donor already beyond the vanished piece; neither exists for this isolated source. A side whose only long connection was to \emph{another} exact owner is handled by its own strong backing at that owner boundary and the same alternative. With no time before normalization for the full pass, retain only the proven localized-wound bound and restart ordinary geometry on normalization.

\textbf{A ballot inside the response.} Under conditions (i)--(iv) of the statement a ballot is a no-birth step whose non-native decisions lie in K. By (iv) the count to the first complete clean Purge uses at most four birth opportunities, and every non-native coordinate lies in K, so the bounds \(25d+100\) and \(3d+12\) hold as above; the last exceptional date is cleared by \(t_c\), so the same pass is selected. By (ii) a panel loses at most one further site per side across the ballot, which the ten-site margin pays together with one initial omission and three withdrawals or Shrinks; by (iii) no old temporary mark crosses it. The remaining steps apply unchanged. The issued-source application verifies (i)--(iv) at a positive vote.

For the ordinary instantiation, there are no true Grow births or withdrawals. Each of at most three real Shrinks removes at most one layer on either exposed face, and an initially improper tip may additionally be omitted; the same ten-site reserve pays for these losses. All other Main instructions here are nonrecruiting. In a trapped internal cavity both inward ordinary bids are legal because its addresses form an interval in \([0,P)\). If one flank disappears, the gap is exterior and no two-sided closure is asserted. No same-rank issuing conclusion is used.
\end{proof}

\begin{lemma}[The ordinary damage and recovery bounds]\label{lem:ordinary-response}
In the ordinary precommit calculation from the lower-rank hypotheses, a home satisfies the erased branch or the surviving backed-interval branch of Lemma~\ref{lem:ordinary-start}. For a survivor the single interruption costs one localized response of length at most \(w^6\). At clean retired strobes outside it the stage-zero and stage-one end budgets are \(3A+3\Delta\) and \(6A+3\Delta\), respectively. By \(s_c-w^6\), either authentic material is extinct or \(C=[12A,P-12A)\) is covered outside the one local wound through pre-ballot work. The wound is confined to K enlarged on each side by \(2B+40q+40\). To the first complete clean Purge there are at most \(25d+100\) new or reassigned positions and \(3d+12\) components; every surviving component of length at least \(q-6\) contains 24 consecutive old positions held continuously since before S. An unmarked contiguous run of \(q+1\) proper old incumbents inside one ordinary home and disjoint from K survives a full initialized Purge with arbitrary exterior states, and also a Purge already in progress at the first exceptional transition if the run existed through its clean prefix. Such an unmarked run of \(q+5\) sites retains its middle \(q+1\) through any interval containing at most two real Shrinks. The erased branch can have only its bounded column remnant until an initialized Purge.
\end{lemma}
\begin{proof}
All times below count child updates, with S a set of consecutive exceptional transitions and K its interval of child coordinates. The common enlargement for shifted-grid queries is applied before using \(d=4w+8\). Assume the earlier-rank interpretation of the native rule, including proper survivors, proper births, and the nominal triple rule at both the old and new row.

\paragraph{1. Counts to the first clean Purge}

Fix one ordinary home and clock, before its ballot or a change of address convention. Immediately before S its actual matching occupancy is one interval, possibly empty. Let \(t_E\) be the end of the first complete initialized Purge starting after S. With \(H_0=2q+2\) an instruction occurs only once per \(H_0\) ticks. Starting after the last exceptional transition one waits at most \(H_0- 1\) transitions for the first clean initialization and then takes \(2q+1\) transitions through its last killing tick. Including S gives at most \(d+4q+2\) transitions. Since \(d<q\), a sequence of four distinct instruction dates (first and last at distance \(3H_0\)) cannot fit entirely inside these transitions. Keeping four opportunities is conservative even if both bounding rows are counted.

During this time let k be the number of canonical recruitment opportunities elapsed; for purposes of the proof count a full instruction even when it does not recruit. Apart from K enlarged by k positions on both sides, each of the two remaining half-lines meets at most one interval of the actual phase. At \(k=0\) initially this follows by intersecting the single initial interval with the two rays. Consider one future update. On a nonbirth tick an old interval within a ray can lose only endpoints. Its strict interior has two matching neighbors, hence is proper by the lower nominal triple property, computes its true update, and persists when the native rule holds it. Purge kills only at a right break, Shrink only at a zero face, and \(\mathsf{Temp}\) retreat only at a zero face. On a birth tick clip each ray by one additional inner position. A child of a site outside the old ray cannot reach the new ray. Descendants of the one old ray interval still constitute an interval, possibly empty. An old interior has a proper positive hold. If a previously exposed tip with an old next-inner neighbor issues an actual outer child, then the old tip, its inner neighbor and the child all have matching nominal outputs (the first two have actual old served states and compute positive native holds). The middle of this next nominal triple is therefore proper and served. A singleton with children on both sides has the same argument; with only one child its descendant set is connected whether or not the old tip persists. A right- or left-end loss with no outer child adds no component. At the old ray boundary, a previously excluded parent's child can occur at most in the one newly clipped position. This proves the induction.

Inside the enlarged band there are at most \(d+2k\) occupied coordinates, giving at most this many components there and at most \(d+2k+2\) components overall. For \(k\le 4\) this is \(\le\)\(d+10\)\(\le\)\(3d+12\). On each of the at most four opportunities a clean birth can only adjoin one of the at most two exposed ends per canonical component at the previous row. At most \(8(d+10)\) assignments are therefore clean new births, counted with multiplicity. All remaining forged canonical sites, or lost-and-recovered canonical sites whose recovery is forged, occupy one of d coordinates. At most \(9d+80\le 25d+100\) distinct coordinates in total can be newly attached or forged. If a surviving component at \(t_E\) has length at least \(q- 6\), delete these coordinates. Deleting k'\(\le\)\(25d+100\) coordinates leaves at most k'+1 runs. If every run had length at most 23, the total initial size would be at most k'\(+23(k\)'\(+1)=24k\)'+23. The inequality \(q- 6>24(25d+100)+23\) holds for \(w\ge 100000\) by division by w (right-hand linear coefficient 2400). Hence some 24 consecutive original coordinates in the component have survived continuously since before S. Each lost original site can only have returned by one of the counted assignments.

A canonical fragment with no outside flank at the start of that clean Purge is erased if its whole strong-link component has \(\le\)q sites. Along any rightward string of strong links a surviving \(\mathsf{Cn}\)=1 after q propagation ticks requires q strong predecessors on that interval: unfold the copies of positive left \(\mathsf{Cn}\) one tick at a time. The finite string lacks such predecessors for every site if its length is at most q. No new site is born in the Purge. Starting from the resulting zero flags, the rightmost remaining site of each component has a weak or broken right link and is removed at each killing tick. Optional lower omissions only shorten the strings. This proves extinction within q killing ticks. A genuinely raw-only independent phase confined throughout to K expanded by four has width \(\le\)\(d+8\)\textless{}q. A newly joined flank instead belongs to the resulting whole component.

\paragraph{2. Fixed backing in the old genuine interval}

An unmarked contiguous interval of \(q+1\) proper old incumbents inside one ordinary home, outside the damage, survives a full initialized Purge with arbitrary exterior states. Induct through its propagation: its last site has a length-q chain of strong left predecessors within that interval. At the end of propagation its \(\mathsf{Cn}\) is one. It cannot die during the killing ticks, while each earlier site has a strong right link to the next and remains. A Purge already partly executed at the first exceptional transition has the same conclusion when the old interval existed through its preceding clean prefix. Preceding clean birth instructions, old proper exact holds and a reset on an already occupied interval do not mark these original sites. On a \(\mathsf{Temp}\) retirement their own original mark remains zero. On an ordinary Shrink either exposed edge of this fixed interval may be lost, costing at most one from either side on one such update. Hence for a time interval containing at most two real Shrinks, any original piece of \(q+5\) such sites retains at least its middle \(q+1\), including when a Shrink occurs before the first clean Purge. Original strictly inner sites at an old single interval are proper by lower triples; discarding one old endpoint on each side gives the needed originally proper section. This proves survival of deep old backing without using recovery.

\paragraph{3. A cavity opening onto a true end}

A vanished distal flank leaves an exterior gap of a single surviving interval. There is no requirement that a single donor close distance \(B+30q+40\) in one B-opportunity pass. Ordinary proposals are in-home, so the adjacent home across a true end is never a legal donor for the remaining in-home vacancies. If the gap is not relatching and no other canonical material exists ahead, the authentic strobe end moves inward by at least one on each eligible Shrink: a provisional trial ending strictly short of the true end has a zero face there and is fully removed by B retreats. Failure to advance on an attempted birth only reduces trial length. This is the open-end episode counted in Lemma~\ref{lem:seed}. A positive true seam already fully occupied in both homes has the stronger fixed-backing repair discussed in Lemma~\ref{lem:flanks}.

\paragraph{4. Relative size of the threshold constants}

With \(B=128q\) and \(\Delta =1024q\), the separate small-run test \(16B+2d+24\) is strictly greater than \(\Delta\) and smaller than 3\(\Delta\). A guarded strip beyond a threshold uses the 3\(\Delta\) comparison in the guarded erasure calculation. The tighter \(\Delta\) in the survivor calculation bounds only the extra displacement of an authentic restored frontier during the short local response.

\paragraph{5. Guarded erasure with the pulse at an arbitrary point of the timetable}

Use fully retired ordinary strobes, also between two successive real Shrink calls. In particular one can select a retired strobe at distance at most \(d_0\) from any ordinary fine date. Unless explicitly called an inspection, such a strobe need not be followed by a real Shrink. On a fully retired interval, a missing end at positive depth cannot be advanced permanently toward its true edge unless the edge relatches or another matching component ahead of that end has been encountered. The foremost site of the whole population on either side of a permanent sterile strip never has a component ahead of it on that same side. A trial of \(\le\)B marked sites into that open gap, stopping short of a true-end exemption and any other matching component, is erased by the B no-birth retreats: at each step its exposed forward site has grade zero and is nonexempt. This proves the assertion even for failed attempts or if an honest lower end is independently omitted.

Let \(W=4B+4\). A W-site vacant strip at a fully retired strobe, entirely inside the ordinary home, disjoint from K and from the two true endpoint sites, is a permanent barrier. At most B marked sites enter per side on a pass. Its central \(W- 2B>0\) sites remain vacant; every still entering tip is nonexempt and withdraws by at least one per removal until its entrance is empty again. Other phases yield zero grade to its advancing tip unless they count as this matching population. Lower omissions can delete more sites but cannot create an unmarked new site inside. Thus its middle stays empty at every fine date and the entire band is empty again at the next strobe, including during the pulse elsewhere. Across a transition from one ordinary stage to the next the same true homes, reset and trial restrictions apply.

Here is a quantitative interruption charge that does not require a macroscopic old interval. At the last retired pre-S strobe let f denote the nearest position of matching material on one side of a permanent sterile band, ordered by distance toward that band. If no old material exists on that side, only a raw-only family of width \(\le\)\(d+8\) is possible up to the first complete clean Purge, and vanishes there. Otherwise at the end of that Purge genuine descendants can be farther toward the band than f by at most \(B+d+16\). From the old strobe up to S there is at most one partial ordinary pass, giving \(\le\)B steps beyond the former largest coordinate; if the chosen strobe is just before the previous ordinary handle and a real Shrink intervenes it can only subtract sites. From S to Purge end there are at most four additional birth opportunities on this clock. Descendants with no raw assignment in their ancestry therefore lie at most \(B+4\) beyond f. Those with a raw ancestor and now outside that previous half-line lie in K enlarged by four. If this last band has no contact with any genuine old-descendant site by the time of the complete Purge, its entire width \(\le\)\(d+8\) dies. For a contact its far side can be at most \(d+8\) farther than \(f+B+4\), plus one for adjacency. A pre-existing old group at a different location behind f does not increase this maximum. This proves the stated \(B+d+16\) bound, whether or not the pulse creates numerous components in K. Through the first following retired strobe pay at most another B for the remainder of a pass. Thereafter the foremost frontier can only withdraw again at the clean strobes. Between the old and the new strobes at most three real Shrinks have been missed: the elapsed time is \(\le\)\(2d_0\)+\(d+4q\)+4\textless{}\(3d_0\), and true Shrinks are \(d_1=q\) \(d_0\) apart. Use the symmetric argument for a leftward order; this is an order calculation on front positions, while the actual small-component Purge proof still kills from the right.

Suppose the left endgap at an inspection (or any retired strobe in the eligible stage) exceeds \(\theta +3\Delta \), where \(\theta =3A\) or \(6A\). In the vacant part strictly between \(\theta\) and the actual old population choose a W-site strip not meeting K or a true endpoint. There is room for two disjoint W strips separated by \(d+2\) and both avoiding the true endpoints since \(3\Delta >2W+d+4\). K cannot meet both. No old material lies to its left. There can only be an independent, first-Purge-extinct colony there after S. Every true Shrink with threshold \(\theta\) removes the leftmost remaining site to its right, as its zero left face is at depth \(>\theta \). The sterile band prevents any probe from permanently latching at the left true end or encountering another matching site ahead. On clean strobes the ordinal position of the nearest surviving authentic front increases by \(\ge\)1 per eligible real call; disappearance of a component exposes a farther front and only speeds this up. Counting from the start of the \(\theta\) stage, at most \(P+2B+d+19\) eligible inspections suffice even if S falls in its middle, by the preceding interruption charge \(2B+d+16\) and three potentially missed calls. This is \(<1.1P\); there are \(2P\) such calls in a whole stage. When the large gap was already seen in the preceding lower-threshold stage, the same band remains empty and all authentic sites on its populated side still have depth \(>\theta \) when the \(\theta\) stage begins. The independent late raw-only case does not produce a genuine late comparison. The reflected end uses the depth \((P- 1)- x\) in the eligibility test and the rightmost remaining authentic front facing the strip. If the sole survivor is exactly at the opposite true endpoint, the face toward the strip is still internal and eligible.

For an initially arbitrary normalized home at retired strobes, each interior gap surviving n clean stage-0 inspections contains at least n old vacant coordinates in the following precise size sense. At an inspection a still-internal gap has unjoined, retired banks: no clean birth creates a new internal gap, so trace one preexisting internal gap inside its predecessor whenever two gaps merged. At each true Shrink its exposed bank loses its inside tip (or the entire bank). A recurrence on the ancestral gap's length gives the next internal gap at least one longer than that chosen gap. If the gap becomes exterior, it can never become an internal gap again without an independent newly born component, unavailable on the clean ordinary evolution. Therefore after \(20A\) clean real inspections the home is one interval or has a gap of length \(>20A\). Choose a W-site sterile band within that gap avoiding true ends and K as above. On each side the matching front facing this strip erodes by \(\ge\)1 per remaining \(\operatorname{Shrink}(0)\) call, using the interruption charge above on both sides concurrently. More than \(3P\) stage-0 calls remain, so both sides are gone in at most \(P+2B+d+19\) of those calls, except short immediate raw-only debris. If S intersects the initial \(20A\) inspections, restart once after S and a clean ordinary reset; \(50A\) \(d_1\) covers both attempts and the boundary spacings for \(w\ge 100000\). If the sole surviving interval on that early clean row misses a true end by more than \(16B+2d+24\), there is again room for two disjoint W candidates separated by \(d+2\) between that end and its front; a sterile strip forces the same stage-0 extinction. Thus on a branch capable of authentic late occupancy the simultaneous early row is a genuinely long interval and both \(\theta\) thresholds lie strictly in its proper interior. A site outside K which was a proper incumbent stays proper on every subsequent live native hold; an honest re-entry after absence is a proper birth by the previous-rank proper successor clause.

In the later stage-0 interval, \(P+2B+d+19\) calls in the following whole stage 1 dispose of any authentic population whose endgap ever exceeded \(3A+3\Delta \) at a genuine retired strobe. Likewise a stage-1 depth above \(6A+3\Delta \) is excluded by the whole following stage 2. This establishes the guarded-threshold bounds on every surviving clean retired strobe without presuming a finished local refill. In particular prior to an unspent S in stage 0 or 1, either one of these erasure tests applies or the two true ends enclose an old interval of length at least \(P- 12A- 6\Delta - 2\). A virgin late pulse in stage 2 or computation follows a wholly clean stage 1; an initially active left stage-1 gap \(p>3A\) cannot latch because \(p>B\) and grows each Shrink until extinction, and an inactive \(p\le 3A\) cannot first lose a proper occupied \(3A\) threshold spontaneously. In detail, while the interval is long, Purge misses the strongly backed last site, the endpoint at depth \(\theta\) has no eligible outward Shrink, and a genuine strobe end at positive depth is unmarked (a nonjoining B-long outward trial has completed all its removals). An uncharged omission there would require improperness, excluded at an initially proper or honestly reborn threshold site. If either end was already active, once the old interval becomes short its continuing exposed active face only accelerates its exhaustion; no clean outer trial from its interior can cover an end gap already exceeding \(3A>B\). Therefore in this clean-before-S case either the home is already extinct or its stage-1 exit has both end depths at most \(3A\). The same clean comparison applies during stage 2 up to the first later pulse.

\paragraph{6. Local response on a survivor and final-stage timing}

Apply Lemma~\ref{lem:local-wound} to the old macroscopic ordinary I just obtained; its length is at least \(P-12A-6\Delta-2>2q+d+44\), its strict interiors are proper by lower triples, and its evolution before S is clean. The issued-source application is given later in Lemma~\ref{lem:attachment}; that is a pointer, not an input to the ordinary calculation. There are no true Grow births or withdrawals in these ordinary stages. At most three actual Shrink instructions occur in the response. Each can remove one layer on either exposed face, and an initially improper tip can additionally be omitted. The ten-site reserve on each face pays for these losses; all other Main instructions are free of structural attack or withdrawal. Both in-home inward proposals in a trapped internal cavity are legal, since the intervening ordinary addresses form an interval in \([0,P)\). No ballot occurs on this ordinary response if it is used to establish geometry before computation. For a pulse near the computation ballot, the ordinary assertions here stop at its pre-ballot row with the local wound still charged. The two event horizons and partial pass in that calculation yield at most \(B+30q+40\) sites of persistent internal cavity, or disappearance of the distal piece. The successful-pass estimate uses \(B- 20\) instead of \(B- 3\); \(2(B-20)-(B+30q+40)=98q-80>0\). If one side is lost and the hole is exterior, use the earlier true-end cavity calculation. Deep surviving old strips have no interior loss, and the discrepancy before restoration is confined to K enlarged by \(2B+40q+40<.2A\). The full response uses \(<2d_0+10(2q+2)<w^6\) ticks and at most three true Shrinks. To be precise, the portion of a fresh pass used after its uniform clearance may skip up to ten initial growth instructions, already included in \(B- 20\). This local argument does not import any same-rank strong vote or any Grow-specific conclusion into an ordinary precommit calculation. The guarded-erasure and survivor steps use the held-wall implication of Lemma~\ref{lem:flanks} and the supplied-interval lemma only when their wall and old-interval premises have been established.

At either original strobe end of I, a surviving old proper strip begins within \(q+d+40\) sites of that end throughout the response on the macroscopic branch. Indeed choose two candidate strips of length \(q+12\), separated by \(d+1\), strictly inside the old interval beginning one site from that end. K can touch at most one. Trim four old positions at each end of the surviving choice for potential Shrinks and an unprotected extremity, leaving more than \(q+1\) old sites. The nearest preserved point has depth at most \(q+d+18\) beyond the previous old endpoint. Conversely, by the first-Purge contact calculation in step 5, a \emph{persistent} edge cannot shift outward beyond the old retired edge by \(>B+d+16\) at first cleanup. Up to the completion of the chosen pass there are at most two full or partial additional B-growth passes and up to three Main steps; all ordinary Main steps here are nonrecruiting. Including one more B safety allowance, the absolute displacement of a restored strobe end beyond the extrema of the original I is \(<4B+d+40<\Delta \), with the inward loss paid by the preserved strip. If a true-edge latch occurs during this interval the newly latched endpoint must also lie within this paid distance, since each recruited layer is counted. These assertions are for the authentic strobe frontier, not for a distant unjoined raw-only satellite.

For the final threshold test choose one true end, put \(h=3A\), h'\(=6A\) in its directed depth coordinate, and separately test whether K lies within 2\(\Delta\) of h'. Since h'\(- h=3A>5\Delta \), this cannot also put K within 2\(\Delta\) of h. Work on the potentially surviving macroscopic branch above; if an active facing population instead disappears only the immediate small raw alternative remains.

\begin{enumerate}
\item When S starts before the beginning of stage 2 and K is farther than 2\(\Delta\) from h', the site at depth h' was proper on the early long row. Any later authentic re-entry there was honest and proper, because the site is outside K. If both ends of the restored interval have depth \(\le\)h' at stage-2 entry they cannot first cross that bound: on a long interval at an inactive threshold its genuine outward Shrink is exempt, interior bonds prevent other removal, no improper omission is possible there and backing prevents purge. The two-end first-crossing argument uses length \(\ge\)\(P- 2h\)' until the first such crossing. If an end started at \(p>h\)', it faces a gap of width \(>B\); unsuccessful clean trials fully retire, each stage-2 Shrink removes that foremost member, and by the first \(P+2\) genuine calls the old population is gone. The other end's behavior cannot slow this one. If S itself straddles the stage-1/stage-2 transition or began fewer than \(2d_0+10H_0\) ticks before stage-2 entry and restoration continues into it, delete the \(\le\)3 initial interrupted calls and start this same clean test at restoration. Original properness at h' is preserved through the job or supplied anew by honest re-entry outside K. This leaves far more than \(P+2\) calls.
\item When K lies within 2\(\Delta\) of h', it is farther than 2\(\Delta\) from h. On a clean inactive stage-1 trajectory \(p\le h\) the proper h site cannot fail uncharged. If S comes while this boundary is still inactive, h lies in an old protected strip directed inward of length \(q+12\) beyond h disjoint from K. Through the local response the left face at h remains exempt in stage 1, its inward \(q+1\) backing persists across the \(\le\)3 Shrinks by the reserves, and if h was a strict interior it either remains so or becomes that exempt exposed endpoint. The properness of its old incumbent follows at every live step since h is outside K. When the interval is restored, \(p\le h\) still. If \(p>h\) already at stage-1 entry, clean calls increase p until S. The interruption cannot relatch it across h. For \(p\ge h+\Delta \) at the last pre-S strobe, the permanent outward gain proved above is strictly \(<\Delta \). For \(h<p<h+\Delta \), K is \(>\Delta \) from the frontier and the \(\le\)B partial probe from it. A raw-only patch here cannot contact that frontier: a birth history of \(\le\)four instructions before a complete Purge stays in K expanded by four, and no persistent independent short patch is possible. Nor can a change deeper than \(\Delta\) force a new inwardly supported outward latch from this frontier: there are true old proper sites of length \(q+12\) immediately inward of it, and the exposed depth \(>h>B\) forbids a completed true-end latch; a genuine vacancy ahead has no other old matching component. It therefore remains \(>h\). The same argument applies when S was already spent in stage 0: if h is occupied at stage-1 entry it is proper, and if the end is beyond h it dies through the clean stage. Accordingly an active stage-1 end facing a wide gap throughout pays at most \(\Delta\) extra retreats plus three missed calls and is exhausted within \(P+\Delta +4\) of the \(2P\) calls. On the surviving alternative this side has depth \(\le\)h at stage-2 entry. If damage is still being repaired over the boundary, the surviving old strip beginning within \(q+d+40\) of its former end gives depth \(<h+\Delta \), far short of h'. Thereafter at most the one uncharged loss of a previously improper exposed tip can occur before any fresh authentic endpoint birth, and that loss while backed cannot open a depth-h' breach. Thus a possibly forged h' site is never asked to retain an already marginal h' edge arbitrarily late in stage 2.
\item If S has not started before stage 2, the wholly clean stage-1 exit has \(p\le h\) on both ends, or has gone extinct. Until S any exposed proper h' site cannot be endangered from an edge of depth at most h plus one optional endpoint loss. The macroscopic old strips bounding inward setback throughout S put the new strobe depths below \(h+\Delta +2<h\)'. Their original interior extent is again much larger than q, and the at most \(d+8\) subsequent improper-tip omissions cannot move an edge from this depth to h'. Erasure of a remote deep piece has already been localized and repaired by the internal-cavity lemma. An S in computation instead encounters the same long old interval under no true Shrink and the same one local wound.

\end{enumerate}

Both ends undergo the same numerical test, sharing a single K rather than independent exceptional columns. In ruling out a spontaneous first crossing while applying the two tests, stop at the first crossing at either end, at which all protected opposite-side buffers are still macroscopic. If an end begins a stage active (\textgreater{}threshold), its exhaustion argument depends only on its own positions and remains valid when very little material remains. By the time of \(s_c- w^6\) either authentic material is extinct (temporary raw occupancy only), or \(C=[12A,P- 12A)\) is covered outside the one local wound through pre-ballot work. The long stage-2 erasure branches terminate by \(P+\Delta +10\) calls, leaving an exponential margin before the end of its \(2P\) calls. The two stage-0 and stage-1 bounds of step 5 and the clean interval persistence used for the nonfull-date count now apply to this latter alternative. The true-end one-sided case and the distinction between a grade-one bridge and an independent Purge backing are unchanged.
\end{proof}

A home is \emph{nonfull} at a fine date if at least one of its P positions lacks an actual matching served child.
\begin{lemma}[Number of nonfull dates and a common full seed]\label{lem:seed}
Work after the early stage-0 relaxation. Assume the survivor assertion of Lemma~\ref{lem:ordinary-response}: apart from one response interval of at most \(w^6\) fine dates, this is one genuinely contiguous ordinary interval at retired strobes, with each endgap \(p\le L=3A+3\Delta \) at each clean stage-0 Shrink strobe. Assume also that on a clean partial pass lasting through a strobe, a trial which fails to latch its true end returns to the previous exposed frontier, and that through any fully covered true end a \(q+1\) backing of previously proper incumbents is maintained on this branch. Then after the relaxation and throughout stage 0 the number of nonfull fine dates is at most \(100A\) \(d_1\). In particular \(N=200A\) \(d_1\) is a uniform search allowance. With the same hypotheses for up to six neighboring homes and the common forbidden dates below, a full common seed of \(5d_0\) consecutive ticks occurs.
\end{lemma}
\begin{proof}


Cut the calendar at the response interval, leaving at most two clean time intervals. Consider one true end in one such interval. After a genuine birth authenticating the true-end site the newborn is proper; it remains served while its native update holds. Its outward zero face is exempt from \(\mathsf{Temp}\) and Shrink, and its in-home link is positive. The \(q+1\) same-home true-end backing prevents an erasure from the tail of a Purge, even though the seam's wrap link itself is only grade one, by the last-site backing argument of Lemma~\ref{lem:ordinary-response}. Consequently if the true end is initially present and improper, that initial site may be lost at most once, and after any subsequent relatching birth there is no further clean loss in this fixed clean interval. An initially absent endpoint counts as one initial missing episode. Allow two missing episodes per end per clean time interval, including a partially open first or last episode.

Until relatching, any incomplete attack into the remaining true-end gap has only B marked sites along its nonjoining trial; each removal tick retires its newly exposed eligible tip, hence at a clean fully retired strobe this trial cannot decrease the old depth. At each \(\operatorname{Shrink}(0)\) with positive depth the end with its outward zero face is removed, increasing the depth by at least one before the following strobe, unless a true-end latch has since occurred. Thus successive fully internal strobe depths of a missing episode are distinct increasing positive integers \(\le\)L, so there are at most L such strobes. With \(d_1\) between strobes an episode including the two arbitrary partial spacings takes \(\le\)(\(L+2\))\(d_1\) dates. Fine rows temporarily covering the true end inside such an episode only reduce its nonfull count. Counting both ends, both clean intervals and two episodes gives \(\le\)8(\(L+2\))\(d_1\)+\(w^6\)+\(4d_1\), padding both exterior sides and both response boundaries by a strobe spacing. At \(w\ge 100000\), \(3\Delta /A=3072/w<0.031\); \(d_1=q(2A+4)(2q+2)>4w^7\) implies \(w^6<d_1\). Hence \(8(L+2)+5<100A\). A hidden interior gap would invalidate this argument; the genuine interval premise and the response charge are essential.

For up to six neighboring positive branches, let \(n_i\) denote their nonfull dates. Each \(|n_i|\le N\). For each common raw temporal interval requiring exclusion, enlarge it into the full forbidden date set, of total cardinality at most \(4w^6+20\) since the single immediate response has length at most \(w^6\), any associated delayed ballot response is confined to the same wound and has length at most \(w^6\), and adding a \(w^6\) preparation halo on each temporal side, with 20 rounding dates in total, covers all queries. Within stage zero itself there is no new ballot. Define \(b=6N+4w^6+20\). After the early cutoff take \(b+1\) adjacent disjoint trials, each of length \(5d_0+1\) integer dates. Every failed trial contains at least one distinct forbidden date (use its first such date). The union of six nonfull sets and the slabs has size \(\le\)b. Therefore at least one whole trial is full in every needed home. In particular, this costs \((b+1)(5d_0+1)\)=poly(w) dates, while the remaining \(\operatorname{Shrink}(0)\) timetable has exponentially many calls and \(d_1\ge 1\). A full trial of this length includes a \(\mathsf{Temp}\) reset and a complete Purge whose initialization follows that reset, because resets recur every \(d_0\) ticks, and a Purge takes \(2q+1<d_0\) ticks. For adjacent full interiors, the lower triples which straddle a positive seam put both seam endpoints in proper status during the seed interval. Coverage after a later raw interruption requires the paired repair of Lemma~\ref{lem:flanks}.

The common full trial supplies the paired input corridors and the six-home service calculation. The survivor and backing premises were established by the ordinary relaxation and damage bounds.
\paragraph{Discharge on the surviving stage-0 branch}

By Lemma~\ref{lem:ordinary-start} and steps 5--6 of Lemma~\ref{lem:ordinary-response}, away from the one response interval there is a genuinely single ordinary interval at each retired strobe and each of its two directed end depths is at most \(3A+3\Delta \). Its length is therefore at least \(P- 6A- 6\Delta >2q+30\). When it reaches a true end by growth the new end has an actual proper birth. Its near in-home \(q+1\) sites form the beginning or end of this interval; the true outward \(\mathsf{Temp}\)/Shrink exemption, positive inward link, old proper holds and rightmost q-Purge backing imply that an already achieved latch cannot be lost through a later clean retreat. For an initially full but not-yet-proper endpoint, allow its first free loss; strict old interior sites were already proper by triples and later legitimate births are proper. A previously nonjoining trial on an exterior gap that never reaches the true edge has no matching population ahead and regresses its still-marked facing tip at all B removals. The current single interval could not meet another matching component across that exterior gap. Thus precisely the clean-episode premises of the counting lemma are discharged by the preceding ordinary geometry. If the end was attained for the first time on the Bth attempt, the facing true-site exemption already applies at the very first retreat, and its positive inward link provides the other face. With a full seeded pair any ensuing seam interruption is treated by Lemma~\ref{lem:flanks}.
\end{proof}

\subsection{Transactions and classification}
The weak-description then strong-issue order follows G\'acs \cite[Lemmas~14.1--14.2]{gacs86}.
\begin{lemma}[Weak descriptions and the strong issuing cut]\label{lem:ballot}
Use the ordinary precommit alternatives and the lower-rank transition and regularity clauses for the owner and its three suppliers. There is a weak positive description by one common live-letter string on the central fifths, or a weak dead cut with actual support in the exceptional column. Three prior weak descriptions give independent parsing of the served suppliers and the strong cut of Proposition~\ref{prop:organization}\ref{org:issue}. The latter includes the noncomputing branch and a possibly halted true literal.
\end{lemma}
\begin{proof}
In each unaffected evaluation the full central three working fifths are occupied with scratch reset. The three LOAD sites and both neighbors on their actual shifts are lower-rank live matching triples. Their exact native step loads the true assigned operand at that position even from arbitrary pre-existing latches and program tape. In every such repetition the same three fixed geometric positions and parameter position are used. Conditional on the specified frozen environment the target F is therefore a single fixed owner-specific map, including on forged legal mutable states. Section~\ref{sec:environment} identifies the actual operands and proves localization.

\textbf{Geometric inputs.} A weak positive description means one common live-letter string on the central fifths off the common physical column, with the narrow support test specified below; a weak dead description means actual matching support confined to that column. The inputs are the late core of Lemma~\ref{lem:ordinary-response}, the common full seed of Lemma~\ref{lem:seed}, and the paired corridor of Lemma~\ref{lem:flanks}. Use the lower exact-or-dead transition, true lower birth/hold and proper-successor clauses. For the current centered transaction and its three suppliers assume the ordinary dichotomy, the late-date core C, and the continuous paired corridors of Lemma~\ref{lem:ordinary-response}, Lemma~\ref{lem:seed} and Lemma~\ref{lem:flanks}, wherever the full time interval of an operation misses the raw slab enlarged by \(w^6\). The empty ordinary branch has no matching material beyond the immediate \(5q\) allowance. A disturbance during ordinary work changes an old survivor's interior only in K enlarged by \(2B+40q+40\) (or the wider \(.2A\) used here); outside it the survivor endpoint ledger keeps C protected after the common seed. Work in the one rectangle K\(\times\)S in this entire buffered domain; the possible delayed wound at a subsequent vote is in the very same K. A complete good repetition spans its reset, input gather, immutable loads, computation, result writing and final output writes. The separation P idle Main instructions gives pairwise disjoint \(2w^6\)-padded transaction intervals and permits the stipulated margins before the first and after the last prior cut. No claim here assumes the current nominal output.

\textbf{Old carrier invariant.} Number fifths 0,...,4; the three inputs of each supplier are in fifths 1,2,3. In a previous weak positive description there is one common live-letter string z, possibly malformed, on all its central fifths off K, and \(C_*\) is served and matched off one common interval of \(\le\)\(2q\) containing its possible missing positions and any wrong tested letters. The narrow actual issuing cut localizes any such loss at K and its immediately adjacent short piece, or in the \(4q\) reserve outside \(C_*\). In the absence of a raw column there is no such defect. In particular the error sources for X in the central fifths have the one common physical column; this is an explicit part of the narrow weak cut proved below.

Any \(q+1\) consecutive old positions deep in \(C_*\) outside those exceptions and K, with one adjoining old match at both ends, are initially proper by lower triples. Their original scratch, including \(\mathsf{Temp}\), is reset at the preceding healthy ballot and their original addresses are ordinary. Through an entire Grow no true Grow withdrawal reaches those addresses, recruits ignore incumbents, and no Shrink is active. A complete Purge preserves the whole old strip: the rightmost \(\mathsf{Cn}\) remains one after its q internal strong predecessors, and all earlier sites retain strong right bonds in every killing tick. \(\mathsf{Temp}\) stays zero on old incumbents, even when nearby births are marked. The fixed-strip proof and lower proper-successor clause apply inductively through all handles to every old strip and are independent of the exterior. This preserves actual held X through the long Grow. Normalization keeps these ordinary addresses unchanged.

For an explicit later spatial margin set \(N=200A\) \(d_1\), \(D_{\rm tr}=50A\) \(d_1+(6N+4w^6+21)(5d_0+1)+w^6+2d_0+20\). This exceeds the upper bound from Grow-to-ordinary normalization to the selected common seed. For fixed w all these quantities are polynomial expressions in w; in particular choose w so that \(2(D_{\rm tr}+10A+10q)+d<P_0/2\) and \(12A+4q+D_{\rm tr}+10A+10q<P_0/2\). In ordinary times up to this seed, any middle-fifth cell whose \(D_{\rm tr}\)-neighborhood is entirely old correct and excludes K retains X: in each next step the strict interior of the previously held segment has two strong neighbors and is proper and held by the exact native transition. No clean interior death occurs in this interval, including in a partially completed Purge, because its old right bond is strong. Trimming one site on each side per step proves the claim. In particular all three central fifths except one physical interval K expanded by \(D_{\rm tr}+10A+10q\) carry the original z. Its width is \(<P_0\). The outer ends of \(C_*\) have more than \(D_{\rm tr}\) margin to the middle fifths. After the seed and through good reads every such old unhit position remains present by the earlier surviving-home ledger and the localized single-wound repair; a proper living incumbent does not rewrite X on a nonballot step. At any offset u, positions \(sP_0+u\) for \(s=1,2,3\) are \(P_0\) apart, so at least two \emph{of this supplier's coordinates at that offset} continuously carry z. A different u may choose another pair. When the preceding weak decision was D, matching occupancy at the cut outside K was empty. A subsequent matching X at such a site comes from a native birth, writing the playable neutral carrier at its own residue; thereafter it is again held or replaced by a true birth with the same neutral letter. A foreign surviving incumbent cannot switch into canonical matching status across the off-Grow normalization without the correct original clock and residue; at the preceding cut all off-K survivors had normal addresses. This proves the neutral or physically absent alternative.

\textbf{Paths and validity.} For receiver home h and supplier offset j in \{-1,0,1\}, let \(h_0=hP\) and source coordinate \(a=(h+j)P+sP_0+u\) and destination \(b_{\rm dst}=h_0+(2+j)P_0+u\). During supplier j/copy s transport each hop at step k has prescribed coordinate \(a+k\) sign(\(b_{\rm dst}- a)\), \(0\le k\le |b_{\rm dst}- a|\), on its assigned lane. The displacement is an integer multiple of \(P_0\) independent of u. At injection only a site with the programmed source residue places its X there; every uninjected lane writes message-empty. At zero displacement the destination takes the newly assigned local value. At a positive displacement count (in either orientation), the old upstream signal from the unique predecessor enters the next slot if the receiver survives that hop and their grade is positive. A missing link gives message-empty, and a newborn's mail fields start as message-empty. No stale destination value is read because receive unconditionally overwrites the indicated bank at that offset. Induction on k therefore gives X(a) on an unbroken path or message-empty on an interrupted one. Any injected garbage on a physically damaged hop belongs to the one affected repetition and is not included in either good read. Across computing homes every hop is within their joined corridor (including a grade-one wrap, which transmits); all paths between homologous data slots remain inside the core-to-core corridor. Every required write and hold is the native update of the actual supplied slot. For \(g=0\) the entire source fifth at a good injection is absent, so all these receives get message-empty or their birth/reset default. The centered evaluation itself is full on the required destination fifths, validator and gathering tape.

For \(g=1\) and a prior weak positive, the at-least-two surviving copies per offset yield z exactly at every offset after the local vote. If z were malformed, any candidate live decode outside its own exceptional interval would be contradicted at some offset: of three homologues one avoids both that interval and K. So malformed z gives \(x=D\), exactly the fresh parser output; a legal live z instead yields its serialized x. For a weak dead input, the proper parse of a neutral packet or absence is D. For \(g=0\) the supplier value y is D irrespective of preceding formatting. This proves that each supplier independently parses to its actual y in a good repetition, assuming its preceding weak description. Two good repetitions therefore evaluate exactly the same total native \(F(y- ,y,y+)\) on the fixed physical loads. The computation work tape and mailbox overwrites are freshly initialized in each repetition. Within output readout apply the same path induction with source \(hP+2P_0+u\) and each destination \(hP+sP_0+u\) for \(s=0\),...,4. A true dead result cannot insert a live letter on a broken path. Delivery in the middle three fifths is exact at all good writes.

\textbf{No backward assumption of a strong cut.} First obtain a weak description when the center has no direct exceptional update in its late ordinary computation, verification or ballot region (include a \(2w^6\) temporal margin before its first read and the entire actual home through its ballot). Incoming data from neighbors may change between all three evaluations. Each fresh full center still outputs one legal live packet with every letter live or all message-empty on its central fifth, so the central three fifths see the same generated letter at each homologue in each repetition. The weak liveness decision is their common majority of the three liveness bits; if live, the vote gives the same string z at every homologue. No validity is claimed for z. A repetition with a dead result can put only dead banks in any unhit shifted path; a direct event in another home cannot inject into the all-in-home output shifts. A direct event confined to an adjacent home has broad damage radius \(.2A\) and cannot touch C of this center from outside. For a center whose direct late ordinary work does meet J, the preceding periods of all three suppliers ended a full Grow/Shrink duration before that late work. Thus their late ordinary weak descriptions are all in the preceding clean-center case. Apply the two-good-repetition calculation just proved to this affected center, yielding in particular a weak description without assuming one for any arbitrary earlier predecessor. If S straddles the bottom of the current period or occurs in its early Grow, the present late-ordinary center is clean; if S straddles its final ballot, it is in the other case and the three earlier supplier work computations are clean. The one separated delayed footprint at a vote uses the same physical column, and its earliest read margins exclude it. Having proved weak descriptions for all internally queried prior periods by this one backward step, apply the good-read argument to the actual current period for its strong conclusion.

\textbf{Retaining the central votes.} In ordinary computation no true Shrink is active. On a late raw pulse take the original interval \([P_0- q- 10,4P_0+q+10)\) of proper held center slots just before the pulse; on an early pulse instead take the good complete writes after its recovery. Split the original interval on K. Any component of its complement containing a middle data coordinate has \(\ge\)\(q+3\) positions, or its short boundary piece lies entirely in the \(q+10\) margin. Such an old component has proper middle triples and had persisted through the initialization of a partial Purge. It survives all later Purges by a contained \(q+1\) strip through the tested position (using the margins for initially exposed sites); unmarked incumbents hold their banks. Thus outside K an earlier good central write persists or a later good write unconditionally restores that bank. In readout step 4 each repetition writes the result to all five fifths, not just the three middle ones. Its transaction interval is padded by \(2w^6\), while the single local response lasts less than \(w^6\); hence at most one repetition can meet the enlarged raw slab. In the positive case each of the other two repetitions writes a live liveness bit at every site of C outside K and the refill band. That band is K enlarged by \(2B+40q+40<.2A\) on each side, and has width at most \(.4A+d+2\) with rounding. Each such write travels an unbroken in-home path inside C, by the path induction above: its transaction misses the raw slab enlarged by \(w^6\), and at such times the geometric inputs give the late core C fully covered. At every such site an unmarked incumbent retains its bank by a Purge-backed \(q+1\) strip; recruitment cannot replace it. Thus two of its three liveness bits remain intact and live at verifier initialization, modulo the already consumed edge reserves. For the clean-center weak vote all three central writes give the specified weak vote, with the same bound on untouched sites.

\textbf{Liveness.} For a dead majority all initially active sites are in K, because both reliable dead repetitions can only deliver a dead or interrupted dead output, even to an outer site. Raw bank forgery/reset is charged in K. Up to the first purported activation outside it, an \(M=128d\) step sampling epoch has at most d original positive injection coordinates. Any forged directed transmission follows a line \(x- t\)=constant in the positive direction, \(x+t\)=constant in the negative direction; each constant gives at most one arrival at a given receiver in that epoch. K\(\times\)S supplies at most \(2d- 1\) values per direction. The resulting \(\le\)\(3d- 1\) count is below the required \(>4d\) threshold. A birth resets lanes, counters and \(\mathsf{Active}\). This rules out a spurious positive outside K at any later time.

For a live majority choose a clean no-birth verification interval. The genuine surviving home there decreases as an interval containing C. Any eventual survivor lies within \(20A+2M\) of an intact run of M initially certified positions, excluding the one \(.4A+d+2\) historic band and the \(12A\) margin. Along the interval joining this run to the target, all surviving positions and their connecting lanes remain present throughout. One complete epoch certifies the next \(M- 4d- 1\) positions: at the furthest new site the incoming distances to at least \(4d+1\) old positive sources are between 1 and M, so those distinct shifted impulses all arrive. Repeating epochs costs at most \((ceil((20A+2M)/(M- 4d- 1))+2)(M+1)\) ticks, bounded by \(100A+4(M+1)\) since \((M+1)/(M- 4d- 1)<2\). Choose the second plateau if J has expired more than \(w^6\) beforehand. Otherwise the first is clean by the \(>3w^6\) separation, and covers every survivor of that earlier plateau.

After this first certification and before a later J, at each complete Purge all eventual survivors are again certified. In each interval between Purges at most one new clean recruitment instruction adds two previously uncertified sites to the single clean component. A component beginning the next initialized Purge of length \(\le\)q dies. For every eventual survivor x of a longer component, take its subinterval on the row \(2(M+1)+2\) steps into the Purge. During these initial propagation steps there is no birth and only endpoint omissions, at most two losses per step. Length is at least \(q+1- 4(M+1)- 4>2M+1\). It therefore contains a solid M-site one-sided interval adjacent to x, held since the start. This time span contains a full sampling epoch, with arbitrary epoch phase. At least \(M- 2>4d\) donors were previously positive. The receiver is certified by the end of the epoch unless already certified. The same proof applies on the first clean Purge beginning just before a later J with its conclusion restricted to the prior clean times as necessary; just before J at most two sites from the last intervening birth need remain uncertified.

At the first complete Purge wholly following J add the at-most-\(25d+100\) raw/new coordinates proved in Lemma~\ref{lem:ordinary-response}. At most \(25d+102\) of the donors can then be uncertified; \(M- (25d+102)>4d\). The identical surviving-M-run argument certifies every eventual survivor of this Purge. In each later clean Purge there are at most \(2(d+10)\) newly attached positions over the whole home, because the first-clean-Purge component count was \(\le\)\(d+10\) (step~1 of Lemma~\ref{lem:ordinary-response}; its statement records the weaker \(3d+12\)) and clean updates can only merge components. This bound covers several small components merging into one, and again leaves more than \(4d\) donors in the one-sided M-run. Finally, if there is no complete clean Purge after J before the ballot, J began at most \(d+4q+6\) ticks before it. The final plateau reaches the ballot and has duration \(\ge\)\(100A+4\)(\(M+1\)), so it already included a complete clean Purge under the preceding certification invariant before this J; throughout the remaining time there are no births. Only K can be reset or forged. At the ballot every surviving site outside K with a live target thus acts on true \(\mathsf{Active}\); a missing payload bank at the fringe gives the legal live filler. Central sites off K retain their voted letters. A late raw damage with earlier, widely separated verification interval cannot quietly reuse an uncertified site after the last Purge.

\textbf{Support.} In a positive result split the old C on K; if a side has \(\le\)\(q+2\) sites its union with K has length \(\le\)\(q+2\)+d\textless{}\(2q\) and lies at the C boundary. Longer old sides have long proper held strips. The \(4q\) additional trim to \(C_*\) pays the two unprotected outer boundary sites; the same interval around K covers the support test and all potentially incorrect central letters. Absent a recent pulse there is no such actual loss. For the no-computation branch a surviving same-ordinary-home raw-only colony at the ballot has its origin and its residual spatial footprint to check separately: Lemma~\ref{lem:ordinary-response} gives the preballot width and dates, and the last no-birth period rules out an off-K newly recruited survivor. In the strong computed-D branch all matching lower occupancy outside K is absent by the negative verifier proof. At an empty K the entire cut is physically empty for this home. A remnant in K can add at most four layers per side before its first complete clean postcut Purge, to width \(\le\)\(d+8\)\textless{}q; different source labels cannot lend a strong link at their contact, so the zero/right-break calculation eliminates it. All ranks use the actual served data, not an assertion that a supplied symbol is immutable after its own \(g=0\).

The late-event verification uses the final no-birth interval reaching the ballot.
\paragraph{Positive ballot after a completed restoration}

The liveness proof restores \textbf{all actual survivors}, including a previously hit site in K, at the end of a complete clean post-event Purge once a clean earlier verifier plateau has certified the home. If the event instead precedes that bootstrap, the later clean plateau ending at the ballot certifies all its eventual survivors directly. Between two complete Purges the subsequent restoration argument again applies to every survivor; the final plateau has no births. Thus, when there is at least one clean restoring Purge wholly after S before the vote (or the clean final plateau itself supplies the bootstrap), every incumbent present in the preballot row takes the native positive decision. At a clean ballot a strict interior incumbent was proper by the previous-rank matching triple. Its positive successor is proper and served by the rank-\((j-1)\) proper-successor and service clauses of Proposition~\ref{prop:organization}\ref{org:successor}. Only already exposed improper endpoints can therefore be omitted. In particular such a clean positive ballot cannot split an existing canonical interval anew inside K. This statement is about liveness, without asserting correct restored payloads at a previously forged site.

If the positive ballot can make a new internal hole in K and the first clean plateau did not reach this vote, no complete clean restoring Purge had time to fit after the actual event. The last event thus ended fewer than \((2q+2)+(2q+1)+2=4q+5\) ticks before the vote (two endpoint-rounding ticks allowed); its start is within \(d+4q+6\). This whole interval is inside the final no-birth plateau. It permits the spatially common original K hit and a ballot loss in K to be handled as one short two-date disturbance without a new canonical recruitment between them. An intervening gap already opened by S can of course persist; the statement excludes only a further internal cut after old liveness had been restored. On a dead decision a forged \(\mathsf{Active}\) inside K can remain sticky until the ballot and the stronger conclusion is unavailable.
\paragraph{Mail between nonconsecutive transport instructions}

Successive logical hop instructions can be \(d_0\) fine ticks apart. Following a nondead received value backward across the intervening handles is valid only at a coordinate where the same canonical incumbent survived continuously and its designated lane was in hold mode on all those handles. A birth, including a reappearance in a formerly absent slot, initializes its mail and stored output banks to message-empty; the receiver's current native injection/overwrite clears any further stale value. A physically distinct predecessor on another geographical grid cannot supply this ordinary positive link. Within the same grid, an unhit continuation advances the clock by one modulo T and preserves its address residue even at normalization. Reading the grade at each preceding hop consequently forces the same programmed injection schedule all the way back. During actual input transfer and readout all unhit sites have ordinary normalized addresses; for a fixed clock, residue and location the owner is unique. A wrong clock cannot become synchronized merely by wraparound because every genuine continuation advances once. Therefore a prior message on a different lane, repetition, source phase or former extended label cannot furnish a nondead word after the required path has broken. When a good path is intact, induction through these fine holds followed by its next unique unit hop yields exactly the designated source coordinate from the path equation. This is also the reason an empty supplier cannot inject a live value through a different center's concurrent transaction.
\end{proof}

\begin{lemma}[Classification at an issued face]\label{lem:classification}
On a face for which its immediately preceding strong issuing cut is valid, exactly one of nominal live and proper dead holds. This implication also holds at a cap cut with the ordinary transition replaced by the constant D. It applies to arbitrary payloads entering the weak-cut calculation.
\end{lemma}
\begin{proof}


\textbf{Precise issuing-cut hypotheses.} On the actual aligned child grid the predecessor has either (i) no computation, (ii) computation with a dead result, or (iii) computation with a live result s in the native state alphabet. In (i) and (ii) all observed matching occupancy on the new face is confined to its exceptional column K, of width at most d, including arbitrary forged live children inside it. An empty exception gives empty occupancy. In (iii) \(C_*\) is matched off one interval \(I_{\rm dec}\) of width at most \(2q\), and every central-fifth coordinate outside K is present with the correct issued letter of the encoding of s. The exceptions in the two tests admit one common interval of width at most \(2q\). In particular this follows if the possible \(C_*\) loss is a terminal piece at K whose union with K has width at most \(2q\). These are conclusions required of the preceding strong cut, not premises about this tile's classification.

Put \(P_0=P/5\). For \(w\ge 100000\), \(d=4w+8<q/2\), \(2q<P_0\) and \(|C_*|>3q\). At a fixed packet offset u in \{0,...\(,P_0- 1\)\} the three tested positions are \(P_0+u\), \(2P_0+u\) and \(3P_0+u\). An interval of width \(\le\)\(2q\) can hide at most one of them. If two live-state decoding witnesses use different exceptional intervals, at least one of the three positions avoids both. Its carrier therefore equals the two alleged encoded letters at u. Apply this to each u. The packets agree at every offset. The live structural tags and fixed field positions of the serializer determine every scalar and every pair of payload halves, so the two candidate live states coincide.

In cases (i) and (ii) all observed matched children fit a single interval of \(\le\)d\textless{}q/2 sites. This is proper death by definition. Even deleting an arbitrary \(\le\)\(2q\) positions from the \(C_*\) test cannot hide its missing sites: at most \(d+2q<3q\) positions are accounted for. Hence nominal decoding yields D. In (iii) choose the common interval of the issuing cut as decoding exception. The issued live state s passes both central tests, so uniqueness gives \(x=s\). At least \(|C_*|- 2q>q>q/2\) matched positions exist, excluding proper death.

The pointwise weak vote has no independent formatting consequence. For example if it yields a common malformed live packet z in the three middle fifths off K, it cannot encode a live nominal state: for any candidate decoding exception \(I'_{\rm dec}\) and each offset, some replica avoids both K and \(I'_{\rm dec}\), forcing z to agree with the candidate's canonical packet at that offset. Malformation rules this out. This weak description alone gives no proper-death conclusion. When used as a supplier to two unaffected strong repetitions, each correct read receives z at each offset from at least two of the three homologues (one interval shorter than \(P_0\) covers the compromised coordinates). Both fresh parsers on that supplier return D. A weak empty supplier represented by three neutral D packets or absent messages also returns D. Two identical independent evaluations on the resulting three supplied states and fixed operands issue the same canonical packet of a legal live state, or both issue all message-empty. Thus a malformed weak string is never promoted merely because its letters voted live. The geometric/transaction premises that actually provide those two reads, maintain their banks and impose the strong cut are separate obligations.

\textbf{Shifted faces.} For an internally queried protected face run the preceding issuing calculation in its buffered history, with earlier weak descriptions at the three supplier predecessors of that calculation. Apply this lemma after the strong cut at this rank and before invoking classification for proper-death, disjointness, properness, or service of the current face. In constructing that issuing cut the no-split and protected-geometry premises refer to the rank below. Its reading of a weak malformed supplier uses the elementary replica and parser argument above, not classification of that same-rank weak face.
\end{proof}

\subsection{Grow sources, density and separation}
The terminal geometry adapts G\'acs \cite[Section~15]{gacs86}.
\begin{lemma}[Attachment and local repair of an issued Grow source]\label{lem:attachment}
Away from the actual or delayed wound in this epoch, an issued source is an interval attached to its original ordinary core. Its core \(C_*\) is immune outside its localized wound, and is wholly present after repair. A genuine distal piece cut off from that interval is included in its ancestry. During repair each of its internal holes lies in K enlarged by \(2B+40q+40\). If normalization occurs first the assertion ends there with this last bound. An unrelated short satellite may be arbitrarily far from the authentic tip, but cannot supply a long continued arm by itself. Two sources can change the exact division of an adjacent repaired bridge.
\end{lemma}
\begin{proof}
Use child units. A source is the exact integer base \(hP=x-\pi\), with its genuine clock and address residue. Until the final normalization, an authentic child inherits this same base. A temporary raw-only satellite is counted separately until its first clean Purge. The hypotheses are the preceding-rank nominal and served clauses on all shifted neighborhoods, the established previous ordinary single-interval dichotomy, the prior issuing cut with live decision at all its surviving proper incumbents outside K, and the buffered two-flank clearance implication in Lemma~\ref{lem:flanks} (including the previous-row and four-neighbor tests). Lemma~\ref{lem:local-wound} supplies the actual birth and assignment counts to the first clean Purge for an exact source (the same count as step~1 of Lemma~\ref{lem:ordinary-response} for an ordinary home). This statement is about a previously \emph{issued} positive, not a merely formatted array put on an arbitrary initial row.



\textbf{Source at the vote and real backing.} Fix the true old ordinary home of a positive cut. When that cut and its preceding ordinary work are clean the genuine pre-vote interval is nonempty by positivity; no birth occurs in the ballot, scratch is reset and the live decisions retain its proper sites. Only an initially improper extremity may be lost on either boundary of a genuine interval. The new source is its resulting interval. With a late J retain the last pre-J retired ordinary strobe I. If I were empty, the only independent raw support before the first clean Purge would have span \(\le\)\(d+8\); this could not account for a macroscopic positive cut. More precisely choose a positive cut marker in \([P/5,P/3]\) and one in \([2P/3,4P/5]\), avoiding the formatting exception and K enlarged by \(10B\). Each band has width \(\ge\)\(2P\)/15-2\textgreater{}\(20B\)+\(d+2q\)+10. Each chosen marker is within \(8B\) of I: otherwise a centered vacant interval of \(4B+4\) positions within its host home avoids I,K and both true-edge residues, and the empty-band induction of Lemma~\ref{lem:ordinary-start} (also rederived in Lemma~\ref{lem:remote}) uses only provisional birth/retirement in this ordinary part. Written independently here, each complete attempt enters \(\le\)B sites from each direction and every such still separated entrant has a marked, unexempt, zero-grade face directed to the central empty strip; B removals clear its entrances. Its center cannot be occupied at the ballot. Hence \(|I|\ge P/3-16B-4\). This argument only uses the native empty strip and does not assume wide ancestry. When J is in a subsequent Grow round with a clean preceding ballot, the retained ordinary proper pieces of the preceding productive core give the same macroscopic backing.

\textbf{Run protection.} Apply Lemma~\ref{lem:local-wound} to the actual macroscopic interval just obtained. Its length exceeds \(2q+d+44\) (it is at least \(P/3-16B-4\) for a late J, and contains the retained ordinary core pieces in a later Grow round), its strict interiors are proper by lower triples, and its evolution is clean from its latest scratch reset until S: a \(\mathsf{Temp}\) reset, or the fresh reset of the positive vote, whose live decisions keep the proper incumbents outside K by the prior issuing cut. By the run-protection clause of that lemma each original \(q+1\) ordinary run of C away from K, trimmed by one site at each end, persists throughout Grow.

\textbf{Clean source.} Inside an exact source every adjacent pair is a strong link. At a structural step its strict interior is proper and computes a live result. On a birth tick it also holds nominally, and the next nominal triple of an inner neighbor, old endpoint, and a new outside child makes the old endpoint proper and served. For a singleton with two children use the same triple; for one child, dropping the parent still leaves one interval. Without births only an end of the exact source interval can be lost. The independent clean-evolution paragraph in Lemma~\ref{lem:components} expands these arguments, including length two; it is a forward pointer and uses no attachment conclusion. Thus a clean genuine component cannot split. Zero-grade true-end exemptions apply to \(\mathsf{Temp}\) only at the outward true face. A distinct source, including one of the same clock with a positive wrap, has a nonstrong bond to this source at a contact: equality of successive extended addresses would assert equality of \(x-\pi\). Therefore a short distal source piece whose only support is the other source's grade-one junction still has a right break and a left break for Purge. When its length is \(\le\)q at an initialized reset, its entire strong component clears in that Purge. Pinning by a true-end exemption does not affect the propagation or right-break killing instructions either.

\textbf{Bookkeeping and refill.} Lemma~\ref{lem:local-wound} supplies the \(B+30q+40\) cavity bound, the K enlargement, and reunion or distal extinction within \(2d_0+10H_0<w^6\), stopping at normalization when necessary.

The original C had an extra \(4q\) reserve beyond \(C_*\). At the cut the issuing lemma retains the long ordinary pieces of C outside K, losing at most the one short edge piece of length \(\le\)\(q+2\) there. Every \(C_*\) position outside the wound therefore has one of the fixed long ordinary runs of Lemma~\ref{lem:local-wound}, and a missing \(C_*\) position at an internal wound has long ordinary pieces on both sides in C, so it is filled by the genuine two-sided case. The only other possible loss of a short edge piece capped at \(q+2\) lies within the outer \(4q\) reserve. Thus \(C_*\) persists except during its localized repair. A raw-only label outside the authenticated interval is precisely the short cleared satellite case.

\textbf{Uses.} This provides the earlier attached interval, surviving core, localization of real gaps and short-pinned-flank alternatives for density, wide comparison, root identity and true-end relief. It does not require the contested geometric seam between two positive homes to retain its old division into exact source labels.
\paragraph{Wound on both sides of a positive vote}

There is a case in addition to an actually straddling S and two fully separated repairs. Suppose S ends before the positive vote but a completed geometric refill has not yet occurred at that vote. In particular a terminal ballot could appear to reopen K during the previous wound's response time. The following calculation uses the preceding positive-verifier proof in Lemma~\ref{lem:ballot}; this is available before the current Grow source lemma.

If the ballot's immediately preceding complete Purge started after the last affected transition and the new S is this close to the vote, the first of the two separated verifier plateaus had finished before S (their separation exceeds \(3w^6\) while this case lasts less than \(2d_0+10H_0)\). Every surviving position at the final Purge, including positions in K, is certified. Indeed the first complete post-S Purge before the ballot certifies every eventual survivor: in its initial \(2(M+1)+2\) propagation rows each such survivor has an intact M-site one-sided interval, with at most \(25d+102\) positions not certified by the earlier invariant, and \(M- 25d- 102>4d\). A short strong component dies. Each intervening complete Purge recertifies newly recruited neighbors by the same inequality, and the instruction immediately before the final ballot is the complete Purge in question. All these certifications are sticky in the absence of another S. The actual positive ballot thus makes no new inactive hole, even if the payload at K is filler or wrong. Lower masking can remove only exposed ends of an already existing ordinary exact component: strict interiors are proper and have a positive certified native ballot, propagating by the rank-\((j-1)\) proper-successor and service clauses of Proposition~\ref{prop:organization}\ref{org:successor}. The ordinary true addresses become the new epoch's normalized addresses at the same coordinates. A panel of length at least q-1 thus loses at most one position on either side across this ballot. Its interior wall and exclusion of incoming paths persist. A short raw-created component already joined to a genuine side does not detach from that side on the clean certified vote. In this case use a single ordinary/Grow response beginning at the original S and continue its birth-handle accounting to the first chosen growth pass; the ballot itself is a no-birth, no-interior-split instruction. The regular handle cadence is continuous at that instant and the fresh reset precludes an old temporary mark crossing the vote. The total time from that S to the end of the selected pass and extra Purge remains \(\le\)\(2d_0+10H_0\). Thus conditions (i)--(iv) of Lemma~\ref{lem:local-wound} hold at this ballot: no ballot decision is non-native, a panel loses at most one site per side, the reset and cadence are as just stated, and (iv) is the standard count because the first complete clean Purge after S precedes the ballot.

If instead the last preballot Purge intersects S or S includes the ballot, then the first raw transition is within \(d+2q+3\) ticks of the vote. This interval is strictly inside the final no-birth plateau since its length exceeds \(100A+4(M+1)\). Use the clean genuine ordinary interval just before S. The raw updates and the exceptional decisions at the following ballot are all at positions in K. The latter assertion uses the first clean verifier plateau already completed before S, no subsequent births near this final Purge, and sticky certification off K: the whole immediately preceding plateau contains an earlier full clean Purge that certified its entire surviving group. From the start of S to the first complete clean Purge of the new Grow takes at most \(2d+6q+10\) ticks, with no native birth before the vote and at most two after it. The authentic intervals outside K are therefore shortened only at exposed ends. A clean proper original ordinary run of \(q+1\) positions off K survives both cuts by the fixed-run proof, and its reset sites remain ordinary incumbents. At most d coordinates inside K have been reassigned at either exceptional operation, regardless of how often a coordinate is replaced, and counting four honest growth opportunities still bounds the assignments by \(25d+100\). A shifted foreign family in the earlier cavity either comes from actual S or an old shifted ballot descendant of an S ancestor; no foreign genuine external lineage crossed a surviving wall. A potentially spurious \emph{canonical} decision at the vote is confined to the same K and the ensuing proper Purge clears its separate short component. Choose clearance after the later of the actual S end and the vote, plus \(4q+14\) rows. This incurs \(\le\)\(2d+9q\)+40 ticks after the original S, within the \(t_c\) of Lemma~\ref{lem:local-wound}. There is no additional full B-probe pass before the selected one. Hence conditions (i)--(iv) of Lemma~\ref{lem:local-wound} hold at this ballot: (i) and (iv) were just shown, the authentic intervals outside K are shortened only at exposed ends, at most one site per side as in the previous case, and the reset and cadence are as there. Thus the same bound \(B+30q+40\) for an enduring internal cavity and \(2d_0+10H_0\) for total elapsed time apply. No passage combines two complete B-growth passes into a budget of four birth instructions.

For completeness, the first raw-only enumeration can leave additional \textbf{small canonical inner fragments}: the two rays outside K enlarged by four contain at most one true same-source interval each. A separated component not reaching either ray has at most \(d+8<q-2\) members. Before its next whole initialized Purge it can add at most two new members. It therefore dies in that Purge or joins an adjacent compatible flank, including two fragments coalescing inside the narrow middle. The union of all fragments that have not contacted a flank remains in the same original column with one further opportunity on each side, of width at most \(d+10<q\), so coalescing among themselves does not rescue them. Any contact with an enduring genuine side only enlarges that side. Allow this one additional Purge in the \(9q+40\) buffer of \(t_c\) in Lemma~\ref{lem:local-wound}: an initialization occurs every \(H_0\) ticks, so two clean complete purges have ended by last \(S+6q+4\) (overlapping the elapsed time with the displaced-grid guards). Genuine two-sided pressure may come from different source labels at a repaired seam: the source lemma requires its two internal walls to be the same exact label. When it is applied to a pair across a true wrap it gives only persistence of the two individual backings.

The spatial wall test for the final attempted birth requires only past and present rows. A shifted nominal tile overlapping canonical row t has a preceding served tile whose time interval intersects canonical row t-1 or t for positive duration. Choose a time in that intersection and trace its ancestors backward one unit at a time. A shifted actual blocker is tested directly at its overlap with t. No row later than t is involved; the lower proper-successor test begins one canonical row earlier. At the last attempted birth before a join or extinction the wall need endure only through the old row of that update.
\end{proof}

\begin{lemma}[Deficit from a component count]\label{lem:components}
Let \(H=[0,P-1]\) be the home of a positive issue. In its comparison period from \(r+\lceil v/20 \rceil d_0\) through early \(\operatorname{Shrink}(0)\), suppose (i) every authentic source before the single exceptional rectangle was one interval, possibly empty, (ii) the only additional exceptional cut is one affected ballot at which only positions in K can make erroneous decisions among the ordinary pre-vote survivors, and (iii) outside those exceptions the preceding-rank regularity, nominal triple and service clauses hold. A previously earned cut may initialize a source with at most one interval of erroneous positions in K, or two genuine flanks around K left by a straddling pulse. All intervals here refer to a single exact extended label until normalization. If every interval of missing positions in H has length less than \(L=60B+3d+40\), the total deficit in H is at most \((5(4d+12)+1)L\). In particular it is at most 2000000 A for \(w\ge 1000\). The same bound holds before, within and after the event, and on entering ordinary work. An exact count with \(d+2\) components per source is unnecessary.
\end{lemma}
\begin{proof}
Time and space are integer child-grid indices. The initial source interval comes from the already established issuing cut and the source geometry of Lemma~\ref{lem:attachment}; the clean-evolution paragraph below is independent of that input.



\textbf{Proof of clean evolution.} Follow a connected interval of one source whose neighbors at both ends do not belong to the interval. During a birth step each old site has a live nominal successor with the correct clock and address and retains its source label. If there are three or more old sites the strict interior is proper by the nominal triple, so each interior site supplies its successor. If the newly born child immediately beyond an end is served, the former end has a live nominal successor between the old inner neighbor's live nominal successor and the live newborn. It is therefore proper and served at the new date. In length two this same argument applies at each side that gives birth. Any subset of a two-site old interval is an interval; two attached children cannot detach. For length one with two born children, the middle has a nominal matching triple and is supplied. With just one child the union of that child and the surviving old singleton is an interval even if the singleton drops. On a no-birth step at fixed source labels all strict interior cells were proper. A strict interior faces grade two on both sides; neither Purge killing, either directional temporary removal, Shrink nor native withdrawal acts there. Thus only interval endpoints can be removed, and every surviving interval is connected. These facts also apply to an interval artificially truncated by a moving boundary if a one-cell margin is lost rather than gained. Two old intervals may merge, but this never increases the count. Matching incumbents belonging to other labels never switch label except at normalization and the ballot.

\textbf{Effect of a pulse.} Enclose the raw spatial interval K (length \(\le\)d) at the start in \([a,b]\). At its ith subsequent fine date enclose \([a-i,b+i]\), for \(0\le i\le d+1\). The last interval has length \(\le\)\(3d+2\). No locus in the remaining left half-line at time \(i+1\) has an old neighbor in the excised region at i; the same holds on the right. On either exterior half-line the surviving descendants of each of its old intervals thus form at most one interval by the preceding argument and restriction to a half-line. This remains valid whichever admissible end omissions occur, without comparing to any counterfactual history. Inside the cut-out region any subset has at most as many connected components as it has sites. Starting with a single old interval, the total is at most \(3d+4\) per source at every intermediate date. Starting with two intervals costs two additional exterior components. A delayed affected ballot can delete any subset of K or forge a matching member there. For arbitrary membership replaced on an interval of length d the number of components increases by at most \(d+1\): restriction outside it can split at most one existing component into two, and the new inside has at most d components. If the ballot is inside the excision it has no exterior cost. An honest positive ballot preserves the occupied interval apart from possible improper ends (all incumbent live decisions off K are valid); an honest negative ballot leaves it empty. Counting both exceptional operations, including a vote outside the pulse, gives at most \(4d+7\) per source. We allow \(4d+12\). If the changed cut is first, cutting K out leaves at most two exterior pieces of the preceding genuine source and interior K at that cut. The expanding cone argument then applies directly since it encloses precisely the same column K. A distinct actually empty source can start only in K. After this pulse clean evolution cannot add a component. At normalization the matching material of H is the union of its contributing sources intersected with H. Taking a union does not increase the sum of component counts.

\textbf{Counting homes and gaps.} Translate so a home starts at zero. The base of a possible exact source is hP for an integer h. An occupant x in H has legal address x-hP in \([-1.1P,2.1P)\). If \(h\le -3\) this is \(\ge\)\(3P\), and if \(h\ge 3\) this is \(\le\)-\(2P\)-1. Hence h belongs to \{-2,-1,0,1,2\}. Intersecting an interval with H leaves an interval. There are at most \(5(4d+12)\) occupied components, hence at most \(5(4d+12)+1\) missing intervals counting both boundary gaps. Multiplying by L proves the estimate. If normalization occurs before the pulse, apply the expanding cone to the two remainders of the whole matching home. Splitting old components at the two walls adds at most two, then the cone adds at most \(3d+2\); before a first pulse at most five authentic intervals have entered and at most five remain. A delayed cut in the same column satisfies the same replacement bound. All these counts are smaller than \(5(4d+12)\). For \(w\ge 1000\) one has \(d\le 4.008w\) and \(B=128w^2\), so the first factor is at most \(81w\), \(L\le 7681w^2\) and their product \(\le\)\(622161A\).

\textbf{Uses and dependencies.} Combined with the early productive-home sterile-strip implication (no interval of L missing positions), this estimates a single deficit on a common date for all markers. During later ordinary dates the computing core bound supplies density. After the ballot use the held core of each issued source. These conclusions enter nominal separation and the comparison against terminal contenders during end relief. This lemma does not derive its starting premise of genuine source intervals, the true classification of decisions at the vote, or the productive-home sterile-strip implication. A source-preserving component bound remains valid even if the source has not yet healed its one wound.
\end{proof}

\begin{lemma}[Thin ancestry and productive-home density]\label{lem:thin-density}
Fix one matching tribe in its final v-round withdrawal and ensuing ordinary period before its next ballot, under the protected lower structural clauses outside K\(\times\)S. Every genuine terminal observation at n has a member of the same tribe within \(100B\) of n at each earlier date of this phase interval; an isolated immediate raw-only fragment is instead allowed in the short column transient. In a home H whose ordinary work issues positive at T, no interval of \(60B+3d+40\) sites in H is wholly vacant at any fine date of the final withdrawal or of ordinary work before the early seed cutoff (at all such dates at least \(3P+10\) \(\operatorname{Shrink}(0)\) inspections remain). With Lemma~\ref{lem:components}, the end budgets and late core of Lemma~\ref{lem:ordinary-response} and the held core of Lemma~\ref{lem:attachment}, at every date of \([r+\lceil v/20\rceil d_0,T+s_0)\) the deficit in H is at most \(2{,}000{,}000A\); no dates in that interval are excepted from this deficit bound.
\end{lemma}
\begin{proof}
Use one matching child tribe in its final v-round withdrawal and following ordinary period before its next ballot. Work on a buffered test for which the previous-rank exact-or-dead and structural conclusions hold outside K\(\times\)S. Write J empty if no top exception is needed. Only the B provisional attempts per Main cycle can recruit in this interval. Normalization preserves residues and changes no coordinate. An unflagged newly matching state must arise from a prior matching incumbent or donor. Distinct exact extended sources of this tribe have different \(x- \pi \), so cannot lend strong backing across a source boundary.

\textbf{Thin trace.} For a genuine terminal observation n at t, each earlier date u of this phase interval has a member of this same tribe within \(100B\) of n. An observation which has only isolated recent raw ancestry in a span of \(d+8\) at K, before its first clean complete Purge, is allowed as an alternative.

Suppose the interval of radius \(100B\) is empty at u. Work separately in the integer bands of distances from \(3B\) through \(90B\) on either side of n. Remove the coordinates in K and those congruent to either true endpoint residue. Each candidate band is shorter than \(P- 2\) and contains at most one adjacent endpoint pair or single endpoint. The deletion leaves at most three intervals of total length at least \(87B- d- 2\), so one has length at least \(8B+8\). Select such a band on each side. Starting from u the remaining first partial probe can move at most B places towards the central portion of a chosen band from either edge, and the two directions together cannot reach its central \(4B+4\) places. No true Grow attack is scheduled here. Only provisional recruitment enters these particular vacant strips, and its marked entrants cannot join across the vacant center or reach a true endpoint. The empty-band clause of Lemma~\ref{lem:ordinary-start} gives a fixed vacant central site at every subsequent fine date and restores a \(4B+4\) strip to vacancy at every strobe. The intermediate region between the two permanent barriers was empty at u. No actual matching descendant can arrive from the exterior on any tick: it moves by at most one, its boundary crossing requires the fixed vacant site, and a true nonbirth cannot spontaneously create a matching state. Thus any later occupant near n needs a raw assignment in this interior. From the first exceptional step through the first initialized clean Purge after its end there are at most four genuine recruitment instructions \((d+4q+6<3(2q+2)\) for sufficiently large w, allowing rounding at both end dates). Its spatial support lies in K enlarged by four on either side throughout. At complete purification every same-label strong component has length \(\le\)\(d+8\)\textless{}q; by q propagation and q killing ticks it is empty. If normalization intervenes before the full Purge, count the whole matching union, still in the same \(d+8\) span. No scheduled Purge straddles the normalization. The first complete clean Purge still finishes within \(d+4q+6\) fine ticks of the first exceptional step. This gives precisely the stated short alternative or a contradiction. The proof fixes u arbitrarily, so it does not require a single material particle surviving at every earlier date. If the observation immediately follows a no-birth cut, use its actual preballot incumbent or a recent raw substitution before applying this statement.

\textbf{Empty interval incompatible with a positive issue.} In the same home H whose ordinary work issues positive at T, no interval of at least \(L=60B+3d+40\) sites in H is wholly vacant at a fine date in the final withdrawal or before the early ordinary seed cutoff, provided at least \(3P+10\) \(\operatorname{Shrink}(0)\) inspections remain. Removing K and the two true boundary sites leaves at most two nonempty interior pieces in H (the two endpoint deletions in H do not cut an interior piece). Their total length is \(\ge\)L\(-\)\(d- 2\), and one contains an empty band of length \(8B+8\). The first partial phase and sterile-band proof carry its central vacancy through normalization and all ensuing stage-0 inspections. On either side of that strip the nearest occupied strobe frontier recedes by at least one per eligible clean \(\operatorname{Shrink}(0)\), by the erasure proof in Lemma~\ref{lem:ordinary-start}. After discarding at most three inspection spacings intersecting the one pulse and its boundary strobes, one of the two remaining stretches has more than \(P+1\) inspections. If it follows the pulse all matching home material is extinguished. If it precedes the pulse only the later raw-only family can repopulate H, and it disappears in its first complete clean Purge unless the terminal no-birth interval and the pulse overlap. In that last case only cells in K can remain at the ballot. Such a group cannot meet the positive nominal support test. This contradicts the positive issue. The six-home seed search uses only a polynomial number of fine ticks past normalization, leaving more than \(3P+10\) calls. Once the ordinary interval and threshold ledger apply at clean strobes, occupied C and its possible spatial band of width less than \(.4A+d+2\) give deficit \(\le\)\(24A\)+.\(4A\)+\(d+4\). The band argument applies at every date with more than \(3P+10\) \(\operatorname{Shrink}(0)\) calls left. This includes the entire initial \(50A\) \(d_1\) relaxation and the polynomial common-seed search, so it overlaps the strobe ledger. In the later stage-zero and stage-one work the respective end budgets are \(3A+3\Delta \) and \(6A+3\Delta \). Provisional attempts between strobes do not decrease old occupancy, and the next true Shrink and optional old-end loss cost at most the stated edge allowances. The single internal wound lies within K enlarged by \(.2A\) on either side. Stage two and computation use the \(12A\) core bound. After T the positive source's held \(C_*\) across Grow gives deficit at most \(24A+8q+2\) plus its single localized wound of width \(\le\)\(4B+80q\)+\(d+82\). Thus throughout \([r+\lceil v/20 \rceil d_0,T+s_0)\) a constant \(2,000,000A\) bounds the common deficit, using the explicit component count of Lemma~\ref{lem:components} for the earlier dates. All those counting hypotheses concern the already proved previous ordinary cut and the genuine-source geometry, not the conclusion of disturbance inheritance.

\textbf{Uses.} The thin conclusion supplies an actual earlier witness for the remote-marker estimate and for end relief. The positive-issue deficit is counted once per date against an entire competing interval.
\end{proof}

\begin{lemma}[Remote sources during terminal withdrawal]\label{lem:remote}
Fix one child-grid tribe in an epoch whose final v-round withdrawal starts at r. Let n be the coordinate of a genuinely supported matching member observed at withdrawal age at least \(7v/10\), or later through the final ordinary preballot date. Assume authentic-source attachment to enduring cores reaching within \(12A+4q+2\) of both true edges, K-only extinction at the first complete clean Purge, repair length at most \(w^6+2d_0\), and the thin trace of Lemma~\ref{lem:thin-density}, under the protected lower structural clauses. If the spatial home H containing n has dead nominal input, then a live home is within \(300B\) of n, or matching material covers \([n-v/5,n+v/5]\) at every fine date in \([r+\lceil v/20\rceil d_0,r+\lfloor3v/10\rfloor d_0]\), apart from at most \((220A+2w+20)d_0\) fine dates. An immediate cut remnant is first traced to its actual preceding preballot observation.
\end{lemma}
\begin{proof}
All spatial coordinates and fine times refer to one child grid. A strobe is the row immediately after its true Grow withdrawal. Write \(k=12A+4q+2\), \(W=4B+4\). Enlarge w to at least \(10^9\). Hypotheses supplied by Lemma~\ref{lem:attachment} and the earlier structural lemmas: throughout the final withdrawal, the authentic sources have enduring ordinary cores reaching within k of their two true edges at each clean strobe, each source is connected to its core outside the single repair interval, and an independent K-only source perishes at its first complete clean Purge. The repair interval has length at most \(w^6+2d_0\). Lemma~\ref{lem:thin-density}, including its raw-only alternative, gives an actual member within \(100B\) of a genuine late marker n at each earlier date in this withdrawal. We use that lemma on a date away from repair, not an inference that the same individual lived throughout.



\textbf{Sterile band.} Start from W consecutive empty sites at a strobe in a band disjoint from K and positions congruent to 0 or P-1. In the final withdrawal no Main instruction recruits. At each subsequent cycle the only matching new material is a marked Heal birth, which requires a previously occupied adjacent donor. In a growth pass at most B sites can be entered from either side. The central \(W-2B\) sites stay empty. Each entering one-sided component is marked at every interior point of the band. Its inward extremity has grade zero facing the empty center and is not exempt there. At each of the B temporary-removal instructions the foremost remaining site is removed, unless an intervening Purge or a prior omission already removed it. Its distance from the corresponding band boundary decreases by at least one until the entrance is empty. Raw sites and true-end exemptions never occur in the band. This induction applies to the whole matching tribe, so a different same-phase label cannot make the two trials meet. The central subband stays empty at all fine dates. The full band returns to empty at every next strobe. Across normalization, if needed for the thin lemma, use its separate strobe boundary with a new \(\mathsf{Cn}\) initialization.

\textbf{Geometry at a clean strobe.} Take the nearest authentic home to the left of H and to its right, omitting a side which has no contributing home. Their order follows from the enduring core and attachment hypotheses: a further connected source would have to occupy the core of the nearer one. At a joined row their seam is a positive bond, since the left address is at least \(P-1-k>P/2\) and the right address at most \(k<P/2\). Each side has its own long strong backing, and all seam sites are in a nominal triple. Neither a true Grow, a temporary removal, a Purge nor a clean service omission can open this seam. At a first contact the two provisional chains already have positive facing bonds, so the B temporary removals do not undo their join. The next \(\mathsf{Temp}\) reset makes the retained sites unmarked. The whole interval between the two deep cores is thus occupied until the next event. It contains the claimed band at every such early joined date, because each core extends into its own home a distance greater than \(P/2\) and H is between the homes.

Otherwise write p for the left source's rightmost actual site, z for the other's leftmost site, e,f for their facing true edges. Define \(U=\max (e,p)\) and \(V=\min (f,z)\). The actual open gap has length z-p-1. Its excess over V-U-1 is \((U-p)+(z-V)\), an integer in \([0,2k]\), since the cores are occupied. Here \(2k=24A+8q+4\) \(\le\)\(25A\) for \(w\ge 10\). There is at least one whole home between e and f, hence \(f-e-1\ge P\). On a retained clean unjoined cycle an exposed terminal tip can have a positive outward bond only to material of the same tribe, with the clock and address-residue alignment required by the bond rule. A neighbor with consecutive exact address belongs to the same root, contrary to extremality; a grade-one neighbor from the opposite authentic root is the joined case. The only other possible supporting label is a J-made remnant of width at most q. Its repair cycles are discarded, and at its next initialized Purge it has no q strong predecessors within its own exact label and dies; a grade-one junction supplies no strong Purge backing. Thus no foreign tribe or short supporting label can pin the tip on a retained unjoined cycle. Each extended exposed site then has a free outward face and withdraws or has already disappeared. No new unjoined temporary site can remain beyond the true edge, since each faces an empty site, carries a mark and has the full B withdrawals. If a true endpoint is reached from within its own home the projection already pays for that advance. In particular
\(U_{next}\le \max (e,U-1)\), and \(V_{next}\ge \min (f,V+1)\).
If the projected gap is below \(P/2\), both variables cannot simultaneously equal e,f, so the projected gap grows by at least one. On the first open strobe \(V-U-1\ge -2k\) because \(z-p-1\ge 0\). Consequently after \(110A\) clean unjoined rounds, either \(V-U-1\ge P/2\) or the actual gap is at least \(110A-2k>40B+2d+40\). An interval of length \(40B+2d+40\) is shorter than P-2; deleting K and the pair of true-end residues removes at most \(d+2\) places and leaves at most three intervals, unless the residues lie at opposite endpoints of the tested interval. In the latter case deleting endpoint residues creates no additional intervals. Some remaining interval has at least \((40B+d+38)/3>W\) places. It contains a sterile band. Once such a band is certified the two sides cannot subsequently meet. Discard the first \(110A+1\) comparison cycles and the first \(110A+1\) after the single disruption, plus all cycles meeting the repair interval. This works for every good early strobe. With only one side, no separating band is needed for its clean-tip recurrence.

\textbf{Permanent upper bound across the disruption.} This argument uses a second band and does not bound the number of provisional passes by a constant during the broad \(w^6\) repair. Fix a retained unjoined early strobe. If the disruption lies between it and the late comparison, take U in the gain estimate at the last clean strobe before the disruption: the band certified at the early strobe persists to that strobe. If no disruption intervenes, use the clean recurrence alone. Consider the left root at the selected strobe. The two-root case at such a date has already waited out the \(110A\) grace period and has a certified band; joined dates already supply the wide coverage. If a previously certified separating band exists, call its leftmost site h. It is empty at this strobe, and the genuine left root lies to its left. If \(h\le U+30B\), the farthest connected site of that root after restoration is strictly less than \(h+W\), since it cannot cross the center of the sterile band. Thus \(U_{after}\le U+35B\). If \(h>U+30B\), or no opposite root exists, the first \(30B\) places to the right of U contain no matching member at this strobe. In the two-root case the right source lies to the right of the original band. In the one-root case there is no other connected root on this side of its own core, and the observation window need only extend \(30B\) beyond the relevant tip. Remove K and the pair of true-end residues from these \(30B\) places. The remaining pieces number at most three and total at least \(30B-d-2>3W\). Select a new sterile band there. Its central gap stays empty throughout the disturbance, and the restored genuine source is connected to its original core, which lies behind that band. Its rightmost postrepair site is at most \(U+30B\). A temporary same-label source forged on the far side must either die, or merge with a similarly labeled authentic source. Once it is restored and connected to this root it cannot be on the far side. In both cases \(U_{after}\le U+35B\). Exchanging left and right proves the opposite bound. The test uses only old certified vacant space and postrepair attachment; it needs no claim that a particular exposed tip remained the tip during the pulse.

\textbf{Iteration and dates.} Choose the late comparison strobe between rounds \(\lceil .64v \rceil \) and \(\lfloor .66v \rfloor \). Fewer than \(w+6\) strobes meet the repair interval, by \(d_0=(2A+4)(2q+2)\ge 4w^5\), even including two rounded endpoints and the liberal extra \(2d_0\). Such a clean strobe exists since \(.02v-2>w+6\). Thin tracing supplies a live site within \(100B\) of n there. Its source is authentic by clearance of the short transient. A clean joined early strobe was already dealt with, irrespective of a later break. At an unjoined retained early strobe choose its permanent separating band as above. The late supported root cannot have passed through this band, and it is connected back to its original core. Thus the root on the late side really is the same rooted label on that side at the early date. This also treats an earlier joined configuration opened by J: its early joined dates already have full coverage, its first \(110A\) postrepair rounds are discarded, and every subsequent separated date bears its own newly certified sterile band.

For the late left root put \(N=w+6\), \(D=35B\), and M equal to the difference of the two strobe indices. At clean rounds the projected recurrence holds; during the single interruption bounded by N skipped rounds use the permanent gain D once. Composing a clean block of a steps maps u to at most max(e,u-a). Composing two such blocks with the gain between them gives \(U_{late}\le \max (e+D,U_{early}-M+N+D)\), also valid if the gain is absent. If all live homes are more than \(300B\) from n, left-side supply within \(100B\) forces \(U_{late}\ge n-100B>e+D\). Therefore \(U_{early}\ge n+(.64-.30)v-(N+135B+2)>n+.33v\). This is a truly extended tip, so its connected interval extending back to the old core covers the required band, even allowing \(B+2\) places of loss when passing from a strobe to a fine date. More precisely, on a full clean round a macroscopic connected proper side cannot suffer a native interior loss. Its post-withdrawal tip and its preceding tip differ inward by at most two (one improper endpoint and one actual withdrawal); on partial rounds no interior point common to those strobe intervals can disappear and reappear behind an unfilled break. Alternatively repeat the calculation at both endpoint strobes and use persistence of the large strict interior. The right-root derivation follows upon replacing every position x by -x and defining its projected outward extent as -V; the estimate concerns the true Grow and temporary directions on exposed sites and relies on separate q-backing for Purge on each actual large source. The count of discarded cycles is at most \(220A+2w+20\), including neighboring strobe cycles for all sampled fine dates. The sets deleted depend on the common disruption and the epoch start, so they apply simultaneously to every n.

\textbf{Uses.} Nominal inheritance at a deep issued marker, and terminal exclusion for end relief, retain the \(300B\) near-source option and the \(.2v\)-wide band. Both tolerate the explicit \(35B\) permanent gain, which is already paid by the \(300B\) cutoff. The linear recurrence and the two sterile bands require the preceding authentic-source attachment and core lemmas. Neither hypothesis follows from formatting an arbitrary input cut.
\end{proof}

\begin{lemma}[Backward inheritance across an incoming pad]\label{lem:inheritance}
Suppose two nominally live tiles U',V' of different geographical origins have overlapping temporal rectangles and one spatial home meets the other's \(64A\) pad. Then their same-home predecessors U,V one epoch earlier were both nominal live, and the translated pair retains the same overlap/pad relation.
\end{lemma}
\begin{proof}
This proof uses the earlier positive issuing cut of Lemma~\ref{lem:ballot}, productive-home density of Lemma~\ref{lem:thin-density}, protected lower exclusion and Lemma~\ref{lem:remote}. All examined work rectangles and the preceding supplier computations are inside the common sparse test. Positions and dates below are child coordinates; fractional offsets of different child grids are retained when intersecting rectangles. At most two positions at the ends of any interval need be discarded to select whole child tiles of one grid. Select all comparison dates in the interiors of fine tiles. Choose \(w\ge 10^{10}\), large enough for the explicit component-deficit constant 2000000 and the at most \((220A+2w+20)d_0\) date deletions for each of the fixed comparisons.

\textbf{Identity.} There is no assertion that an individual physical site carries one exact extended address from the terminal ordinary date through Grow normalization. At a genuine late marker thin tracing supplies a supported matching member within \(100B\) at a clean strobe in \([.64v,.66v]\) of the same final withdrawal. Before normalization its extended integer identifies one exact source. At each good early strobe, there are the following exhaustive cases.

\begin{enumerate}
\item There are two roots and their intervals join. Their union at that very date is the required band; no terminal-label conclusion is needed then, even if J later opens a break at their seam.
\item There are two separated roots. The \(110A\) clean-round waiting bound yields a fixed previously empty \(4B+4\) strip between their tips disjoint from K and all true ends. It is empty in its center throughout the remaining withdrawal. The two cores persist on opposite sides of this strip. A same-label forged distal site cannot change this order after restoration, since the restored interval connects to its core and cannot cross the empty center. The assignment at the late clean strobe therefore identifies the same authentic side and same exact label at the early comparison. A root farther out is separated by the nearer core and cannot be the source at the late strobe. Apply the projected-tip recurrence to this particular side, charging at most the single \(35B\) gain across J by the second sterile strip. This covers a raw-created fragment meeting a genuine source, since a postrepair connected member cannot be beyond the certified band. A gap first opened by the disruption is either sampled in its \(110A\) discarded rounds or already admits its own newly certified band.
\item Only one side has an authentic source. With no opposing root the single free-tip recurrence applies, using the second sterile strip to bound the gain over J.
\item Neither side has a source. At a clean strobe with H dead this contradicts the existence of a genuinely supported marker.

\end{enumerate}

These are the alternatives at the two sampled unexceptional times; no side attribution at an exceptional strobe is presumed. In particular a live-state label at a sampled damaged physical position is never treated as a certificate of descent.



Translate previous issuing periods of U' and V' to \([0,T)\) and \([b,b+T)\) with \(0\le b<T\), interchanging names if needed. Let \(l=\lceil v/20 \rceil \), \(h=\lfloor 3v/10 \rfloor \); the reference comparison window is \([r+l\) \(d_0,r+h\) \(d_0]\). For each positive issue the earlier density lemma covers the whole time interval \([r+l\) \(d_0,T+s_0)\) in its issuing coordinates with deficit at most C A at every date in its spatial home. Here \(C=2{,}000{,}000\); if necessary enlarge w so \(w>100(C+1000)\). This choice can be made after the constants are established. We have \(r+h\) \(d_0<s_0\), hence the entire shifted comparison window starting \(b+r+l\) \(d_0\) is in both issue-density spans (a fixed one-tick rounding margin may be deleted). Pick a common retained date outside S, deleting this additional interval of at most d dates. If the lower grids are different, lower nominal exclusion forbids physically overlapping occupied lower rectangles from the two tribes. If the lower grids are identical, the required lower age or address residues differ, because simultaneous equality of both residues at a common lower index would identify the two upper origins modulo their time and space periods. The same lower state cannot realize both. At most C A lower cells are missing in each home, so spatial overlap has length at most \(2C\) \(A+4\); otherwise a common occupied physical position would remain after removal of both sets of missing intervals and the two boundary roundings. If the homes are disjoint, their gap is already \(\le\)\(64A\). Their facing edges thus differ by at most \((2C+64)A+4\).

Suppose first that V was nominally dead. Choose a terminal marker inside V between \(14A\) and \(18A\) from its facing edge, avoiding K enlarged by \(100B\) and the at-most-\(2q\) exception. The available length \(4A\) exceeds \(200B+d+2q+10\). Use the last preballot row: the positive cut supplies occupancy off the common \(2q\) defect throughout this band, its vote cannot recruit into a vacant slot and any affected coordinate is in K. Consequently this is a genuinely supported terminal member. Since V is dead its own home is excluded as a positive input; any other true home is at least \(14A- 2\) coordinates from this marker, exceeding \(300B\). The remote alternative yields actual matching coverage across its \(.2v\) band at all retained dates of the shifted comparison window. Towards U its center is at most \((2C+82)A+6\) from U's facing edge; it extends at least \(.2v- (2C+82)A- 8\) into U. This exceeds C \(A+2\) by the choice of w, and the remaining width of U exceeds v. The band and U's density at a common retained date contradict lower exclusion. Thus V was live.

If U was dead, choose the matching terminal marker in U identically. Its own remote band now lies in the unshifted comparison window. The density spans of the positive issues V and V' begin at \(b- T+r+l\) \(d_0\) and \(b+r+l\) \(d_0\), respectively, and their ends are \(b+s_0\) and \(b+T+s_0\). They overlap since \(r+l\) \(d_0<s_0\). Their union contains the whole unshifted comparison window (use \(b<T\) and \(r+h\) \(d_0<s_0)\). Removing the remote-comparison date set for U (the density spans of V and V' have no excepted dates) and the at most d dates of S still leaves a common date, since \(h-l-2>.24v>220A+2w+21\) for the stated large w. At that date either V or V' supplies density over their common home. The identical spatial calculation contradicts lower exclusion. U is therefore live. Translating both tiles back by T preserves their physical overlap and the \(64A\) pad, proving inheritance.

\textbf{Uses.} Apply inheritance to genuinely overlapping U',V'. Positive predecessors force the preceding computing periods, which allows the separate full-home common-date exclusion using the long \(\operatorname{Shrink}(0)\) stages. The contrapositive supplies incoming clearance of a positive successor when its predecessor was proper dead or undisturbed live, once competing nonproper-dead tiles are classified. A positive current input alone does not assert service of a successor.
\end{proof}

In the next lemma use the earlier rank-j inheritance, seed, issuing, classification and held-flank implications as well as the rank-\((j-1)\) hypotheses.
\begin{lemma}[Overlap exclusion and properness of a matching middle]\label{lem:exclusion}
Under the rank-(\(j-1\)) hypotheses, distinct nominally live rank-(j) tiles in protected work rectangles do not overlap in space--time. The middle of three adjacent matching nominal positive tiles is proper live.
\end{lemma}
\begin{proof}
Suppose two genuinely overlapping positive homes U',V' of different origins existed at this rank. Their input faces and their predecessors, shifted back one whole epoch, have the same relative overlap. Lemma~\ref{lem:inheritance} supplies live U,V and, applied to them, their two live predecessors. In the ordering of the inheritance proof one foreign issuing period \([b- T,b)\) lies in the union of the two canonical issuing periods \([- T,0)\) and \([0,T)\). All three works compute, since each has a positive nominal issue. Foreign \(\operatorname{Shrink}(0)\) lasts more than \(.85T\). At least one of the two canonical periods contains at least \(.425T\) of it and less than \(.15T\) of this intersection lies outside that period's \(\operatorname{Shrink}(0)\). Exclude each home's first \(50A\) \(d_1\) dates, its at most \(N=200A\) \(d_1\) nonfull fine dates afterwards, and at most \(4w^6+20\) for the common slabs and four fine boundary rows. The excluded length is at most \(2(50A\) \(d_1+N+4w^6+24)\), a polynomial in w smaller than \(.275T\) for sufficiently large w. A common remaining physical date has both whole homes occupied. If the grids are displaced the overlapping physical spatial intervals contain a point on both occupied lower rectangles, violating lower exclusion. On a shared child grid two distinct higher origins require different child age or address residues at a common index, impossible for a single served state. Half-open rectangles with mere touching at a single endpoint are excluded by the overlap hypothesis.

Now take three consecutive matching nominal positives. All three prior works were computing. Their common full ordinary seed from Lemma~\ref{lem:seed} and two inside bridges keep the entire middle home and the two neighboring rims joined through the ordinary work except for the one localized response. Its true end sites start proper by the adjoining lower triples at the seed. A late split at K leaves its long proper old sides intact in computation; either both reach length \(q+3\) or the only short side lies next to K, of length at most \(q+2\). The paired repair of Lemma~\ref{lem:flanks} otherwise restores both in-home sides well before the ballot. If K straddles an edge use the independent backing of the adjacent home; grade one at the wrap does not carry \(\mathsf{Cn}\). The positive ballot preserves all unhit properly covered cells (Lemma~\ref{lem:ballot}). Hence \([1,P- 1)\) is present off one interval consisting of K and its possible adjacent short piece, of length \(\le\)d+\(q+2\)\textless{}\(2q\). For an early deep Shrink hit the at-most-two additional losses are healed by the buffered repair before the long computation. A different origin with a nonproper-dead tile touching the middle home's \(64A\) pad is classified as nominal live by its preceding issuing work (Lemmas~\ref{lem:ballot} and~\ref{lem:classification}). Since \(64A<P\), it would overlap the middle home or one of the two neighboring live homes in space and time. This contradicts the overlap exclusion just proved. The middle nominal positive is therefore proper. Neither separation nor this triple argument uses current-rank service.
\end{proof}

\subsection{Current service and proper successors}
The end-relief strategy adapts G\'acs \cite[Section~16]{gacs86}.
\begin{lemma}[Service of a proper endpoint under live holds]\label{lem:endpoint}
Let a lower site at space i, time a have a proper live nominal state. On any integer interval \([a,b]\), suppose at each \(t<b\) that the actual native instruction F evaluated on the three served lower operands at (i,t) prescribes a live output at i. Then at every t in \([a,b]\) the site is proper live and served, and for \(a\le t<b\) its next state is the specified output. Neither spatial neighbor is required to be live after the transition. The hypothesis is pointwise on F and thus includes a positive ballot.
\end{lemma}
\noindent\emph{Full-home specialization.} A full home of proper ordinary sites with a genuine common clock and a fresh \(\mathsf{Cn}\) reset remains full, proper and served up to its ballot; the same is true for a mid-period start. If every native ballot result is live, it stays so across that ballot and its scratch reset, independently of exterior neighbors.
\begin{proof}
Fix a step of the rank-j induction. Work on the common aligned rank-\((j- 1)\) grid, inside the domain where the preceding-rank organization conclusions apply at both times. When a rank-j exception is allowed this lemma applies only to steps so protected. At rank zero take actual unflagged transitions and identity service. The following uses lower proper-successor clause in exactly the form ``proper-dead or undisturbed live input has a proper nominal successor.''



Write x,g,y for nominal state, service indicator and served state. Suppose x(i,t) is proper live. Its definition includes incoming clearance, so it is undisturbed. Lower proper-live service gives \(g(i,t)=1\); consequently \(y(i,t)=x(i,t)\). Lower exact computation then gives \(x(i,t+1)=F(y(i- 1,t),y(i,t),y(i+1,t))\). This value is live by hypothesis. Apply lower proper-successor clause to the undisturbed input at (i,t). Its actual nominal successor is proper; the classification and its live value place it in the proper-live case. Lower proper-live service at \((i,t+1)\) serves it. The first application uses the given start; iteration up to b proves every assertion, and for \(b=a\) lower proper-live service alone suffices. No step applied the triple rule at time \(t+1\). At the physical base \(x=y\) identically and the unflagged rule performs F directly; no lower properness conclusion is needed there. All invocations below positive rank use the already established preceding-rank conclusions, never a same-rank assertion.

\textbf{Native full-home specialization.} On the old row suppose all P cells of one home have proper ordinary addresses 0 through \(P- 1\) and the same genuine clock, and are freshly reset at the start of a work period. Induct up to but excluding its ballot. The same induction applies from a mid-period start, which needs only ordinary addresses, a common genuine clock and a \(\mathsf{Cn}\) reset in addition to fullness and properness: the \(\mathsf{Temp}\)-removal clause exempts the outward true face even for an initially marked endpoint, and every other face is positively bonded. At a birth instruction an occupied destination ignores proposals. A true Grow withdrawal cannot target any of these ordinary addresses. A temporary retirement at an interior face sees a positive bond. An outer face at an extreme address is exempt in its outward direction, and the opposite face has a positive bond. For Shrink(i) each interior face has positive grade and a true outward face has depth zero, which fails the strict comparison with \(3iA\) for \(i=0,1,2\). At the start of every Purge the sites have \(\mathsf{Cn}\)=1. In its q propagation ticks a site at address \(k\ge q\) has \(\mathsf{Cn}\)=1 afterwards: after h ticks this holds for all \(k\ge h\) by induction through the exact integer increment from \(k- 1\). During killing, every \(k<P- 1\) has a strong right neighbor on the old row and thus cannot die. At \(P- 1\) the bit is one since \(P- 1\ge q\). These facts hold successively on all q killing rows since each row retained all sites and killing does not change \(\mathsf{Cn}\) on a survivor. Normalization fixes these addresses; other pre-ballot instructions leave incumbents live. The established pointwise lemma transfers each native hold to a proper served next row. Synchronous induction over the entire home therefore proves all P sites served until the pre-ballot row for any outside neighbors. If each native vote at the ballot returns live, the same pointwise lemma transfers all P votes across that ballot, even when the two neighboring homes vote dead. The ballot's scratch reset, ordinary addresses and zero age give the initial condition for another work period.

For a newborn site, proper deadness of its predecessor and the lower successor clause yield properness once the positive birth is known. The endpoint lemma applies to protected updates with an established positive native result. It does not assert that an eligible birth occurred, that a physical exception leaves a site proper, or that an arbitrary ballot votes live.
\end{proof}

\begin{lemma}[Terminal competitors at an ordinary end]\label{lem:competitors}
Use the issued all-date density, lower overlap exclusion, and the terminal near-source or remote-band alternatives, with the proper-current-input or known-positive-successor anchor described in the proof. For the four-site and four-row stencil about a true end, at every query date of a recovery search whose request lies in \([s_0,s_1+20v d_0]\), through the following \(3v d_0+2w^6+12d_0\) search (all such dates are below \(T-2\)), terminal and early-age contenders are excluded off one common physical raw/repair date deletion, of length at most \(w^6+2d+50\) including its four-row halos. Remaining genuine contenders are attached active Grow sources whose positive homes are disjoint from the canonical \(64A\) pad and, given the retained middle wall, lie on the far side of the pad from the home (strictly to the right of the pad at a right end), with intermediate Grow age and epoch start strictly between 0 and T. The end's retained near and middle walls are hypotheses up to the first excessive-loss test.
\end{lemma}
\begin{proof}
This is the terminal and early-age elimination used in the uniform end-relief argument, prior to selecting a single active competitor. It uses all-date issuing density (Lemma~\ref{lem:thin-density}), held cores, the whole-tribe thin trace and the near/remote/short alternative in Lemma~\ref{lem:remote}, the protected lower exclusion, and backward inheritance in Lemma~\ref{lem:inheritance}. None assumes current-rank service or directed advance. The right end is P and its prospective query strip \(Z=[P- 6A,P+15A]\). For a left end use the physical interval \([- 15A,6A]\) and distances from zero; spatial overlap, held cores, thin tracing and birth counts below do not depend on pressure propagation direction. In proofs using Purge on that side, it is still the rightmost position with q old predecessors that prevents the purge of a long run.

\textbf{Claims and surroundings.} The canonical work period of U is \([0,T)\). The query has time t between \(s_0\) and \(s_1+20v\) \(d_0\), allowing the following \(3v\) \(d_0+2w^6+12d_0\) search. These dates are all below \(T- 2\): after \(\operatorname{Shrink}(0)\) the later two stages alone have 4Pq \(d_0\) ticks, exceeding this extension. The following anchors are distinguished.

\begin{enumerate}
\item A proper live U has its own current incoming clearance. Up to the first excessive loss its near-end ordinary wall is sustained. Its issued-source preceding \(C_*\) through Grow and its preceding productive epoch's density give actual canonical material in a band of width C \(A+O(A)\) near this edge at the earlier comparison dates.
\item A proper-dead or incoming-clear live U has a known positive successor U'. Disturbance inheritance gives U' incoming-clear. Only if U is live does one additionally have direct current-time clearance of U. U' supplies density of U's actual productive period from age \(r+\lceil v/20 \rceil d_0\) onwards; the simultaneous full stage-0 seed supplies the near wall at the start of this search.
\end{enumerate}

Take w large enough that \(v/5>(10C+500)A\), where C is the fixed density constant. A nominal blocker on a lower displaced tile is converted to a served native predecessor by its nonzero nominal issue and the native no-spontaneous-birth rule. If a native increment at the boundary crosses a ballot, use the incumbent at the pre-ballot row. Radius one in the foreign grid, the pad \(64A/P<1\) and two displaced-grid roundings cost at most four canonical rows and four sites each way. Use the following argument for the served observations with that extra buffer. Any physical wrong-age occupant in the opposite neighbor of a proposed vacancy must be considered a served contender as well.

\textbf{Early, raw, and shifted-cut observations.} Delete the one common physical date interval running from four ticks before the common raw rectangle S to \(w^6+d+12\) past its last date, with an additional four-row halo. This covers all actual raw-only new families at a late age and their bounded source repair. At a foreign vote at b an older perturbation may give a delayed small source in K. If its new age is \(\le\)\(5q+2d\)+12 then trace the particular served occupant backward to the genuine final incumbent of \([b- T,b)\), unless this meets the deleted raw dates. Each step is a prevailing member or a birth from an adjacent served donor; across the vote a vacancy cannot give a live ballot. In these \(\le\)\(5q+2d\)+12 ticks there are \(\le\)4 native recruitment dates for the fixed foreign phase since \(5q+2d+12<4(2q+2)\). Consequently the ancestor lies within four sites of the query. The old epoch meets \([0,T)\) because \(0<b<T\) \((t\ge s_0- 4\) and \(s_0\)\(\gg\)\(5q)\). Its terminal-age incumbent is considered below. After its full initialized Purge the independent small source has disappeared. If instead a genuine current source is responsible and age \(a\le w^6+6q+20\), its authentic frontier could advance beyond its ordinary home by no more than \((B+1)(ceil((w^6+6q+20)/d_0)+2)\) plus \(d+8\) for a short forged group joining it. This counts every provisional birth as permanent and every true attempt once per cycle, and includes partial cycles. As \(d_0\ge 4w^5\) the bound is \(\le\)\(32A+400q\) for \(w\ge 100000\). In Z with the four-site buffer this forces the candidate's positive home to intersect the canonical \(64A\) pad. Its start \(b=t-a\) satisfies \(0<b<T\) and its epoch meets both \([0,T)\) and \([T,2T)\). This contradicts the pertinent incoming-clear anchor in cases 1 and 2. This treats each early phase at its actual time without deleting its individual delayed ballot time from the prospective common window.

\textbf{Terminal age} means \(a\ge r+.7v\) \(d_0\), before the candidate's ballot; permit the previous epoch obtained in the shifted-cut case. Its period begins at \(b\le t\) and the query date, including its four-site ancestry, exceeds \(s_0- 5q- 2d- 20\). For a live spatial home use its core; for a dead spatial home apply the remote-marker result, yielding a live home within \(300B\), actual wide coverage at every retained comparison date in \(b+r+[.05v,.30v]d_0\), or the raw-only transient already removed. The bounds stated here use the nearest edge of the true home and permit a fixed site of its held core within \(300B+12A+4q+4\) of that edge, outside the bounded wound on retained dates.

When U is live and incoming-clear, a positive foreign home within \(300B\) of Z (including the observed spatial home if live) directly violates its pad since \(15A+300B+10<64A\). Its period still intersects \([0,T)\). Otherwise the foreign wide band contains canonical positions near the right end. The earliest possible comparison is at least \(s_0- 5q- 2d- 20- T+r+.05v\) \(d_0\). In the preceding epoch this is later than \(- T+r+.05v\) \(d_0\), so the density of U's positive preceding issue is available, or, if in \([0,s_0]\), the held core of U is available. Between \(s_0\) and the starting ordinary strip use the known U' density if U' is assumed. For the current-input service branch the near-wall bootstrap instead supplies an actual member in \([P- 4A,P- A]\) up to the first-loss test and the wall retained over the pending search supplies it on later dates. With a deficit of at most C A one obtains an occupied site in the portion of the foreign \(.2v\) band inside U by counting \(>C\) \(A+2\) whole canonical sites. At a held-core or near-wall date there is such an occupied site directly. Remove the at most two earlier canonical source-repair date sets (issuing density itself has no excepted dates) and the candidate's wide date set from the order-\(.25v\) \(d_0\) comparison. Their total length is at most \(3(220A+2w+20)d_0+3w^6+20\), smaller than \(.24v\) \(d_0\) for the chosen w. Disjoint lower grids cannot occupy the same physical point; two same lower grids requiring different native time/address residues cannot supply both values either.

When U is proper dead and U' positive, select the canonical deep productive late marker in U at distance between \(14A\) and \(18A\) from P, avoiding K enlarged by \(100B\). Its previous-ballot existence is proved in Lemma~\ref{lem:inheritance}'s inheritance argument by the positive output's cut coverage. The distance to every positive current canonical home is \(>300B\) since its own home is dead. Hence there is a canonical wide band in \(r+[.05v,.30v]d_0\). Consider a foreign terminal contender.

\begin{enumerate}
\item If its alternative is wide and \(b\ge 0\), the foreign comparison dates fall after the beginning of the canonical issuing-density interval for U'. They also fall before the observation by at least \((.7- .3)v\) \(d_0\), so before T. Its band reaches more than C A inside U. The same common-date count contradicts exclusion.
\item If its alternative is wide and \(b<0\), the canonical wide comparison is later than \(b+r+.05v\) \(d_0\) and earlier than the observed date, which exceeds \(s_0- 5q- 2d- 20\). At each such canonical comparison date the whole-tribe foreign terminal thin lemma gives an occupied point within \(100B\) of the observed coordinate. This thin implication covers the final withdrawal as well as the following ordinary part and remains valid across normalization by the sterile-band proof for the whole tribe. The marker's distance to the canonical deep center is \(\le\)\(33A+100B\)+12\textless{}v/5, contradicting the canonical band.
\item A nearby foreign live home with \(b>0\) overlaps U' and the \(64A\) pad, directly contradicting U' clearance.
\item For a nearby foreign live home with \(b\le 0\), the canonical comparison dates are after b and before the foreign observation. Until its own normalization \(b+s_0\) use a point of its old held core at distance \(\le\)\(300B+12A\)+\(4q+4\) from the observed coordinate. If the observation is later than that normalization, on its ordinary prefix from \(b+s_0\) onward use the whole-tribe thin trace of the observation itself, within \(100B\). If the observation is during Grow, use just the held core. Thus on each retained canonical comparison date there is a foreign point at most \(300B+45A+4q+20<v/5\) from the deep center, again contradicting exclusion. At \(b=0\) no overlap of successive half-open parent epochs was assumed.

\end{enumerate}

In every case the earlier witness date may depend on this specific competing origin. The conclusion holds for every terminal or early contender at every prospective candidate time outside the \emph{single physical} deletion, and requires no union of historical foreign bad dates over all clocks. A genuine remaining contender must have age strictly between the early cutoff and \(r+.7v\) \(d_0\) and epoch start strictly between zero and T because \(t\ge s_0- O(q)=r+v\) \(d_0- O(q)\). It is an attached positive Grow source away from its early/delayed and actual repair. Its live home cannot touch the canonical \(64A\) pad. At the right end it must approach from the right: a left home outside the pad would have a connected arm passing through the old canonical middle wall (whose independent strip selection is given in Lemma~\ref{lem:calendar}), forbidden by lower exclusion or incompatible physical native residues. This is the input to Lemma~\ref{lem:calendar}.
\end{proof}

\begin{lemma}[A common clear recovery interval]\label{lem:calendar}
At a true end, let the request be the first date of a one-ended ordinary recovery search: an existing end deficit may grow at \(\operatorname{Shrink}(0)\), an interior edge position is missing, or the proper-input bootstrap has its initial shallow shortage. Require its date in \([s_0,s_1+20v d_0]\). Assume source attachment, protected lower exclusion, ordinary precommit, the preceding strong issuing cut, the terminal exclusions of Lemma~\ref{lem:competitors}, and the old near wall and interior support for its proper-current-input or known-positive-successor anchor. Then either the first full ordinary Heal after the request whose complete pass and four prior rows avoid the common raw/repair deletion of Lemma~\ref{lem:competitors} is clear and is completed within \(w^6+2d+4d_0+50\) ticks of the request, or an interval of length at least \(0.05v d_0\) starting within \(3v d_0+2w^6+12d_0\) of the request is simultaneously clear throughout the \(100q\) end strip and its stencil halos.
\end{lemma}
\begin{proof}
The hypotheses give a fixed near wall and an interior support, using the two alternatives for the current input and the known positive successor, respectively. The distance to an actual donor when the later pass starts is a separate premise for filling: Lemma~\ref{lem:end-ledger} supplies an enduring \(q+1\) guard within depth \(97q\). An initially shallower missing strip does not by itself give that later bound. All rank-level clearances, including the old-row halo, refer to original protected physical rectangles.

\textbf{Walls independent of successful births.} For the right end work on \(Z=[P- 6A,P+15A]\) with the four-site/row guard of that lemma; the narrow tested strip for actual births and ancestry lies inside \([P- 100q- 20,P+100q+20]\). At the ordinary full seed, or in the proper-input bootstrap before excessive debt, the region \([P- 4A,P- A]\) has old proper ordinary pieces on both sides of a single possible short missing region. An authentic initial local wound and a later K use that same spatial column; any initial edge shortage \(<60q\) is outside this region. Remove K, the possible single \(2q\) initial support exception with its neighbors, and, if using the temporary expanded wound before its clearance, its band of length \(\le\)\(4B+80q\)+80+d. The three-A region has at most four remaining pieces and retains more than \(2A\) sites, hence an intact piece of length \(>.1A\). Trim the exterior two sites to ensure proper old incumbents. If one piece boundary is an active probe rather than an old site, take its old interior before this event by the core/old-run bounds from Lemma~\ref{lem:attachment} and the ordinary seed. Old proper, unmarked ordinary positions survive from one true Shrink to the next by the \(q+1\) strong-run calculation: their rightmost \(\mathsf{Cn}\) retains q predecessor support and each earlier site has a right neighbor. One Shrink can remove at most one layer at each end. During \(3v\) \(d_0+2w^6+12d_0\) there are at most \(ceil((3v+2w^6/d_0+12)/q)+1<4q+4\) actual Shrinks, since \(v/q=q\) and \(d_0\ge 4w^5\). Taking the original subrun trimmed by \(4q+6\) each way leaves a nonempty fixed wall in \([P- 4A,P- A]\). Source continuation and lower-rank properness apply to that wall throughout the prospective search. A canonical middle barrier in \([P/4,3P/4]\) is obtained by the identical selection and trimming on that much wider interval. Both barriers therefore exist without already having a successful one-sided repair.

Skip the common physical raw/repair interval of the preceding lemma if necessary and test the next full ordinary Heal together with its four prior rows. At most \(w^6+2d+4d_0+50\) ticks are needed for this first trial even if the one deleted interval cuts the first candidate pass. If it is clear this is the immediate calendar alternative. Given a surviving in-home guard within depth \(97q\) on this pass, each step to the true edge is mandatory and needs at most \(97q<B\) opportunities, as proved at the end below. Otherwise take an actual incompatible served witness to failure at a particular query time \(t0\), including the predecessor of a conflicting nominal test. All four-halo enlargements are already in Z. By Lemma~\ref{lem:competitors} this witness belongs to a genuine active Grow source Q, whose home lies strictly to the right of the \(64A\) pad. Its age exceeds the early cutoff and precedes the terminal threshold. It is attached to an enduring core on this clean date. In particular its interval from the tip to the core fills \(P+[.25P,.65P]\), up to a four-child rounding. This follows whether its home starts inside this strip (its core's far edge is beyond \(P+64A+P- 12A- 4q)\) or farther to the right (its authentic arm must traverse the whole strip).

Let \(L=ceil(3v\) \(d_0+2w^6+12d_0)+20\). Since \(P=5\)\(\cdot\)\(2^w\) and L is polynomial in w, choose w so \(10L+d+100<.01P\). Remove K enlarged by \(L+4\) to both sides from \(P+[.25P,.65P]\), and trim the ends by \(L+4\). The remaining length exceeds \(.35P\) in at most two intervals, one of which contains an interval of length \(\ge\).\(1P\). In its interior retain an interval D of width \(.05P\) so every point of D during the next L ticks has its backward cone to the witness date in the original occupied region avoiding K. The date of the recovery below is before Q's own next normalization and ballot. At each Grow update in this cone, a strict interior of three contiguous members of Q is proper, has exact grades, and prescribes a positive hold. This is deterministic also while the source is interrupted elsewhere. Induct on ticks, shrinking the protected occupied witness by one site at each end. Since the cone intervals were chosen with margin L, the whole of D is continually occupied by the same genuine exact label. This is the rear wall of Q.

\textbf{Universal exclusion over a fixed horizon.} Fix an arbitrary later date u in the selected portion of Q's current withdrawal (specified next), excluding the common physical slab and a four-row guard. Inspect any actual incompatible member in Z. At u apply Lemma~\ref{lem:competitors} anew. It rules out the tentative member if it is early, terminal, or a pure postcut residue. A residue at someone else's fresh vote is judged by its actual past incumbent; the historical comparison dates used to exclude that incumbent may vary, but this does not change u. In the remaining case it belongs to a connected active source with a positive home disjoint from the canonical pad, as at \(t0\). A left home would have a connected interval crossing the canonical middle wall. A right home must reach past D's right end from Z for the same spatial reason that Q supplied the original D. Thus its interval contains D at this very date u. Different lower grids violate lower exclusion there. On a shared lower grid, at a given index (x,u) the actual lower state has a single age and residue, so a second different upper phase cannot match both. If both lie on the same phase and grid, its exact extended address at x in D fixes the source base \(x-\pi\); it is Q's base, so the second source is Q. A contender already inside the wall when certified, or one born on its near side later, is subject to the same argument at u: the new age/old-vote cases have already been eliminated, a still-unattached raw lineage requires the removed common slab, and every remaining member must be connected to its own far core. No transport of a distinct source THROUGH the wall between \(t0\) and u is assumed. This proves the universal statement: in the single chosen clean portion, for each u and every conflicting member in Z (and its specified halo), the only possible source is Q.

\textbf{Retraction and uninterrupted window.} If Q is in attack, its next withdrawal begins within \(2v\) \(d_0+d_0\). If in withdrawal with at least \(.3v\) \(d_0\) available, begin now, allowing one partial strobe. If in a nonfinal withdrawal with less remaining, the next full one begins within \(2.3v\) \(d_0+2d_0\). This last case cannot occur in the final withdrawal because the source would already be terminal. In all cases restrict to the first \(.3v\) \(d_0\) of this available portion and remain inside Q's same Grow. Delete the physical raw/recovery interval(s), with integer four-row margins and round to complete strobes. Even allowing two deleted intervals of total length \(2w^6+4d+100\), at least one of at most three pieces has length greater than \(.09v\) \(d_0+2d_0\), since \((.3v\) \(d_0- 2w^6- 4d- 100)/3>.09v\) \(d_0+2d_0\). Choose its first complete strobe and retain \(.09v\) \(d_0\) of clean Q withdrawal (drop endpoint fractions with an additional \(d_0\) guard if needed).

At the selected strobe Q's left tip is to the right of the retained near wall, since otherwise Q's connected interval would occupy a position of that wall. For a tip between \(P- 4A- 4\) and \(P+12A+4\), its address is negative because its home starts beyond \(P+64A\). If a neighbor on its left has positive grade, it either has the immediately preceding integer address, in which case it is another member of Q contrary to extremality, or its address exceeds \(P/2\) by the grade-one definition. In the latter case its own base is less than \(P/2+12A+4\). If its positive home meets the canonical pad it violates the selected anchor (same active clock \(b>0)\). If the home is wholly left of that pad its authentic arm must pass through the canonical middle wall. In particular an honest left neighbor from a second phase cannot pin Q; positive grade requires the very same clock and residue alignment. A raw-only supporting label could last only to its next full Purge and is absent in this chosen clean portion. Therefore at each true Grow there is a zero face at its leftmost extended site or the tip is already clear of this entire strip. The B attempts of an unjoined provisional advance were marked and fully withdraw before this true Grow, as long as the entered front is in this guard, where no true-end latch exists. A site at most B beyond the guard cannot evade the argument by carrying an untested tip through the guard: reaching the narrow tested strip would first require traversing the additional gap of order A outside it. The first \(\lceil .02v \rceil \) clean complete withdrawals exceed \(16A+B+15\). They carry the tip past \(P+12A\), or find it already beyond. Choose a \(4B+4\) virgin band near \(P+8A\) inside this now empty margin, avoiding all true endpoint residues of Q's tribe. Such a band exists, for instance choose within \([P+7A,P+9A]\), since its length \(2A\) is \(<P-2\) and the possible paired true ends remove at most two places and split it into at most two nontrivial intervals. This band contains no competing matching source: the universal exclusion just proved applies to any Q-tribe member in Z that could positively pin/repopulate the narrow strip, since that tribe differs from the canonical one. The whole-tribe sterile-band argument prevents later return through its center at a strobe or fine tick during this withdrawal.

There remain at least \(.06v\) \(d_0\) clean ticks, which contain a contiguous interval of length \(.05v\) \(d_0\) after deleting four-unit rounding and four preceding halo rows. At all its dates every competing served member is excluded by the universal assertion and Q's retraction. A conflicting nonproper-dead nominal tile at any needed prior stencil would have an actual conflicting served predecessor within the already cleared halo, impossible. On the preceding row each of the nine ordinary lower cells was proper-dead or live incoming-clear; lower proper-successor clause makes its successor proper. The initial search + \(2.3v\) wait + \(.3v\) chosen portion totals \(<3v\) \(d_0+2w^6+12d_0\) even after the stated rounding. \textbf{Conditional filling remark.} The current nine cells are incoming-clear and their true central births are consequently mandatory. A same canonical clock/residue neighbor has the unique actual ordinary address at its position by off-Grow normalization. At the actual outer donor-age test it has the correct age if occupied. Across the true wrap its out-of-home ordinary proposal is illegal, and any legitimate in-home proposal installs the intended matching address. A \(q+1\) old backing persists through every intermediate Purge, whether advancing to the left or the right. Each of the B attempts recruits the next vacancy or joins an already properly aged in-home incumbent, so the distance from the enduring guard to the true endpoint decreases at least one per successful attempt. With the separate \(97q\) guard supplied by Lemma~\ref{lem:end-ledger} this repair takes at most \(97q<B\) attempts, irrespective of the end depth at the earlier first-loss request. At its own true endpoint the new trial sees a positive inward face and the outward directional exemption; it persists at retirement.

For a left-end recovery all temporal and home comparisons use \(Z=[- 15A,6A]\). A right-side hostile root would have to cross the canonical middle. The actual attacker lies left of \(- 64A\) and has positive extended rightmost address \(\ge\)P. A protecting right neighbor of a different exact base must have address \(<P/2\) at the high-to-low wrap, so its home meets the pad or its arm crosses the canonical middle from the right. Select the permanent rear wall as the reflected subinterval of width \(.05P\) in \([- .65P,- .25P]\), with the same collar avoiding K; the near wall lies in \([A,4A]\). Replace the empty band location by \([- 9A,- 7A]\) and require the tip pass \(- 12A\). The same distances and withdrawal counts apply. In the Purge and one-sided fill arguments continue to use the rightmost incumbent's q strong predecessors, exactly as in the right-end proof.
\end{proof}

\begin{lemma}[Current service and a proper positive successor]\label{lem:end-ledger}
In a protected rank-(j) work period with the lower-rank hypotheses and the previously issued cut, an incoming-clear proper-live input is served. If a proper-dead or incoming-clear live input has a positive issue, assume its work is computing with the full seed and that the strong positive cut is available. The positive nominal successor is proper live. The two conclusions use the respective end-relief anchors, and the first uses no hypothesis on the successor.
\end{lemma}
\begin{proof}
Use parent home \([0,P)\), write depth \(k=x\) at its left end and \(k=P- 1- x\) at its right end. A ``missing end'' below concerns the retired genuine front. Provisional B-long probes are never used as donor anchors. Inputs for this deduction are Lemma~\ref{lem:endpoint} (applied to the rank-\((j-1)\) children), the all-shift lower organization clauses, actual previous issuing cut, Lemma~\ref{lem:attachment} for genuine source intervals and postpulse holes, the precommit common full seed of Lemma~\ref{lem:seed} in the known-computing case, the buffered two-wall repair for previously occupied interior regions in Lemma~\ref{lem:flanks}, and the two distinct relief anchors in Lemma~\ref{lem:competitors}/Lemma~\ref{lem:calendar}. The latter supplies its common calendar for births testing sites of depth up to \(100q\), including prior stencil rows. Write \(W=3v\) \(d_0+2w^6+12d_0\) for one complete search-and-fill time. These are conditional induction inputs; in particular a positive successor is not assumed in proving service of a proper current input.

\paragraph{Starting material, and two different depths}

For current-input service start at the genuine issued cut of a proper live parent. In \([1,P- 2]\) all positions except a span H of length \(\le\)\(2q\) match. Outside K their surviving true addresses are the ordinary addresses and their scratch is the specified positive reset. A run avoiding H and K and their one-site collars has proper lower incumbents by the matching triple. A proper unmarked run of length \(\ge\)\(q+1\) survives every Grow Purge against arbitrary exterior by its rightmost site's q strong predecessors; all other sites have a strong right neighbor. Incumbent ordinary addresses are exempt from true Grow withdrawal and no incumbent acquires a \(\mathsf{Temp}\) mark. Lower proper successor and service turn each genuine positive transition into a proper served next row. Thus such initial intact strips are available for the first ordinary-work search at \(s_0\). An old true edge itself need not have been proper at the initial cut. A real earlier gap well inside this near-full input has the preceding genuine issuing geometry: if no raw has yet occurred in the buffered history, the previous ordinary interval and ballot cannot create such an internal gap. The ballot has no births and all strict interior proper occupants of the single interval take its positive decision. If a cut-induced defect actually exists in the interior, it is in the single K wound, which Lemma~\ref{lem:attachment} repairs in Grow before normalization unless the actual event is too late. This is where free formatting alone would be insufficient.

For proper successor start instead at the genuine \(5d_0\) fully occupied, clean \(\operatorname{Shrink}(0)\) seed of a \emph{computing} parent whose next issued state is known positive. Choose a \(\mathsf{Temp}\)-reset instruction in this window and the following clean full Purge. Every strict interior of the home was proper before and after this reset, by lower triples and lower proper-successor clause. Both true endpoints may be improper. Here H is empty and the argument starts at this seed rather than at \(s_0\). Backward disturbance inheritance (Lemma~\ref{lem:inheritance}) gives incoming clearance of the positive successor when the input was proper dead or undisturbed live, as required by the second relief case.

A search of length W together with waiting at most two extra ordinary rounds for a complete reset has fewer than \(4q+4\) native Shrink inspections: the spacing is \(qd_0\) and \((W+2d_0)/(qd_0)\le3q+(w/2+14)/q\) since \(d_0\ge 4w^5\). Two such searches and the short pre/postpulse interior cleanup \((\le 2w^6+4d_0+20\) fine steps) cost at most \(n=8q+12\) inspections, counting a partial inspection at each of the three interval ends. On a certified connected ordinary interval, an arbitrarily long \emph{intervening} stretch costs no interior losses: at a true restored end the outward face is exempt and every other face through the attached core has positive grade. The distinct reason that at most two searches can intervene is verified below.

For any contiguous old candidate strip of \(L=q+1+2n=17q+25\) proper unmarked sites missing K, each of these n inspections can remove at most one site from each end. Between inspections the \(q+1\)-residue proof prevents all native Purge losses, and \(\mathsf{Temp}\) never acts on those old sites. Induction through the inspections leaves its central \(q+1\) as a fixed old backing. If the time interval starts with a partial Purge the old strip is present through that sweep's clean prefix; if crossing normalization start afresh at its actual first initialization. When only one prospective search and K are relevant, the smaller \(L1=q+1+2(4q+6)=9q+13\) works. In material starting at depth \(p+1\) two such strips separated by \(d+3\) allow a K-free one with closest fixed point of depth at most \(p+1+L1+d+3+4q+6=p+13q+d+23\le p+14q+d+22\). An H and K hazard together use three candidates separated by \(2q+d+5\).

For the two-search bootstrap select three length-L strips starting at depth 4 and separated by \(2q+d+5\) unused sites. The one-cell-expanded H hits at most one, and the one-cell-expanded K hits at most one. The closest fixed point of a usable third candidate has depth at most
\(4+2(17q+25+2q+d+5)+(8q+12)=46q+2d+76<47q\).
Every site of the selected strip avoids the true endpoints. The same selection at the other end works simultaneously. Before either end's first excessive loss, these retained witnesses prevent any retired deficit from exceeding \(47q<60q\) in the proposed one- or two-search interval, even if no birth worked during its search. This justifies the \(<60q\) depth invocation prospectively from real old material, rather than assuming an arbitrary \(60q\) deficit cannot increase at all during W.

A second selection separates the endpoint used for the \emph{relief calendar} from the farther guard that will actually fill it. If K meets depths 0.\(.58q\), its farthest affected depth is \(<59q\). Take two length-L strips starting at \(60q\) and separated by \(2q+5\). They both avoid K, one avoids H enlarged by one on each side, and both end before \(97q\). A fixed \(q+1\) guard therefore lies beyond K and depth \(60q\) but before depth \(97q\). If K first meets this end beyond \(58q\), use two such strips starting at depth 3, separated by \(2q+5\). Both end before \(37q\), and one avoids H. In this distant-K case the genuine interior around K is handled by the next paragraph. With no K use either selection. All selected runs have room for lower stencil halos in the \(100q\) zone. During the clean calendar each surviving chosen guard can grow all the way to the true end in \(<97q<B\) attempts. Before each attempt there is a complete initialized Purge; an additional correctly aligned occupied piece surviving it is backed, and each newly occupied point attaches by an exact ordinary bond. All queried lower inputs were incoming-clear on the preparatory row and are proper by the rank-\((j-1)\) proper-successor and service clauses of Proposition~\ref{prop:organization}\ref{org:successor}. Every legal in-home birth is consequently realized. At completion the new outward face is exempt, its inward face joins the guard, and every interior bond is positive. This produces every position in the interval from the true end through the fixed deep guard. It does not count an isolated provisional extension into an open interior cavity as a secured end.

\paragraph{Why an interior wound costs no extra search}

The near-K case just considered fills past the affected column to an old ordinary guard beyond \(60q\). If K is farther than \(58q\) from both ends, the near-end bootstrap witnesses lie at depths \(<47q\). Thus there are \(>11q\) sites between such an extreme witness and K. On the last actual full inner region before the pulse, the two sides of K each have a \(q+9\) backing within \(q+15\) of its corresponding edge; take their nearest fixed interior \(q+1\) sites after trimming three on each face. The sites in these deep guards (more than \(58q- q- 15\) from either competing end) were in the initially near-full proper interior, or in the genuine full seed. If K is a new later pulse, no previous unrelated deep interior hole exists; if K gave a delayed cut, Lemma~\ref{lem:attachment} supplies its real previously occupied bracketing geometry. The near-end deficit \(<47q\) cannot have reached either guard. An unhit ordinary part between the previous edge and this raw location has no mechanism for a first clean interior split (positive bonds, lower properness, and a lower hold).

Across an actual pulse these guards bound the previously occupied internal wound. They lose at most two faces to scheduled Shrinks before a complete initialized ordinary repair; the elapsed time is \(<w^6\), and \(w^6/(q\) \(d_0)\le 1/(4w)\). In the last Grow/first ordinary seam allow one final withdrawal on extended neighbors. The ordinary old positions of these guards cannot be removed by it, and normalization preserves their strict bonds. The guarded cavity has length \(\le\)\(d+20\) to its facing old supports. Lemma~\ref{lem:flanks}'s two-wall clearance provides the nine simultaneous stencils after the common alien reset and its ancestry rows; remaining true in-home bids close that cavity in a clean B-pass. If a raw wound was still open across normalization, the same fixed ordinary guards were present before the pulse on the Grow side; reset \(\mathsf{Cn}\) on ordinary entry and run that same repair. If the initial defect was a prior cut far earlier, the established Grow repair has already completed. A wound within \(58q\) of only the other end belongs to that other end's ledger, separated from the present end by \(P- 116q- d\). This covers the full interval between the two protected end regions. A distant pulse therefore does not leave a permanent inner break from which a restored near rim could erode. Only the bounded cleanup interval is charged to its neighboring old strip. The local argument uses actual pre-event occupancy of the inner region; it is not a statement about every interval modulo an arbitrary spatial error mask.

\paragraph{How many waits an end can incur}

On clean ordinary work before its first effective repair, a previously improper original true endpoint can disappear freely one time. If it never disappears and has continuous, properly backed inner material, it causes no lost interior. Otherwise request a relief pass when the existing end deficit could grow at \(\operatorname{Shrink}(0)\) or when an already missing interior edge position must be restored. At the outset it is also legitimate to restore the old shallow pre-Shrink shortage in the proper-input case. During the requested wait count native losses against the old backing above. The age/deficit and the prospective near/middle walls of Lemma~\ref{lem:calendar} are available up to the first claimed excessive loss. That lemma supplies the complete common clear Heal. A repair is considered secured when the end is actually continuously occupied into its old ordinary material and all these preexisting sites are proper or newly born proper. If a newly filled site is still marked, its complete positive inward bonds and the true outward exemption preclude clean retreat. One more clean \(\mathsf{Temp}\) reset is allowed inside the long clear interval. In the immediate-clear alternative, if the pulse intervenes before that reset, include that case in the interrupted first search below.

From a secured rim to the next actual pulse, no second free disappearance is possible at this end. Its old lower cells were proper on the preparing clear row or were proper births, so a hold remains served and proper by Lemma~\ref{lem:endpoint}. Its outward face has zero depth, its interior is strongly joined through the old backing, and a clean Purge cannot propagate a zero through that backed run. A previously shallow frontage cannot reopen spontaneously. In the near-K case the secured rim extends through the only possible wound to the fixed guard beyond it. In the distant-K case any actual interior cut is dealt with in its own \(\le\)\(w^6\) cleanup; the separating \(>11q\) interval and fixed backing prevent destruction of the near-end source. Beyond these transient regions the occupied interior has no first native split. Simultaneously argue for both ends until the first excessive loss; the opposite extremity cannot traverse the interior without first violating its own \(<60q\) bound.

There is one exceptional physical interval. If it precedes the first search, that search uses a clear pass after the event and there is no further occasion to reopen. If it meets the search, the common calendar skips the event and gives a clear pass, or one can restart once after it if the attempted immediate pass was cut. If it first occurs after a secured repair, only that event can reopen this end: use one new pass after the pulse, with a subsequent initialized check. If it never affects this end, no second search is charged. Thus at most two waits are charged per end, including the case in which they are consecutive. During a healthy intermission the secured rim and its attached interior have no exposed inner face except the bounded distant wound treated above. Consequently the old candidate-strip accounting uses at most the n paid inspections even if the time between the two waits is macroscopically long. In a pre-event intermission all the near/deep candidate strips in the old uninterrupted interior lie on that connected ordinary body. Existing uncharged dates before the very first loss likewise have full or previously certified interior and no exposed face at these strips. This proves the prospective bootstrap: a first \(>60q\) loss cannot occur during one or both waits because the old near witness of depth \(<47q\) is still present; on the secured portions the end is restored. A first deep wound continuing beyond its allowed repair contradicts the distinct buffered-cavity argument. A single temporary raw revisit of K by the final ballot is accounted for below, not counted as a new pre-ballot erosion episode.

\paragraph{Service, final time cuts, and the single common ballot exception}

In current-input service an end deficit exposed during \(\operatorname{Shrink}(0)\) has a relevant search starting no later than its last inspection near \(s_1\). Allow \(<W\) for that request, another \(<W\) for a pulse intervening before its successful completion, and the explicit short alignment/cleanup delays. Their sum is \(<8v\) \(d_0\) for this w. A newly free original \emph{true endpoint only} after the last \(\operatorname{Shrink}(0)\) cannot expose another interior position by \(\operatorname{Shrink}(1)\) or \(\operatorname{Shrink}(2)\): depth 1 fails their strict \(3A\) and \(6A\) tests, and the long proper inner interval supplies its own Purge backing. Thus unspent later free losses touch only the two actual endpoint coordinates excluded in properness. If the pulse occurs after the previously needed requests have completed, take that now prepared interior as the pre-event row. In particular when a new event first occurs later than \(s_1+8v\) \(d_0\) the earlier work has already established full \([1,P- 1)\) using the clean-history case. If it occurred no later than \(s_1+4v\) \(d_0\), any last resulting search is complete by \(s_1+8v\) \(d_0\) \((4v\) \(d_0+W\) plus its cleanup is \(<8v\) \(d_0)\); if it lies in the remaining crossover interval, earlier unhit initial repairs were done by \(s_1+4v\) \(d_0\) and the late-event calculation applies. All these initial search dates fall in the horizon \(s_1+20v\) \(d_0\) for which the relief comparison was proved. At \(s_c\) the proper-current home contains more than \(P/2\) matching positions: before the late event its prepared interior has length \(P- 2\), and under the late event the guard split below loses \(<2q\) positions in that interval (during a deep in-progress repair one can also retain both long pieces of that prior interior apart from K and two possible native Shrink layers on each facing edge). No assumption about that home's future simulated output entered the argument.

For proper successor take the prepared proper \([1,P- 1)\) of the full-seed case. Its ultimate positive ballot is known. If the physical pulse was spent during early preparation and fully healed, the only later free losses are endpoints. A delayed ballot effect may still occur inside K. Otherwise split the previously prepared stable old interval \([1,P- 1)\) at the late K. During computation each side of length \(\ge\)\(q+1\) is an unmarked proper ordinary run, so the whole side survives all clean initialized or suitably preinitialized Purges until the ballot. No scheduled Shrink remains in computation. When the pulse instead occurs in \(\operatorname{Shrink}(1)\) or \(\operatorname{Shrink}(2)\) and both sides have \(\ge\)\(q+40\) positions, the two deep guards have room for the possible two native Shrink losses before buffered repair. They restore the entire interior hole before ballot since computation intervenes for more than \(w^6\) fine ticks (if arbitrarily close to the stage boundary the initial computation segment supplies that time). On the outer faces of the two pieces the end depth \(\le\)1 is exempt under both late thresholds. If a piece is shorter than \(q+40\), it is between K and a single true edge. Enclose that whole side with K. The other long piece's exposed face towards K has depth \(\le\)\(q+40\)+\(d+2\) from that same true end, below \(3A\); its opposite true face is also exempt. None of its surviving old interior cells can die by later Shrink or \(\mathsf{Temp}\). Its rightmost Purge backing and all strict bonds persist. At most one side can be short, since \(P- 2- d>2(q+40)\). If K misses the home there is no such interior loss.

At the live ballot the strong cut gives the positive decision at every unhit proper preballot incumbent and the correct central payload outside K. Lower proper successor and proper service transfer that live vote, including an incumbent one position inside a missing true endpoint. The same interval I handles occupancy and payload: take K intersected with the home, enlarged by its one possible short adjacent interior piece and at most two discrete boundary positions, for length \(\le\)d+\(q+42\)\textless{}\(2q\). If no short piece remains, take just K; with no exception use the empty set. In particular combining a missing left flank with a separate independent right-flank exception is unnecessary. The two ends excluded in \([1,P- 1)\) are outside the support requirement. Together with the previously derived incoming clearance of the positive successor, this gives its properness. A dead successor follows instead from the actual small-group cut and needs no occupancy argument here.
\end{proof}

\subsection{Directed advance and initialized descent}
A \emph{willing donor} is an immediate proper-live neighbor whose true served state, in the served input triple with the proper-dead target, makes the prescribed rank-\(j\) map \(F_j\) request a legal live recruitment into that target, with its clock, address and opposite-neighbor tests satisfied. In a left request the proposed address is the donor's address plus one; in a right request it is that address minus one. The nine-input properness and incoming-clear hypotheses are imposed separately at each use.
\begin{lemma}[Directed clearance and a permanent startup anchor]\label{lem:startup}
For a proper-dead target and a willing proper live neighbor in the nine proper incoming-clear input stencil, the lower hypotheses give the clear preparation rows for directed recruitment outside the common short interruption. The donor has a permanent in-home backing and can secure its startup anchor in at most \(6d_0\) clean ticks, even if its original endpoint is absent.
\end{lemma}
\begin{proof}
Use time measured in child steps and the canonical target \(H=[0,P)\), with willing left home \([- P,0)\). A reflected willing donor uses x\(\mapsto\)\(P- 1- x\); Purge still propagates from the physical left, and both proposals, when legal, have the clock needed below. Write \(H_0=2q+2\) for the handle length, \(d=4w+8\). Choose \(w\ge 10^9\). We use the preceding-rank exactness/omission, nominal classification, proper successor, proper service and nine-proper-input birth clauses, under the actual protected stencils including shifts, with the conventions of Lemma~\ref{lem:transition}. Assume a single common enlarged exception K\(\times\)S of both spans \(\le\)d which contains all intersecting failed lower calls after rounding, and the actual previous-cut conclusion that a nonexceptional ordinary vote starts the next Grow with normalized addresses. The nine parent inputs are proper and incoming-clear, including the dead target. No exclusion based solely on vacancy of the target is used.

\paragraph{Directed clearance on the nine-home strip}

Use \(Z=[- 2P- 8B,3P+8B]\). The preparation rows omit the first \(6q\) rows and the slab \([\inf S-8,\sup S+2L+8]\) of at most \(17q\) rows, with \(L=4q+12\). Test actual matching occupants of every conflicting child phase or grid, with an additional eight child sites on both sides and four rows of stencil lookback. A genuine Grow occupant belongs to the home recovered by subtracting its extended address. Its distance to that home is \(\le\)\(1.1P+1\). In ordinary time its current home \([a,a+P)\) accepts no recruits from another ordinary home. At normalization its ancestors can be from \([a+iP,a+(i+1)P)\), \(i=- 2,- 1,0,1,2\): for \(i\le - 3\) the address would be \(\ge\)\(3P\) and for \(i\ge 3\) it would be \(<- 2P\). The nearer edge of each of those five intervals is at distance \(\le\)\(2P\) from any x in the current home. All these tests, also after a time-L predecessor trace of speed one with \(L=4q+12\) defined below, intersect one of the nine parent homes or its \(64A\) pad: at worst their nearer edge is \(8B+8+L+4\) beyond \(- 4P\) or \(5P\), which is \(<64A\). The four extra units allow the two one-child-unit roundings at each end. Thus if the candidate epoch intersects \([0,T)\), every authentic relevant foreign source in that epoch was proper dead by incoming clearance and classification. All lookups described here take place on the named protected spatial and temporal domains.

A proper-dead source on its normalized starting cut occupies a span \(\le\)q/2. No new same-phase member can appear except by a holding predecessor, an old neighboring donor or the event. The first complete native Purge ending after a specified last exceptional fine row finishes within \(2H_0\)=\(4q+4\) further transitions: wait at most \(H_0- 1\) for its initialization and use \(1+2q\) transitions for its three stages. Address normalization or an age-zero vote lies on an elementary instruction, never inside that Purge. For a short component with strong breaks at both ends, no same-clock birth occurs during the pass. After h information ticks any surviving \(\mathsf{Cn}\)=1 requires h immediately consecutive strong predecessors on the coordinates of the old rows. This follows by induction from the propagation rule testing the left bond and old left bit. With at most q original consecutive sites all bits are zero by \(h=q\). The rightmost remaining site in each component has a right break and is killed; successive rows erase every survivor in \(\le\)q ticks, even with extra omissions. An adjoining different exact owner has no strong link to it.

If the pulse meets the initial clearance, then from the initial row to the first complete Purge after its end there are at most six elementary instructions in this clock: the first four-q wait, d duration, and second four-q wait occupy \(\le\)\(8q\)+\(d+8\) ticks, less than \(6H_0\). Honest descendants of the initial span are inside its six-site enlargement on both sides, and those of the raw sites are inside K enlarged by six on both sides. These two regions have combined length at most \(q/2+d+24<q\). Each strong component has at most this length even if the two spanning intervals originated far apart. The initial complete Purge lies strictly inside Grow. If the pulse is later, initial survivors died already; then only the six-site enlargement of K can contain independent descendants on the first complete postpulse Purge. At normalization use the union of the five possible sources, which all died earlier or lie in that same enlarged K. Distinct short groups consequently cannot create one strongly backed union at normalization. If one traces an ordinary occupant back until this normalization, any intervening birth is in its own home; a previous empty home with no raw has only empty subsequent rows. This covers all sources that could have furnished the observed site or its earlier ordinary home.

It follows that at a queried time t a foreign member requires either a pulse within \(4q+8\) prior child units or a cut b with \(0\le t- b\le 4q+8\) (one lower unit of alignment on both boundaries is included). In the latter case trace to \(b- 1\) inside the foreign grid. Every new nonexceptional member requires a served predecessor within one step, and the vote cannot create an incumbent at a vacant site. Either this hits S or gives an actual previous terminal occupant in that halo. If \(t\ge 6q- 4\) and \(\ge\)4 rows of lookback are allowed, then \(b>0\), so the preceding epoch also meets the directed-clear work interval. Its age at this previous terminal row is \(T- 1\), greater than \(L=4q+12\). Apply the same preceding five-source calculation there. Only a raw event within L earlier ticks is possible. Consequently in canonical time, discard the first \(6q\) rows and the one interval [inf \(S- 8\),sup \(S+2L+8]\), inclusive. Its number of integer rows is \(\le\)\(d+2L\)+18\textless{}\(17q\). Outside these slabs there is no conflicting served occupant in Z with the halos just specified. If S is empty the latter slab is empty. The truncation at the upper end of \([0,T)\) in this argument is interior for the dates used, through the service sample \(s_c\).

A conflicting nominal nonproper-dead lower tile is nominal live by lower classification. Its nonzero nominal issue requires a computing predecessor with a genuinely live served donor or incumbent: the transition on three dead arguments returns D and a surviving native continuation or birth inherits the pertinent clock and address residues. Thus a competing nominal blocker has a served predecessor confined to one additional space/time tile. For a shifted cell that meets the pad of a queried canonical child, its near edge can differ by at most one additional child unit for rounding and less than one child unit for the pad \(64A/P\). The three predecessor candidates and these shifts fit the spatial/timing halos above. At a remaining birth date require also one previous clear canonical row and its lower stencils, already paid by the four time units. Every canonical lower input at the previous row was either proper-dead or live undisturbed, so lower proper-successor clause makes each of its successors proper, and all nine are incoming-clear at the new row. Lower directed advance therefore supplies each genuine prescribed birth. The same clearance rules out an opposite \emph{served} state of wrong age in the recruitment vacancy test.

\paragraph{A permanent ordinary startup anchor}

In the last \(20q\) donor positions omit the facing endpoint, the first site if needed to avoid a missing neighbor immediately before this interval, the common input exception of length \(2q\) and K of length d, and one neighbor on both sides of each deleted piece to ensure matching triples initially. Within the \(20q\)-window these omissions leave at most three intervals of total length \(\ge\)\(18q\)\(-\)\(d- 10\). One has length at least \((18q- d- 10)/3>q+3\). Fix \(q+3\) consecutive sites in it. Their initial inputs are proper by lower triples; their genuine survivor fields outside K have the preceding positive ballot's fresh \(\mathsf{Temp}\) and clock/address reset. Every other clean continuation keeps the prescribed result by lower proper-successor and service. The run endures for all Grow: no true Grow withdrawal acts on an ordinary address, \(\mathsf{Temp}\) cannot be set on an incumbent, and in every Purge its rightmost site has q predecessors for all propagation ticks while its other sites have a strong right bond. This argument is independent of gaps at either end of the run or clearance at the start of Grow. It is valid for a right donor with the same rightmost-site proof. Its facing site has distance \(<20q\) to the true facing endpoint.

After the initial clear rows use a complete clean Heal starting at its \(\mathsf{Temp}\) reset. Before the first such Heal, a different true matching source cannot have reached this short rim from a remote home: during at most \(6d_0\) fine ticks its total displacement is at most \(6d_0<P- 40q\), and a genuine incumbent in the target came only from its initial proper-dead source, cleared by the first complete Purge. If S meets this first startup, discard its \(\le\)\(17q\) slab and restart; the new full Heal still finishes by \(6d_0\) since two partial/full cycles plus both short slabs use \(\le\)\(4d_0+23q\)+4\textless{}\(6d_0\). A raw-only same-phase source other than the donor dies at its first complete Purge as above (genuine source cores are a distance \(P- 20q\) away). A forged copy of the donor's own label, if joined to its authentic backing, has the correct structural age/address for this geometric calculation; a fresh reset clears \(\mathsf{Temp}\) and \(\mathsf{Cn}\) before growth. A surviving different label in the facing rim could not have q consecutive genuine predecessors there and is removed by the initialized pass. Each of the B bilateral attempts now preserves the connected prefix from the anchor and either finds the next slot already occupied with the right label or recruits it by the nine-site implication. A newly attached site cannot be purged from this prefix: when its growth is rightward the right tip has q strong predecessors, and when leftward the new left tip has a strong right neighbor and the prefix's right edge still has q predecessors. During attempts no temporary retreat occurs. After \(<20q\) such steps the true endpoint is present. Subsequent temporary retreat can remove the exterior temporary extension, but its first removal from the true endpoint would test either an inward positive bond or the expressly exempt outward face. The same anchor procedure on a genuine opposite donor gives its facing rim.
\end{proof}

\begin{lemma}[Directed advance through one interruption]\label{lem:advance}
In a protected rank-j stencil of nine proper incoming-clear inputs, under the lower-rank clauses and one common exception K\(\times\)S of both spans at most d, let the central target \([0,P)\) be proper dead and an immediate neighbor a willing donor (the left one if both are willing; a right donor is treated by reflection). Then more than \(P/2\) positions match at \(s_c\). The pulse may occur at any date, including startup, first contact, the last withdrawal and normalization. The old-site debit is \(D=d+4(q+20)<5q\), with at most six complete Grow cycles charged when free progress must be repaid. Two B-long provisional probes which were never committed progress are excluded from that debit.
\end{lemma}
\begin{proof}
Let \(q=w^{2}\), \(B=128q\), \(v=w^{4}\), \(P=5\)\(\cdot\)\(2^w\), \(d=4w+8\), \(d_0=(2w^{3}+4)(2q+2)\), \(w\ge 10^9\). A round means one Main cycle of \(d_0\) fine updates; one Grow cycle has \(2v\) attack rounds followed by v withdrawal rounds. There are \(N=floor(7P/(5v))\) Grow cycles. On the nine proper incoming-clear inputs assume the startup and directed-clearance statements of Lemma~\ref{lem:startup}. Away from the common raw rectangle use preceding-rank proper-successor, service, and nominal-triple. The no-competing-phase clearance and lower nine-site advance give every eligible central birth after the designated clear rows, with the stencil conventions of Lemma~\ref{lem:transition}. These are induction hypotheses, not a hypothesis that the injured parent work later computes.



\paragraph{Reference geometry and clean progress}

Use the permanent anchor of Lemma~\ref{lem:startup} in the willing home \(L=- 1\). In each potentially opposing genuine right source choose the reflected ordinary anchor avoiding K. The source +1, if genuine, bars +2 and all more remote sources from crossing this anchor. An incumbent never changes its extended owner during Grow; to cross the fixed nonempty anchor a same-phase lineage must at some date be one of its sites, which already has the anchor owner's address. The same argument at the willing anchor bars a rear left source. If +1 is a proper-dead residue its first full Purge removes its at-most-\(q/2\) span. Up to \(P+10v\) into the rightward progress coordinate, the range bound permits only +1 or +2 in addition to L: a +3 source has smallest possible position \(3P- 1.1P=1.9P>P+10v\). Thus after startup the only genuine labels on the path are L and at most one R from \{1,2\}. Extend the reference on either side to its persistently supported facing rim when there is an opposing source. The nine-home clearance strip contains these working intervals. In the case R does not reach this region, only L has to be followed. A wrong label due to the subsequent pulse is cleared separately below.

A clean source interval from its anchor is continuous up to its first free front or a meeting. At a free face the B provisional attempts grow at most B sites; their B retreats delete the full unjoined outward extension unless it reaches its own true exempt endpoint. Only a genuine true Grow attack gives the required unit of progress into an empty target beyond the willing true endpoint. It creates an unmarked incumbent. On a withdrawal the permanent L face loses at most one; its remaining long run is backed. Between these instructions a clean unmarked proper strong run of more than q members is not pruned: only its rightmost site can have a right break and it has q strong predecessors in the pure information pass. Every clean small component with strong breaks and length \(\le\)q is erased at the Purge immediately before the elementary instruction. Thus retained detached strong pieces at such a cut have \(\ge\)\(q+1\) sites.

For different L,R labels any possible meeting in the target and up to the stated overshoot has a positive grade. At an adjacent L|R junction the extended integer drops by either \(2P- 1\) or \(3P- 1\). Both extended values are in \([- 1.1P,2.1P)\), so the left value is \(>.9P- 2>P/2\) and the right value \(<.1P+2<P/2\). Their residues are consecutive, with equal clock. Each side has its own q-strong backing. A last-slot tie chooses one donor; the filled position has an exact bond on its donor side and on its other side either another exact bond or precisely this positive junction. During an unbroken clean join the marks cannot withdraw from positive faces, neither true Grow outward face sees zero at the seam, and each separate label withstands Purge. In particular a joined prefix cannot break in its interior. A right-hand terminal point of R at \(x<P- 1\) has negative extended address relative to either possible right home. Its homeward right face can recruit at true Grow into \(x+1\) (still outside that home), and cannot be removed by true Grow in that direction. At \(x=P- 1\) or slightly farther, a negatively extended or ordinary R end does not withdraw on its right. This argument works even for a backed R piece temporarily cut away from its distant home.

Follow the longest continuous prefix from the left anchor after each completed round. Disregard unjoined provisional tips. Below \(P- 1\) a free attack extends this prefix by at least one. This follows from the preceding paragraph when its exposed label is R; when it is L its proposed address is outside its home and below its upper clip. If reaching an occupied position, incorporate the whole continuously occupied ordered piece until the next free face. Once an R-ended prefix reaches \(P- 1\) the target is safe; before a pulse this is an actual connection to the genuine opposing rim. If the prefix remains L, continue up to \(S*=P+10v\), inside the legal clip at \(1.1P\) \((10v+2<.1P)\). After first attaining \(S*\) the next withdrawal can lose at most v, and the next \(2v\) attacks restore this loss prior to the next withdrawal. A contact farther out instead pins the attained prefix. The target then has at least \(9v\) free-tip reserve or a positive join. Until one of these alternatives, a clean complete Grow cycle gives net progress \(\ge\)v; the minimum over intermediate committed fronts of that cycle is the initial front, since attack precedes withdrawal. Startup fits in the first Grow cycle. If uninterrupted, by \(\lceil 1.05P/v \rceil \) full cycles the target has the established alternatives: \((\lceil 1.05P/v \rceil - 1)v>P+10v+2\) for this w. By \(\lceil 1.1P/v \rceil \) there has also been a fresh \(\mathsf{Temp}\) reset far past first contact within the target. First contact beyond an attained overshoot has no bearing inside the target.

\paragraph{What is lost near a pulse}

For a pulse before completion of the initial recovery use the startup case below. Otherwise take as reference the latest clean \(\mathsf{Temp}\)-reset row preceding the pulse, at most \(d_0\) steps before it. The designated instruction sets every incumbent mark to zero and the following fine instruction begins a fresh Purge. At this reference the committed L and possible R intervals run from their ordinary anchors to their free tips, or meet with a positive seam. A first temporary contact made since this reference can have up to B additional marked sites per side. It costs zero previously certified sites to discard both of those tails. If the junction existed at the reference it is already unmarked; use the established-junction case.

After the \(\le\)\(17q\) interruption slab, allow one whole clean Purge and one fresh complete clean bilateral Heal. If it can finish within Grow, at most four true Grow instructions occur between the reference and the completion of its repairs: distance to pulse \(\le\)\(d_0\), its slab \(\le\)\(17q\), a complete clean Purge costs \(\le\)2(\(2q+2\)), waiting for a new round costs \(\le\)\(d_0\), and that round costs \(\le\)\(d_0\). Their sum \(<4d_0\). If it cannot finish before normalization, that normalization is at most \(17q+2(2q+2)+2d_0\) beyond the pulse. Continue to the first complete clean ordinary Heal. This adds at most two ordinary Shrink inspections before the end of the recovery, since \(17q+2(2q+2)+3d_0\)\(<qd_0\). Altogether there are fewer than ten true Grow/active-Shrink inspections in either case; allow ten exposed layers \emph{on each face}.

Consider any one contiguous reference stretch of a single label disjoint from K, initially unmarked with all its sites proper. In the normalization case restrict attention to stretches lying in one future ordinary home, so that its interior does not acquire a normalized weak wrap. If it has length \(\ge\)\(q+21\), excise the outer ten sites on each side. The remainder has at least \(q+1\) sites and endures to the clearance and recovery. Indeed, at each true Grow withdrawal or Shrink lose at most one from each exposed face, at all other instructions a reference incumbent is unmarked, and between such losses the rightmost surviving point of any run of length \(\ge\)\(q+1\) has q strong predecessors in that Purge's information pass. During killing all sites left of it have a strong right neighbor. At normalization the incoming Purge initializes \(\mathsf{Cn}\) before its first propagation, and inside this home the consecutive integer relation of the stretch persists. The preceding-rank proper-successor clause transfers every live native hold outside K, including at the two ends. No ``one layer per fine tick'' loss is needed for a genuinely proper unmarked incumbent. Initial anchors were chosen through actual earlier triples and are handled by the permanent ordinary-run proof independently of this trimming.

Restrict the reference interval of interest either to the region between the two anchors at an established join, to a free committed arm ending at its old front, or, for finishing across normalization, to the previously completely occupied target \([0,P)\). There are at most two reference labels and one excluded K. Each label minus one interval K has at most two pieces; hence at most four. For a short piece of length \(<q+21\) debit its whole length, at most \(q+20\). For a long piece debit its twenty excised margins. The total reference-occupied positions outside the resulting guard runs are at most \(D=d+4(q+20)<5q<B- 10\). This counts missing old positions across all affected pieces even when the defect lies next to an old weak join or a new true wrap. It excludes a genuine old vacancy between originally unjoined arms. The bound for a single arm is no larger. In the two-anchor case the invariant anchors themselves may be added to the guard union; both were chosen outside K. At no time can a short detached end of the former permanent progress cost more than the same D, since any extra reference stretch longer than \(q+20\) would contain a retained guard.

After the slab and before measuring fresh progress take the first whole initialized clean Purge. At most 12 hardware elementary instructions can elapse from the first pulse row up to the end of this pass, since \(17q+2(2q+2)+4<12(2q+2)\). A completely new matching raw label has its descendants within K enlarged by twelve per side until that pass, a span \(\le\)\(d+24\)\textless{}q, unless their strong component contains genuine older sites. Different labels cannot provide strong backing. A surviving component on that Purge must have started with \(>q\) sites and, in clean rows where all lower slots are proper, all such sites survive. Thus a surviving exact-label component intersects the complement of the enlarged column and has a witness of its label descended from before the event. The genuine labels in our strip are L and R in that order. This order is inherited by different surviving components: choose their pre-event genuine witnesses. Paths following chosen actual parents cannot cross. Equality at a vertex assigns the same label, and a simultaneous swap of two adjacent occupied positions would require both positions to have been vacant on that old row. For witnesses outside the twelve-site enlargement, their ancestry never intersects a raw assignment. Consequently an R run strictly to the left of a distinct L run cannot both survive this pass. A genuine rear source cannot send a bypassing witness through the permanent facing anchor. Additional honest steps preserve the order by the same argument. The sole lower raw event cannot supply another independent label afterward.

\paragraph{Recovering a free prefix, an old bridge, or an ordinary target}

In a fresh complete Grow Heal, after its \(\mathsf{Temp}\) reset and the Purge before its first attempt, surviving unobstructed label components of length \(\ge\)\(q+1\) are backed. They stay occupied throughout the B attempts. A new same-label site has an exact bond to its component. Its right tip, if exposed, has q strong predecessors, and every other site has a strong right continuation. When two components meet, the joined link is exact for equal labels and has the preceding positive grade for L|R; the long separate backings persist. Therefore no old backed component shrinks during attempts. Any gap between consecutive surviving guards in a previously continuous reference region has length at most D. If both sides have the same label all intermediate addresses of that label are legal because the legal addresses form an interval. If sides are L and R, the physical intervals of legal addresses overlap: their origins differ by \(2P\) or \(3P\) and the legal length is \(3.2P\) up to rounding, with overlap \(\ge\)\(.2P- 2\). If both inward bids were clipped there would be a positive gap between those legal intervals; hence at least one advancing face is legal whenever the gap is nonempty. With foreign phases absent the age test passes, and lower advance supplies a birth. Thus vacancy count falls by at least one per attempt until the guards join, within \(D<B\) attempts. Extra correctly ordered backed pieces in the gap only shorten this count. In a reached join every recruited site has positive bonds on both faces and its label has \(\ge\)\(q+1\) contiguous backing. It survives the B temporary removals. This argument handles several subgaps in parallel (each has length bounded by the \emph{total} D); it does not assume both sides reclaim their original labels.

For the old free L arm this restores its prefix from the fixed facing anchor to within D of its pre-event tip: apply the preceding gap argument to consecutive surviving guards within its reference interval. At most D positions of the reference interval lie to the right of the last guard in that interval. If its whole new exterior extension was temporary only those old D positions are charged, and the originally occupied ordinary anchor still starts the prefix. For a formerly joined arm apply the argument between the two permanent facing anchors. This restores the entire target. Once the join is stable, on the R-labelled right side the homeward face does not withdraw, and the connected target is protected. For a lone prior overshoot the restored prefix reaches at least \(P+9v- D>P+8v\); even a whole first remaining withdrawal cannot expose the target. The clean prefix argument then applies to later cycles.

For a pulse meeting the final cycle too late for a whole Grow Heal the prior late geometry already covered the target. Compute its guards \emph{restricted to H} with the bound D above. These survive including normalization; each is in-H, unmarked before any true ordinary loss, of length \(\ge\)\(q+1\) after trimming, and their union contains \(\ge\)P\(-\)D sites. Thus there is at least one in-home guard. All positions in H to the left of its leftmost guard or right of its rightmost guard have combined length \(\le\)D. After a full clean Purge any additional matching fragments are backed, and every occupied position in H then has the unique ordinary local address. During ordinary attempts join adjacent guards by the same argument. The leftmost guard also reaches address zero, and the rightmost guard reaches \(P- 1\), each within \(D<B\) attempts, by one-sided legal in-home bids. Both outward faces are exempt at their true endpoint. This completes the ordinary in-home repair, permitting several local subgaps with total deficit D. The argument never demands an intact target exactly on the normalization row. If the pulse first occurs in later ordinary work apply this same guard count in H to the previously full ordinary target at the last \(\mathsf{Temp}\) reset. If the service sample arrives before full healing, the old guards alone supply \(\ge\)P\(-\)D\textgreater{}P/2. With the event spent, after normalization with its fresh \(\mathsf{Cn}\) reset or after complete repair, an ordinary full home stays full and proper to \(s_c\) by the native full-home argument of Lemma~\ref{lem:endpoint}, using the same mid-period start as in Lemma~\ref{lem:two-period}. When the sample arrives before the pulse, the same argument applies to the earlier full ordinary home. No subsequent source advance is required.

\paragraph{Exhaustive timing and progress cases}

The cases below refer to the first possibly contaminated fine update. A pulse spanning a boundary uses the case covering its latest required cleanup. In every case in which the target is covered at the end of Grow and the pulse, if any, is spent within Grow (no pulse, or a completed Grow repair), it is full at normalization and the full-home argument just cited carries it, full and proper, to \(s_c\).

\paragraph{Interruption cases.} Appendix~\ref{sec:case-tables} tabulates the associated time and vacancy bounds.
\medskip\noindent\textbf{No pulse, or pulse whose initial/cleanup interval meets startup before the first secured reset.} Use the K-avoiding permanent ordinary anchor(s). Restart the fully clear startup Heal; even a pulse in the last initial recovery round costs \(\le\)\(12d_0\) from the cut, below one \(3v\)-round cycle. A remote rival cannot arrive in \(12d_0<P- 40q\). Drop this first whole cycle. The remaining \(N-1\) clean cycles suffice to attain contact or \(P+10v\).

\medskip\noindent\textbf{Free arms; damage during a \(2v\) attack, including an unreset first \(L/R\) probe join.} Reference taken before the event; debit no pre-reference sites for the two provisional tails, and \(\le\)D for its old L tip. All internal old gaps bridge. Exclude the cycle of the reference and any cycle intersecting the repair. At most two Grow cycles touched; on resuming mid-cycle allow an additional v of retreat before the next whole cycle.

\medskip\noindent\textbf{Free arms; damage in a v withdrawal, including a provisional join first made there.} Same reference and D. At most four true Grow steps before a clean repair; the partially withdrawn reference already counts its old state. The same two-cycle discard and one residual v bound cover an interruption that crosses the attack/withdrawal or whole-cycle boundary.

\medskip\noindent\textbf{Previously committed \(L/R\) contact, in an attack or withdrawal.} The adjacent anchors are restored by one full Grow Heal using \(\le\)D missing old sites; the established weak seam's two possible short cut-off fragments were both paid in D. If enough Grow remains for that Heal there is no distance repayment. Otherwise use the final-cycle row below.

\medskip\noindent\textbf{Lone overshoot already attained; pulse at the tip, the clip, a remote arm, or the true donor border.} Old prefix restored to \(P+9v- D\) or farther, or its genuine joined case applies. The legal clip is \(>S*\). One whole v retreat still leaves target covered. Subsequent complete attacks make up clean withdrawals.

\medskip\noindent\textbf{Pulse in the final withdrawal or straddling normalization, without room for a complete clean Grow Heal.} The undamaged interval before it contains more than \(\lceil 1.1P/v \rceil \) full cycles. The target was covered by committed geometry. Restrict old guards to H, permit D missing even on the normalized cut itself, complete the ordinary in-home Heal. At most two early Shrink inspections in the specified recovery window.

\medskip\noindent\textbf{Pulse starts after Grow.} Old full ordinary H, its D-debit guards or the already observed service sample decide the count. With the event spent, a restored ordinary full home stays full to \(s_c\) by Lemma~\ref{lem:endpoint} (with the mid-period reset used in Lemma~\ref{lem:two-period}); before healing the guards give \(\ge P-D\), and before the pulse the old full home gives P.


For the free-arm rows let c be the complete pre-reference cycles after startup. If the old tip is still unsatisfied its attained displacement is at least cv at the reference, since completed clean cycles give \(\ge\)v and progress never dips below its starting frontier during the following partial clean cycle. After recovery the displacement is at least cv\(-\)D. The reference-to-completion interval is \(<4d_0<3v\) \(d_0\), meeting at most two whole cycles even if it crosses their boundary. Until the next cycle boundary one additional withdrawal costs at most v. The initial discarded cycle, the two interrupted cycles and this one additional v use four cycle-equivalents. Allow two further complete cycle-equivalents for the two integer endpoints in counting disjoint pre- and post-reference complete cycles. Thus the eventual free displacement lower bound, until permanent contact or overshoot, is \((N- 6)v- D\). It exceeds \(P+10v+2\): \(N\ge 1.4P/v- 1\), \(D<5q<v\), and \(.4P>18v+2\). The late threshold used by the final row has spare room since \(1.1P/v+3<N\) for this w. The omitted and the still-open old macroscopic gap never enter the D count; progress into such a gap is paid only from the available honest attack rounds.

Finally if only the right donor is willing, read all prefix statements from the right. At the encounter the label order in physical left-to-right coordinates is still the smaller-origin label followed by the larger, giving the same positive junction and left-priority resolution. A reflected free right worker has legal decreasing addresses beyond its own left edge until it reaches the reflected overshoot. A homeward end of a left source in the target has positive extended address and true Grow cannot kill it on its left. Each Purge proof still uses the physical rightmost site's q physical left predecessors. The clipping bound is again \(\ge\)\(1.1P- 2\). These are the only directional tests used by the argument. Combining the service count with the already derived strong issue and lower proper successor gives the requested central nominal birth. It is served under the ordinary proper-live service clause, once the positive successor has been made proper by the end-repair lemma.

This discharges the induction step using only preceding-rank protected holds, cuts, service and mandatory recruitment. For rank one the child history is the actual radius-one data rule off its flags, with physical service equal to identity. There are no distinct geographic rank-zero grids. Thus each finite rank is treated from its base, with the zero-top-exception variant whenever the entire window is already sparse one level lower.
\end{proof}

\begin{lemma}[Two periods of served children]\label{lem:two-period}
In a fully protected rank-j \(W_0\) test that is \((j- 1)\)-sparse, use lower coordinates and translate the first work interval to \([0,T)\). Number the six consecutive served positive inputs at its bottom 0,...,5. They belong to the specified common tribe. Suppose the two prescribed transitions for owners 2 and 3, using these inputs and the fixed environment, are live with the stipulated next tribe residues. Then every lower site in those two homes is served with the correct lower structural residues on \([0,2T)\), hence on the requested \([T/2,2T)\); their two upper successors at T also serve the prescribed values. The endpoint \(2T\) is excluded. If T is odd the left endpoint of the requested integer interval is its ceiling.
\end{lemma}
\begin{proof}


\textbf{Earlier premises.} The rank-j precommit dichotomy and simultaneous full ordinary seed for up to six computing predecessor homes, with J empty, are supplied by Lemma~\ref{lem:seed}. A noncomputing predecessor returns nominal D. These are rank-j past-epoch conclusions; neither rank-j proper-live service nor rank-j directed advance is assumed. At rank \(j- 1\) use its protected exact-computation, nominal-triple, proper-successor and service clauses, including Lemma~\ref{lem:endpoint}, with the conventions of Lemma~\ref{lem:transition}. At \(j=1\) the exact unflagged native update replaces the lower clauses. Empty J here follows from the stronger \((j- 1)\)-sparsity of the entire test; lower-rank sparse defects are allowed. Without the earlier seed lemma this result is a conditional induction implication.

\begin{enumerate}
\item \textbf{The four full predecessor homes.} The six served positive bottom values imply six positive preceding nominal issues. A noncomputing predecessor issues D, so all six preceding homes computed. Select their simultaneous length-\(5d_0\) full stage-0 interval before the first read. Each handle's Purge reset is a legal starting time for the full-home argument; the distance to the next such time is at most \(2q+2<5d_0\). At that time the four inner homes 1,...,4 have all P occupants, and every one of these \(4P\) sites has two matching neighbors in the six-home stretch, also across the aligned wraps. The lower nominal-triple rule makes each proper; a served matching triple is nominal matching since a served live value equals its nominal value. The next reset gives \(\mathsf{Cn}\)=1 simultaneously. During the reset all four homes still have positive neighbors. Now apply the full-home specialization of Lemma~\ref{lem:endpoint} up to the first ballot at time zero. Its proof requires only ordinary addresses, common genuine clock and a \(\mathsf{Cn}\) reset for a mid-period start: the \(\mathsf{Temp}\)-removal clause exempts the outward true face even for an initially marked endpoint and every other face is positively bonded. All \(4P\) sites stay proper and served.
\item \textbf{All four ballots at zero retain their endpoints.} In each of these four homes the three clean reset evaluations give, separately for each repetition, one packet string on the central output fifth: a serialized live result with every letter live, or all message-empty. This is the total evaluator's property for arbitrary received input including a malformed supplier. Fixed shifts to the owner's five fifths are wholly inside that full home; induction on the shift length transports each intended letter, since every crossed interior bond is strong, every required overwrite occurs and no site is born or dies. Thus every site of that home has the same triple of bank liveness bits (payloads may differ by offset). If at least two bits are live, every site starts verification \(\mathsf{Active}\), stays \(\mathsf{Active}\) and votes live. Otherwise all start inactive with cleared scratch. During verification strict in-owner edges forbid entry from outside, no births occur in a full home, and each epoch starting with no \(\mathsf{Active}\) bits has zero arrivals everywhere. No counter exceeds \(4d\), all sites therefore vote D. Lower exact computation at proper sites makes this nominal result D at every one of the P sites; their served states would all be D, so their parent could not have the assumed positive nominal issue. This excludes the second case individually in each of homes 1,...,4. The first case and lower proper-successor clause now give a proper served reset at every site of those homes, including the two endpoints of each. The argument never supposed that a lower triple survived the vote.
\item \textbf{The first work interval \([0,T)\).} Starting at zero, the four reset homes stay full and proper to the pre-ballot row by Lemma~\ref{lem:endpoint}. Every old carrier retains its X until the ballot. For both computing centers 2 and 3, all three contributing supplier homes are among those four. For each supplier and fixed offset, the three stored homologous coordinates are spaced \(P_0\) apart. Nominal acceptance at time zero certifies that all three carry the appropriate encoded letter off one interval of at most \(2q<P_0\) positions. Hence at least two are correct at every offset. Each inject/shift/receive between these consecutive full homes crosses only strong in-home links or the synchronized positive \(P- 1\) to 0 wrap. Induct on unit shifts from injection to the receiving full site; no shift can substitute another source. The majority gives the exact packet at every offset, all three parsers recover the respective served input, and evaluation with the frozen operands gives the prescribed live transition at the two centers in each repetition. Readout inside each full center is exact at every site by step 2, so all sites there start verification \(\mathsf{Active}\). They remain \(\mathsf{Active}\) to the ballot, return their appropriate live letters with zero age and cleared scratch, and lower proper-successor clause followed by lower proper-live service gives proper service of every endpoint across time T. The exterior neighbors at that vote may die.
\item \textbf{The second work interval \([T,2T)\).} Apply the reset-home proof to homes 2,3, independently, now without using occupancy of homes 1 and 4. On every pre-ballot structural step each proper site prescribes a live result by the instruction calculation and gets a proper served successor from the lower clauses. Their clocks and address residues are those of their own work interval by the true schedule. Thus both central spatial homes are full at every date T through \(2T- 1\). At \(T+s_c\) each has exactly P matching occupants, exceeding \(P/2\), and at T the full correct payload gives each prescribed nominal successor. They therefore serve those upper values. On the physical grid recursive lookahead obeys \(b_0=0\) and \(b_k\le s_c\) \(T^{k- 1}+b_{k- 1}\le (T- 1)T^{k- 1}+T^{k- 1}- 1=T^k- 1\), so every last lower service observation fits its own tile.
\end{enumerate}

\end{proof}

\paragraph{Assembly.} Lemmas~\ref{lem:ballot} and~\ref{lem:classification} give the issue and classification, Lemma~\ref{lem:exclusion} gives separation and the matching middle, Lemma~\ref{lem:end-ledger} gives service and proper successors, and Lemma~\ref{lem:advance} gives the central requested birth. The dead successor uses the small-group cut. The structural hold and ancestry conclusions apply wherever the child step is protected. At rank one the calculation starts with the native step. This closes the induction in Proposition~\ref{prop:organization}.
